\documentclass[11pt,a4paper]{article}

\usepackage[utf8]{inputenc}
\usepackage[T1]{fontenc}
\usepackage{lmodern}
\usepackage{microtype}
\usepackage[a4paper,margin=1in]{geometry}
\usepackage{amsmath,amssymb,amsthm,mathtools}
\usepackage{mathrsfs}
\usepackage{booktabs}
\usepackage{enumitem}
\usepackage{array}
\usepackage{xcolor}
\IfFileExists{orcidlink.sty}{\usepackage{orcidlink}}{}
\usepackage{hyperref}
\usepackage[nameinlink,noabbrev]{cleveref}

\providecommand{\orcidlink}[1]{%
  \href{https://orcid.org/#1}{\textsuperscript{\scriptsize ORCID}}%
}
\setlength{\emergencystretch}{2em}

\hypersetup{
  colorlinks=true,
  linkcolor=blue!55!black,
  citecolor=blue!55!black,
  urlcolor=blue!55!black,
  pdftitle={Finite Dilation Operator Algebra of Homothetic B4/S3 Valuations and Homogeneous Closure Tests for Valuation Dark Energy},
  pdfauthor={Andrei Preece and Boris Batenin},
  pdfsubject={Closed cosmology; quasilocal filling; effective residual sector},
  pdfkeywords={closed cosmology; S3; B4 filling; quasilocal energy; residual dark sector; valuation spectrum; Noether identities; S4 compensation}
}

% --- Fix for hyperref warnings: safely expand math macros in PDF strings ---
\pdfstringdefDisableCommands{%
  \def\Bfour{B4}%
  \def\Sthree{S3}%
  \def\Sfour{S4}%
  \def\dd{d}%
  \def\Dil{D}%
  \def\QU{Q_U}%
  \def\QG{Q_G}%
}

\setlist[itemize]{leftmargin=*,itemsep=0.2em,topsep=0.4em}
\setlist[enumerate]{leftmargin=*,itemsep=0.25em,topsep=0.4em}
\renewcommand{\arraystretch}{1.15}
\numberwithin{equation}{section}

\theoremstyle{definition}
\newtheorem{criterion}{Criterion}
\newtheorem{definition}[criterion]{Definition}
\newtheorem{hypothesis}[criterion]{Hypothesis}
\newtheorem{closurecondition}[criterion]{Closure condition}
\newtheorem{nogotheorem}[criterion]{No-go theorem}
\theoremstyle{plain}
\newtheorem{lemma}[criterion]{Lemma}
\newtheorem{theorem}[criterion]{Theorem}
\newtheorem{proposition}[criterion]{Proposition}
\newtheorem{corollary}[criterion]{Corollary}
\theoremstyle{remark}
\newtheorem{remark}[criterion]{Remark}

\newcommand{\Sthree}{S^3}
\newcommand{\Sfour}{S^4}
\newcommand{\Bfour}{B^4}
\newcommand{\M}{\mathcal{M}}
\newcommand{\Dom}{\mathcal{U}}
\newcommand{\Dil}{\mathscr{D}}
\newcommand{\B}{\mathcal{B}}
\newcommand{\Efour}{E_4}
\newcommand{\rhoFour}{\rho_4}
\newcommand{\epsD}{\varepsilon_D}
\newcommand{\pD}{p_D}
\newcommand{\QD}{Q_D}
\newcommand{\PiA}{\Pi_A}
\newcommand{\vol}{\operatorname{Vol}}
\newcommand{\diag}{\operatorname{diag}}
\newcommand{\dd}{\mathrm{d}}
\newcommand{\QG}{Q_{\rm G}}
\newcommand{\QU}{Q_U}

\title{Finite Dilation Operator Algebra of Homothetic \(\Bfour/\Sthree\) Valuations\\
       and Homogeneous Closure Tests for Valuation Dark Energy}
\author{Andrei~Preece\,\orcidlink{0009-0004-5978-795X}\\
        \small \href{mailto:andre.preece@web.de}{andre.preece@web.de}\\
        \small Independent Researcher, Berlin, Germany
        \and
        Boris~Batenin\,\orcidlink{0009-0000-7027-8583}\\
        \small \href{mailto:batenin@web.de}{batenin@web.de}\\
        \small Independent Researcher, Berlin, Germany}
\date{2026}

\begin{document}
\maketitle



\begin{abstract}
We develop a homogeneous residual-profile framework for closed FLRW
cosmology using a round Euclidean \(\Bfour_A\) as an auxiliary quasilocal
filling of \(\Sthree_A\).  The universal work map
\[
\Efour(A)\mapsto
\left(\epsD=\Efour/V_3,\quad
\pD=-\dd\Efour/\dd V_3\right)
\]
is algebraically invertible at homogeneous level and is therefore a
representation theorem rather than a predictive closure.  Predictivity begins
after imposing the Euclidean convex valuation hypothesis on the auxiliary
round filling.  The induced density profile has the finite dilation spectrum
\[
S_{\rm val}=\{1,0,-1,-2,-3\},
\qquad
P_{\rm val}(z)=(z-1)z(z+1)(z+2)(z+3),
\]
and the polynomial action of \(\Dil=A\dd/\dd A\) factors through the
semisimple operator algebra
\[
\mathfrak A_{\Dil}\simeq \mathbb R[z]/(P_{\rm val})\simeq \mathbb R^5.
\]
Its primitive idempotents give canonical spectral projectors, branch
coefficients, dimensionless fractions, minimal annihilators, and
covariance-aware finite-dimensional null tests.  For nonnegative branch
energies,
\[
\Dil^2\ln(\epsD/\varepsilon_*)=\operatorname{Var}_\pi(s)\ge0,
\qquad
\Dil w_D=-\operatorname{Var}_\pi(s)/3\le0,
\]
with sharper diameter and Bhatia--Davis bounds.

A filling-relative local curvature-measure boundary filter removes the bulk
\(s=1\) branch.  At the level of \(\mathcal H^2(a)\) the remaining nonstandard quotient
direction is \(a^{-1}\), with \(w=-2/3\), but a physical \(a^{-2}\) stress must
be separated from geometric curvature by transverse-distance data.  Thus the
general boundary-profile extension is two-parameter,
\((\Omega_X,\Omega_{2,\rm val})\), while the one-parameter VDE model is the
nested subclass \(\Omega_{2,\rm val}=0\).
Positive \(s=-1\) and \(s=-2\) branches admit independent canonical-scalar
realizations, and the \(s=-3\) branch admits a conserved pressureless-current
realization.  For the nested canonical VDE branch the background trajectory
fixes the potential, rest-frame sound speed, anisotropic stress, and scalar
mass-flow invariant.  These source actions are existence realizations of
selected homogeneous branches; they are not derived by covariant variation of
the auxiliary valuation functional.  Topology alone does not determine
dimensional amplitudes, and a full compact-\(\Sthree\) Boltzmann likelihood
implementation remains open.
\end{abstract}

\noindent\textbf{Keywords:} closed cosmology; \(\Sthree\); \(\Bfour\) filling; quasilocal energy; effective residual dark sector; Noether identities; \(\Sfour\) compensation.

\section*{Summary of the Result}
\begin{center}
\begin{tabular}{@{}p{0.28\linewidth}p{0.65\linewidth}@{}}
\toprule
Layer & Result and status \\
\midrule
Homogeneous representation & \(\Efour\leftrightarrow(\epsD,\pD,\QD)\): exact reconstruction at background level, not source microphysics. \\
Euclidean valuation profile & Five dilation eigenvalues \(S_{\rm val}=\{1,0,-1,-2,-3\}\): proved for the round homothetic image under the Euclidean convex valuation hypothesis. \\
Dilation operator algebra & \(\mathfrak A_{\Dil}\simeq\mathbb R^5\), primitive projectors, minimal annihilators, and spectral lattice: proved; \(\mathcal V_{\rm val}\) is a faithful module, not a pointwise algebra. \\
Positive mixtures & Variance, Bhatia--Davis, endpoint/\(\Lambda+X\) flows, Hankel positivity/rank: proved for nonnegative active branch energies; the bulk \(s=1\) slot is algebraic and not a stable source sector. \\
Boundary filter & Local curvature-measure densities relative to the declared Euclidean filling remove \(s=1\); the filter is conditional and filling-relative. \\
Local flat continuation & Fixed valuation coefficients make the nonconstant boundary branches decouple as \(A_0\to\infty\); finite local amplitudes require curvature-radius-dependent matching coefficients. \\
Expansion quotient and curvature split & \([a^{-1}]\) is the expansion-rate quotient; \(a^{-2}\) stress is separated from geometry by \(\Omega_{2,\rm val}=\mathcal C_{-2}-\Omega_K^{\rm geom}\). \\
Source realizations & Positive \(s=-1,-2\) branches have independent canonical-scalar realizations; the nested \(s=-1\) VDE subclass has a fixed one-field potential and compact scalar source kernel; \(s=-3\) dust has a pressureless-current realization. These are existence lifts inside declared source classes and are not derived from \(\Efour^{\rm val}\). \\
Endpoint classification & Inside the restricted oriented endpoint class, \(E_{\rm end}=\sigma V_3+\sum_A\mu_AQ_A\); charge-current equality is a matching condition, conserved once imposed. \\
Amplitudes & Pure topology cannot determine nonzero dimensional amplitudes; coefficients are dimensional matching/coupling data, Wilson coefficients only after an explicit EFT embedding. \\
Observational status & One-parameter \(\Lambda+X\) is a small nested null-test branch; \(\Omega_{2,\rm val}\) is mainly an auxiliary curvature-splitting parameter; compact \(S^3\) origin is conditional and not empirically established. \\
Analytic reproducibility & The accompanying analytic package supplies TeX/PDF, build metadata, checksums, and machine-readable scope ledgers; no Boltzmann code, likelihood chains, or observational data products are claimed. \\
Transfer level & The canonical scalar source equation and compact scalar kernel are supplied; the complete closed-SVT Einstein--Boltzmann hierarchy, covariance prediction, lensing map, and exact-model likelihood remain not supplied. \\
\bottomrule
\end{tabular}
\end{center}

\tableofcontents

\section{Position of the Work and Scope of the Result}
This paper is a local mathematical node of a broader research corpus. The wider ontological program motivates the construction but is not a theorem-level output of the present work.

The proved core of the present work is the chain
\begin{equation}
\Sthree_A\longrightarrow \Bfour_A\ \text{filling}\longrightarrow \Efour(A,\xi)\longrightarrow (\epsD,\pD,\QD).
\end{equation}
The \(\Sfour\)-compensation construction is retained only as a downstream extension and moved to an appendix. Only orientational cancellation under compatible gluing is proved here. The scale-covariance no-go theorem below prevents pure topology from selecting nonzero dimensional amplitudes without an additional matching scale.

This chain is not presented as an autonomous brane model, a standard extra-dimensional cosmology, or a physical embedding of the observed space into a five-dimensional bulk. Similarity between an isolated formula and a formula in another model does not establish structural equivalence. The adjective quasilocal is used here only in the sense that a bounded spatial domain is assigned an auxiliary filling energy. It does not imply the Brown-York boundary stress tensor, the Hamilton-Jacobi gravitational energy, or a physical bulk embedding. In particular, \(\Efour\) is not identified with the Brown-York gravitational quasilocal energy \cite{BrownYork1993}, whose variational origin and boundary stress tensor are structurally different. No Hamiltonian boundary stress tensor is constructed from \(\Efour\); the term quasilocal refers only to the assignment of an auxiliary filling energy to a spatial domain.

The precise result has two directions. The direct construction produces an effective background fluid from a filling energy. The inverse representation theorem shows that every differentiable homogeneous background density has a unique representing \(\Bfour\) density within the adopted homothetic class; a prescribed pair \((\epsD,\pD)\) is represented by the adiabatic map exactly when it satisfies the corresponding continuity identity. This establishes both the strength of the construction and the exact limit of the universal map: geometry fixes the representation. Predictivity begins only after an invariant closure or additional dynamical principle selects a restricted functional class for \(\Efour\).

The claim hierarchy used throughout the paper is therefore
\begin{equation}
\begin{aligned}
\text{valuation restriction}
&\Rightarrow
\text{homogeneous profile closure},\\
\text{homogeneous profile closure}
&\not\Rightarrow
\text{source action}
\not\Rightarrow
\text{amplitude selector}
\not\Rightarrow
\text{transfer prediction}.
\end{aligned}
\end{equation}
Equivalently, the primitive residual reconstruction, source realization, amplitude matching, transfer evolution, covariance prediction, and likelihood comparison are distinct layers. The paper proves the homogeneous valuation-profile layer and provides independent source realizations for selected branches; it does not claim a covariant derivation of those sources from the auxiliary valuation functional.

The protective separation is summarized by
\begin{equation}
\boxed{
\text{reconstruction}
\neq\text{source}
\neq\text{amplitude}
\neq\text{transfer}
\neq\text{likelihood}.
}
\end{equation}
This separation is part of the model definition, not a disclaimer added after the fact.

\paragraph{Statement of results.}
The universal \(\Bfour/\Sthree\) representation is intentionally separated from the predictive closure. The central invariant objects are the five-dimensional round valuation-profile space and the polynomial action of the dilation operator on it. The resulting five-eigenvalue semisimple operator algebra is
\begin{equation}
\mathfrak A_{\Dil}
=\mathbb R[z]/[(z-1)z(z+1)(z+2)(z+3)]
\simeq\mathbb R^5,
\end{equation}
with canonical primitive projectors. This yields exact linear closure tests, a lattice of spectral subspaces, and--under positivity--nonlinear moment and convexity constraints.

The local curvature-measure boundary restriction selects \(S_\partial=\{0,-1,-2,-3\}\). Its quotient by the standard scalings is one-dimensional only as a statement about the normalized expansion function \(\mathcal H^2(a)=H^2(a)/H_0^2\). Once transverse distances are included, the \(a^{-2}\) valuation fluid must be separated from geometric curvature. The general physical boundary-valuation background extension of curved \(\Lambda\)CDM is therefore generated by \((\Omega_X,\Omega_{2,\rm val})\), while the one-parameter VDE model is the nested subclass \(\Omega_{2,\rm val}=0\).

Separately, inside the oriented minimal endpoint class, a conditional classification theorem gives \(E_{\rm end}=\sigma V_3+\mu\QU\). A scale-covariance no-go closes the topology-only amplitude question negatively: nonzero dimensional amplitudes require external matching data. The global integer charge and the local dust current are not identified pointwise. A conserved pressureless current supplies the endpoint dust dynamics, while canonical scalars provide explicit covariant realizations of the positive \(s=-1\) and \(s=-2\) valuation branches. Observational superiority remains a likelihood-level question.

\subsection{Precedents, non-equivalence, and exact novelty}
The elementary ingredients used here have important precedents.  Hadwiger's
classification and the intrinsic-volume basis are classical results in convex
geometry \cite{Hadwiger1957,Schneider2014}.  Closely related morphometric
reductions occur in morphological thermodynamics, where motion invariance,
continuity, and additivity reduce certain shape-dependent thermodynamic
functionals to finitely many curvature measures \cite{KonigRothMecke2004}.
Inverse reconstruction of quintessence potentials is also a known programme
\cite{SahlenLiddleParkinson2005}, and boundary likelihood-ratio laws for
parameters constrained to a half-line belong to the Chernoff--Self--Liang
statistical setting \cite{Chernoff1954,SelfLiang1987}.  None of these facts is
claimed here as new.

The contribution of this paper is the particular residual-cosmology synthesis:
under the declared Euclidean convex-valuation hypothesis, the round homothetic
\(\Bfour_A/\Sthree_A\) density profile becomes a faithful module over a finite
semisimple dilation-operator algebra.  This turns the homogeneous residual
sector into a spectral calculus with primitive projectors, minimal
annihilators, a positivity cone, curvature-splitting diagnostics,
covariance-aware finite-dimensional tests, and source-realization checks.  The
claim is therefore not
\begin{equation}
\text{topology}\Rightarrow\text{dark energy amplitude},
\end{equation}
but rather
\begin{equation}
\boxed{\begin{gathered}
\text{Euclidean valuation closure}\\[1mm]
\Rightarrow\ \text{finite dilation spectrum}\\[1mm]
\Rightarrow\ \text{exact homogeneous null tests and source diagnostics}.
\end{gathered}}
\end{equation}

This distinction also separates the present construction from Brown--York
boundary stress tensors, AdS/Fefferman--Graham holographic renormalization,
Faddeev--Skyrme Hopfion models, and Wald--Zoupas charge formalisms.  Those may
be useful in other theories or optional extensions, but they are not consequences
of the auxiliary Euclidean valuation framework used here.  The result below is
accordingly a finite homogeneous invariant classification plus a falsification
programme, not a completed microscopic origin theory.


\begin{criterion}[Minimal physical completeness]
A candidate physical dark sector must provide more than a background function \(H(a)\). It must also specify an energy-momentum tensor, a conservation or exchange law, and rules for perturbations. The background map is a necessary but not sufficient level of a physical model.
\end{criterion}

\section{Conventions and Geometric Arena}
\subsection{Carrier, Metric Representative, and Reconstruction Level}
\begin{definition}[Carrier, metric representative, and reconstruction level]
Let
\begin{equation}
\Sigma_{\rm car}\simeq \Sthree
\end{equation}
denote the underlying topological carrier. For each \(A>0\),
\begin{equation}
(\Sthree,g_A),\qquad g_A=A^2\dd\Omega_3^2,
\end{equation}
is its round homothetic metric representative. For an observer \(O\) at proper time \(\tau\),
\begin{equation}
\Dom_O(\tau)\subseteq \Sthree_A(\tau)
\end{equation}
denotes the causally accessible or operationally reconstructed domain.
\end{definition}

These three levels are not identified. In particular, \(A(\tau)\) is a scale parameter of the metric representative. The present background map does not by itself interpret limiting values of \(A\) as creation, destruction, collapse, or rupture of \(\Sigma_{\rm car}\).

We use metric signature \((-+++)\). Covariant matter-action and scalar-field formulae are written in units \(c=\hbar=1\), with \(M_{\rm Pl}^{-2}=8\pi G\). Factors of \(c\) are restored explicitly in geometric distance, curvature, reconstructed-Einstein-tensor, and observational formulae. In the natural-unit matter sector, \(\mu\) is an energy-per-charge matching coefficient; interpreting it as a mass per charge is equivalent only after setting \(c=1\). The physical arena is the closed \(k=+1\) FLRW class,
\begin{equation}
\M=I\times \Sigma,
\qquad
\Sigma\simeq \Sthree,
\qquad
A(t)=a(t)R_0.
\end{equation}
We use the normalization \(a(t_0)=1\), so that
\begin{equation}
A_0=A(t_0)=R_0.
\end{equation}
The metric representative is
\begin{equation}
\dd s^2=-N(t)^2c^2\dd t^2+A(t)^2\dd\Omega_3^2.
\end{equation}
This is the standard closed FLRW convention used in relativistic cosmology \cite{HawkingEllis1973,MTW1973,Weinberg2008}.
Proper time satisfies \(\dd\tau=N\dd t\), and
\begin{equation}
H=\frac{1}{A}\frac{\dd A}{\dd\tau}=\frac{\dot A}{NA},
\end{equation}
where the dot denotes differentiation with respect to coordinate time \(t\). Derivatives with respect to proper time are written explicitly as \(\dd/\dd\tau\).

For the round three-sphere,
\begin{equation}
V_3(A)=\vol(\Sthree_A)=2\pi^2A^3,
\qquad
{}^{(3)}R=\frac{6}{A^2}.
\end{equation}
The Riemannian filling \(\Bfour_A\) has
\begin{equation}
\partial \Bfour_A=\Sthree_A,
\qquad
V_4(A)=\vol(\Bfour_A)=\frac{\pi^2}{2}A^4,
\qquad
\frac{\dd V_4}{\dd A}=V_3.
\end{equation}
Hence
\begin{equation}
\frac{V_4}{V_3}=\frac{A}{4}.
\end{equation}
The factor \(A/4\) is a metric consequence of the round homothetic filling, not a consequence of the topological relation \(\Sthree=\partial\Bfour\) alone.

\begin{remark}[Auxiliary status of the filling]
The Riemannian ball \(\Bfour_A\) is attached to the metric representative \((\Sthree,g_A)\), not to the bare topological carrier alone. It is an auxiliary filling used to define a geometric measure and a variational energy; it is neither an additional physical spatial dimension nor a physical five-dimensional spacetime bulk.
\end{remark}

The \(\Bfour\)-filling can geometrically represent any separately conserved homogeneous background component within the adopted homothetic class, but it does not determine its microscopic origin, perturbations, sound speed, anisotropic stress, particle content, or full spacetime dynamics.

\begin{criterion}[Compactness stress test]
Every residual-sector hypothesis should be testable both in the global compact class \(\Sthree_A\) and in the local flat limit \(A\to\infty\), \(\Sthree_A\to\mathbb{R}^3\) at fixed finite observational scale, understood as a local or pointed geometric limit rather than a change of global topology. A construction that requires openness must identify the physical term or boundary condition that requires it.
\end{criterion}

\section{Quasilocal Map and Geodesic Domains}
For a domain \(\Dom\subset \Sthree_A\), consider the quasilocal map
\begin{equation}
\Dom\longmapsto \B(\Dom)\subset \Bfour_A.
\end{equation}
Let \(\Dom_\chi\) be a geodesic ball of angular radius \(0\le \chi\le\pi\). Then
\begin{align}
V_3(\chi,A)&=4\pi A^3\int_0^\chi \sin^2u\,\dd u
=2\pi A^3(\chi-\sin\chi\cos\chi),\\
A_2(\chi,A)&=4\pi A^2\sin^2\chi.
\end{align}
Define the homothetic cone by transporting the same angular domain homothetically to each intermediate radius:
\begin{equation}
\B(\Dom_\chi):=\{(r,\Omega)\mid 0\le r\le A,\ \Omega\in \Dom_\chi^{\rm ang}\},
\end{equation}
where \(\Dom_\chi^{\rm ang}\subset \Sthree\) denotes the angular geodesic ball of radius \(\chi\) on the unit sphere.

\begin{lemma}[Volume of the conical filling]
For fixed \(\chi\),
\begin{equation}
V_4(\chi,A)=\frac{A}{4}V_3(\chi,A),
\qquad
\frac{\partial V_4(\chi,A)}{\partial A}=V_3(\chi,A).
\end{equation}
\end{lemma}
\begin{proof}
Write \(V_3(\chi,r)=C(\chi)r^3\), where \(C(\chi)=2\pi(\chi-\sin\chi\cos\chi)\). Then
\begin{equation}
V_4(\chi,A)=\int_0^A V_3(\chi,r)\,\dd r
=C(\chi)\int_0^A r^3\,\dd r
=\frac{C(\chi)A^4}{4}
=\frac{A}{4}V_3(\chi,A).
\end{equation}
Differentiation with respect to \(A\) gives the second identity.
\end{proof}

The full boundary of \(\B(\Dom_\chi)\) contains the outer domain \(\Dom_\chi\), the lateral conical part over \(\partial\Dom_\chi\), and their codimension-two corner. The conical-volume identity does not remove possible lateral or corner contributions in an action variation.

\section{Filling Energy and Two Distinct Responses}
Let \(\rhoFour(A)\) be the auxiliary filling density per spatial four-volume, \([\rhoFour]=EL^{-4}\). Define
\begin{equation}
\Efour(A):=\rhoFour(A)V_4(A).
\end{equation}

\begin{definition}[Accumulated boundary density]
\begin{equation}
\epsD(A):=\frac{\Efour(A)}{V_3(A)}=\frac{A}{4}\rhoFour(A).
\end{equation}
\end{definition}

\begin{definition}[Radially conjugate response]
\begin{equation}
\PiA(A):=\frac{1}{V_3(A)}\frac{\dd\Efour}{\dd A}
=\rhoFour(A)+\frac{A}{4}\rhoFour'(A).
\end{equation}
\end{definition}

The dimensions are different:
\begin{equation}
[\epsD]=EL^{-3},
\qquad
[\PiA]=EL^{-4}.
\end{equation}
Thus \(\PiA\) is not a second three-dimensional energy density and is not the physical pressure. It is a radial force density, or equivalently an energetic response per unit boundary three-volume. The physical three-dimensional thermodynamic pressure conjugate to \(V_3\) is instead \(\pD=-(A/3)\PiA\).

For homogeneous \(\rhoFour(A)\), the quasilocal density on every fixed \(\Dom_\chi\) agrees with the global one:
\begin{equation}
\frac{\Efour(\Dom_\chi,A)}{V_3(\chi,A)}
=\frac{\rhoFour(A)V_4(\chi,A)}{V_3(\chi,A)}
=\frac{A}{4}\rhoFour(A).
\end{equation}
The independence from \(\chi\) follows here from homogeneity and homothety; it is not a general law for arbitrary fillings.

\section{Pressure and Direct Reconstruction}
\begin{definition}[Adiabatic pressure]
For a one-parameter energy \(\Efour=\Efour(A)\),
\begin{equation}
\pD(A):=-\frac{\dd\Efour}{\dd V_3}
=-\frac{\dd\Efour/\dd A}{\dd V_3/\dd A}.
\end{equation}
Since \(\dd V_3/\dd A=3V_3/A\),
\begin{equation}
\pD=-\frac{A}{3}\PiA
=-\frac{A}{3}\left(\rhoFour+\frac{A}{4}\rhoFour'\right).
\end{equation}
The factor \(A/3\) is essential: \(\PiA\) is the radial conjugate response per boundary three-volume, whereas \(\pD\) is the thermodynamic pressure conjugate to the physical boundary volume \(V_3\).
\end{definition}

\begin{theorem}[Direct background reconstruction]
Let \(A>0\), let \(\Efour(A)\) be a \(C^1\) function on this domain, and assume the round volume law \(V_3(A)=2\pi^2A^3\), so that \(V_3'(A)=3V_3/A\neq0\). Then \(\epsD=\Efour/V_3\) and \(\pD=-\dd\Efour/\dd V_3\) uniquely define an isotropic background fluid satisfying
\begin{equation}
\frac{\dd\epsD}{\dd\tau}+3H(\epsD+\pD)=0.
\end{equation}
\end{theorem}
\begin{proof}
Write \(V=V_3(A)\) and \(E=\Efour(A)\). Since \(V'/V=3/A\), one has
\begin{equation}
\frac{\dd\epsD}{\dd A}
=\frac{E'}{V}-\frac{EV'}{V^2}
=\frac{E'}{V}-\frac{3E}{AV}.
\end{equation}
Moreover,
\begin{equation}
\pD=-\frac{E'}{V'}
=-\frac{AE'}{3V},
\end{equation}
and therefore
\begin{equation}
\frac{3}{A}(\epsD+\pD)
=\frac{3E}{AV}-\frac{E'}{V}.
\end{equation}
Adding the last two equations gives
\begin{equation}
\frac{\dd\epsD}{\dd A}
+\frac{3}{A}(\epsD+\pD)=0.
\end{equation}
Multiplication by \(\dd A/\dd\tau=HA\) gives the continuity identity. Uniqueness follows from the definitions.
\end{proof}

\section{Internal Variables and Exchange}
Let
\begin{equation}
\Efour=\Efour(A,\xi^1,\ldots,\xi^n).
\end{equation}
Pressure is defined at fixed internal variables:
\begin{equation}
\pD:=-\left.\frac{\partial \Efour}{\partial V_3}\right|_{\xi}.
\end{equation}

\begin{proposition}[Exchange term]
Under evolution \(\xi^I(\tau)\),
\begin{equation}
\frac{\dd\epsD}{\dd\tau}+3H(\epsD+\pD)=\QD,
\qquad
\QD=\frac{1}{V_3}\sum_I\frac{\partial\Efour}{\partial\xi^I}\frac{\dd\xi^I}{\dd\tau}.
\end{equation}
\end{proposition}
\begin{proof}
The total derivative is
\begin{equation}
\frac{\dd\Efour}{\dd\tau}
=\frac{\partial\Efour}{\partial A}\frac{\dd A}{\dd\tau}
+\sum_I\frac{\partial\Efour}{\partial\xi^I}\frac{\dd\xi^I}{\dd\tau}.
\end{equation}
The first terms cancel exactly as in the proof of the continuity identity; the remainder is the exchange law.
\end{proof}

Thus \(\QD\) is not arbitrary within this class: it measures the nonadiabatic work of the internal filling variables.

\section{Inverse Homogeneous Representation and the Limit of Background Predictivity}
\begin{theorem}[Inverse background representation]
Let \(\epsD(A)\) be a differentiable background density. Within the homogeneous homothetic filling class there is a unique representing density
\begin{equation}
\rhoFour(A)=\frac{4}{A}\epsD(A),
\end{equation}
such that \(\Efour=\rhoFour V_4=\epsD V_3\). The pressure induced by the filling-work prescription is
\begin{equation}
\pD(A)=-\epsD(A)-\frac{A}{3}\frac{\dd\epsD}{\dd A}.
\end{equation}
Conversely, a prescribed pair \((\epsD,\pD)\) is representable by this adiabatic filling map if and only if it satisfies
\begin{equation}
\pD(A)=-\epsD(A)-\frac{A}{3}\epsD'(A),
\end{equation}
equivalently the homogeneous continuity identity.
\end{theorem}
\begin{proof}
The relation \(\Efour=\epsD V_3\) together with \(V_3/V_4=4/A\) immediately gives the inverse relation. The same volume ratio gives uniqueness. Furthermore,
\begin{equation}
\pD=-\frac{\dd(\epsD V_3)}{\dd V_3}
=-\epsD-V_3\frac{\dd\epsD/\dd A}{\dd V_3/\dd A}
=-\epsD-\frac{A}{3}\epsD',
\end{equation}
which is the induced pressure relation. The final equivalence follows by multiplying \(\epsD'+3(\epsD+\pD)/A=0\) by \(\dd A/\dd\tau\).
\end{proof}

\begin{remark}[Exact scope of inverse representation]
The inverse representation map
\begin{equation}
\epsD(A)\longmapsto \rhoFour(A)=\frac{4}{A}\epsD(A)
\end{equation}
is homogeneous and algebraic. It establishes uniqueness of the representing density only within the adopted homothetic background class and on the \(A\)-domain where the background history is expressed as a differentiable function of \(A\). It does not reconstruct microscopic fields, particles, perturbations, or a covariant matter action.
\end{remark}

\begin{remark}[No inverse reconstruction of perturbations]
The inverse map reconstructs only the homogeneous representative \(\rhoFour(A)=4\epsD(A)/A\). It is not an inverse map from a background history to a unique stress-energy tensor, perturbation theory, sound speed, anisotropic stress, or microscopic action. These require a separate covariant completion.
\end{remark}

Any interpretation of a singular limit as a boundary of inverse physical reconstruction is an additional interpretive principle, not a consequence of the inverse background representation theorem proved here. Predictivity begins when \(\Efour(A,\xi)\) is independently fixed by a symmetry, an action, a compensation condition, a spectral principle, or equations for the internal variables.

\section{General Equation-of-State Formula}
For \(\epsD>0\), define \(w_D=\pD/\epsD\). The inverse pressure relation gives
\begin{equation}
w_D(A)
=-1-\frac13\frac{\dd\ln(\epsD/\varepsilon_*)}{\dd\ln(A/A_*)}
=-\frac43-\frac13\frac{\dd\ln(\rhoFour/\rho_{4,*})}{\dd\ln(A/A_*)}.
\end{equation}
Define the local scaling exponent
\begin{equation}
\beta_{\rm eff}(A):=-\frac{\dd\ln(\rhoFour/\rho_{4,*})}{\dd\ln(A/A_*)}.
\end{equation}
The fixed positive references \(A_*\), \(\varepsilon_*\), and \(\rho_{4,*}\) only make the logarithms dimensionless; all logarithmic derivatives are independent of their numerical values.
Then
\begin{equation}
w_D(A)=\frac{\beta_{\rm eff}(A)-4}{3}.
\end{equation}

\begin{corollary}[Power-law class]
If
\begin{equation}
\rhoFour(A)=\rho_{4,0}\left(\frac{A}{A_0}\right)^{-\beta},
\end{equation}
then
\begin{equation}
\epsD\propto A^{1-\beta},
\qquad
\pD=\frac{\beta-4}{3}\epsD,
\qquad
w_D=\frac{\beta-4}{3}.
\end{equation}
\end{corollary}

\begin{center}
\begin{tabular}{@{}cccc@{}}
\toprule
\(\beta\) & Scaling of \(\epsD\) & \(w_D\) & Background class \\
\midrule
0 & \(a\)      & \(-4/3\) & phantom-like slot \\
1 & \(a^0\)    & \(-1\)   & vacuum slot \\
2 & \(a^{-1}\) & \(-2/3\) & accelerating slot \\
3 & \(a^{-2}\) & \(-1/3\) & boundary slot \\
4 & \(a^{-3}\) & \(0\)    & dust slot \\
5 & \(a^{-4}\) & \(1/3\)  & radiation-like slot \\
\bottomrule
\end{tabular}
\end{center}

For \(\epsD>0\),
\begin{equation}
\epsD+
\pD=\frac{\beta-1}{3}\epsD,
\qquad
\epsD+3\pD=(\beta-3)\epsD.
\end{equation}
Thus \(\beta<1\) violates the null energy condition of the effective fluid, while \(\beta<3\) gives a negative active gravitational density \(\epsD+3\pD\). These background properties do not replace a microscopic stability analysis.

\section{Dimensional Theorem}
\begin{theorem}[Homothetic filling in arbitrary dimension]
Let
\begin{equation}
V_d(A)=C_dA^d,
\qquad
V_{d+1}(A)=C_{d+1}A^{d+1},
\qquad
\frac{\dd V_{d+1}}{\dd A}=V_d,
\end{equation}
so that \(C_d=(d+1)C_{d+1}\). Define
\begin{equation}
E_{d+1}=\rho_{d+1}V_{d+1},
\qquad
\varepsilon_d=\frac{E_{d+1}}{V_d},
\qquad
p_d=-\frac{\dd E_{d+1}}{\dd V_d}.
\end{equation}
Then
\begin{equation}
\varepsilon_d=\frac{A}{d+1}\rho_{d+1},
\qquad
p_d=-\frac{A}{d}\left(\rho_{d+1}+\frac{A}{d+1}\rho_{d+1}'\right).
\end{equation}
If \(\rho_{d+1}\propto A^{-\beta}\), then
\begin{equation}
w_d=\frac{\beta-d-1}{d}.
\end{equation}
\end{theorem}
\begin{proof}
The identity \(\dd V_{d+1}/\dd A=V_d\) gives \(V_{d+1}/V_d=A/(d+1)\), which proves the dimensional density relation. Since \(V_d'/V_d=d/A\),
\begin{equation}
p_d=-\frac{A}{dV_d}\frac{\dd E_{d+1}}{\dd A}.
\end{equation}
Differentiating \(E_{d+1}=\rho_{d+1}C_{d+1}A^{d+1}\) gives the dimensional pressure relation. For a power-law density, \(\rho_{d+1}'=-(\beta/A)\rho_{d+1}\); division by the dimensional density relation yields the dimensional equation-of-state relation.
\end{proof}

The coefficients \(1/4\) and \(4\) in the physical case \(d=3\) are therefore special values of a general dimensional structure.

\section{Linear Superposition and Protected Slots}
\begin{theorem}[Linear superposition]
If \(\Efour(A)=\sum_i E_{4,i}(A)\), then
\begin{equation}
\epsD=\sum_i \varepsilon_{D,i},
\qquad
\pD=\sum_i p_{D,i}.
\end{equation}
\end{theorem}
\begin{proof}
Both maps, \(\Efour\mapsto \Efour/V_3\) and \(\Efour\mapsto -\dd\Efour/\dd V_3\), are linear.
\end{proof}

For
\begin{equation}
\rhoFour(A)=\sum_i C_iA^{-\beta_i},
\end{equation}
one obtains
\begin{equation}
\epsD(A)=\sum_i \frac{C_i}{4}A^{1-\beta_i},
\qquad
\pD(A)=\sum_i \frac{\beta_i-4}{3}\varepsilon_{D,i}(A).
\end{equation}
The two-branch law
\begin{equation}
\rhoFour(A)=\frac{4\varepsilon_q}{A}+\frac{4\varepsilon_wA_0^3}{A^4}
\end{equation}
is especially important. With \(a=A/A_0\),
\begin{equation}
\epsD(a)=\varepsilon_q+\varepsilon_w a^{-3},
\qquad
\pD(a)=-\varepsilon_q.
\end{equation}
Define the diagnostic invariants
\begin{equation}
\mathcal{A}_q(a):=-\pD(a),
\qquad
\mathcal{A}_w(a):=a^3\bigl(\epsD(a)+\pD(a)\bigr).
\end{equation}

\begin{criterion}[Exact vacuum-dust decomposition]
At the background level, the conditions
\begin{equation}
\frac{\dd\mathcal{A}_q}{\dd a}=0,
\qquad
\frac{\dd\mathcal{A}_w}{\dd a}=0
\end{equation}
are necessary and sufficient for the representation with constant \(\varepsilon_q\) and \(\varepsilon_w\).
\end{criterion}
\begin{proof}
Necessity follows by direct substitution. Conversely, constancy of \(\mathcal{A}_q\) gives \(\pD=-\varepsilon_q\), while constancy of \(\mathcal{A}_w\) gives \(\epsD+
\pD=\varepsilon_w a^{-3}\). Adding the two equations yields the decomposition.
\end{proof}

The dust slot is an exact background class, but it is not yet clustering cold dark matter. That identification requires perturbative closure.



\section{Full Valuation Spectrum, Boundary Filter, and Minimal Endpoint Classification}
The universal background map
\begin{equation}
\Efour(A)\longmapsto(\epsD,\pD)
\end{equation}
is algebraically complete but not predictive by itself. Predictivity begins only after a restricted invariant class of filling energies has been selected.

We do not reprove Hadwiger's theorem here \cite{Hadwiger1957,Schneider2014}. We use its intrinsic-volume classification only to motivate the finite valuation basis and then compute its round homothetic specialization explicitly.

\begin{hypothesis}[Euclidean convex valuation closure]
The auxiliary filling energy is assumed to be a continuous, Euclidean-invariant, additive valuation on convex bodies in \(\mathbb{R}^4\), restricted to the round homothetic family \(\Bfour_A\). For the smooth round representatives considered here, this intrinsic-volume span may be written, after absorbing conventional normalizations into the coefficients, as
\begin{equation}
\Efour^{\rm val}(A)
=\gamma V_4(A)
+\sum_{j=0}^{3}\alpha_j
\int_{\Sthree_A}e_j(K)\,\dd\Sigma .
\end{equation}
\end{hypothesis}

The constants \(\gamma,\alpha_j\) absorb the conventional normalizations of the intrinsic volumes and curvature measures. Only the powers of \(A\) and the induced equations of state are invariant under these normalization choices. Here \(K^a{}_b\) is the Weingarten map of the boundary hypersurface and \(e_j(K)\) are the elementary symmetric polynomials of the principal curvatures:
\begin{equation}
e_0=1,
\qquad
e_1=\operatorname{tr}K,
\qquad
e_2=\sum_{a<b}\kappa_a\kappa_b,
\qquad
e_3=\det(K^a{}_b).
\end{equation}
We use the Weingarten-map sign convention for which the outward normal of the convex round filling gives positive principal curvatures. With the opposite convention the odd elementary symmetric terms change sign; this does not change the valuation spectrum, since the coefficients \(\alpha_j\) absorb the corresponding sign convention. For the round homothetic representative,
\begin{equation}
K^a{}_b=\frac{1}{A}\delta^a_b,
\end{equation}
and hence
\begin{equation}
e_j(K)=\binom{3}{j}A^{-j}.
\end{equation}
Under the standard convex-body valuation assumptions of Euclidean invariance, continuity, and additivity, these terms correspond to the intrinsic-volume basis; the explicit formula below is the homothetic round specialization needed here.

\begin{theorem}[Full homothetic valuation spectrum]
Under the Euclidean convex valuation closure hypothesis \cite{Hadwiger1957,Schneider2014}, for the round homothetic filling \(\Bfour_A\),
\begin{equation}
\Efour^{\rm val}(A)
=\frac{\pi^2}{2}\gamma A^4
+2\pi^2\left(
\alpha_0A^3+3\alpha_1A^2+3\alpha_2A+\alpha_3
\right).
\end{equation}
The induced boundary density is
\begin{equation}
\epsD(A)
=\frac{\gamma A}{4}
+\alpha_0+\frac{3\alpha_1}{A}
+\frac{3\alpha_2}{A^2}
+\frac{\alpha_3}{A^3}.
\end{equation}
The pure valuation branches have equations of state
\begin{equation}
\gamma V_4:\quad w=-\frac{4}{3},
\end{equation}
and
\begin{equation}
e_j:\quad w_j=-1+\frac{j}{3}\qquad (j=0,1,2,3).
\end{equation}
Thus the full homothetic valuation spectrum is
\begin{equation}
w=-\frac{4}{3},\ -1,\ -\frac{2}{3},\ -\frac{1}{3},\ 0.
\end{equation}
For the general five-term valuation sector,
\begin{equation}
\pD(A)=-\frac{\gamma A}{3}-\alpha_0-\frac{2\alpha_1}{A}-\frac{\alpha_2}{A^2}.
\end{equation}
\end{theorem}

\begin{proof}
The bulk term is
\begin{equation}
\gamma V_4(A)=\frac{\pi^2}{2}\gamma A^4.
\end{equation}
For the boundary terms, since the unit three-sphere has volume \(2\pi^2\), the round boundary measure is \(\dd\Sigma=A^3\dd\Omega_3\) and
\begin{equation}
\int_{\Sthree_A}\dd\Sigma=2\pi^2A^3.
\end{equation}
Therefore
\begin{equation}
\int_{\Sthree_A}e_j(K)\,\dd\Sigma
=2\pi^2\binom{3}{j}A^{3-j}.
\end{equation}
Division by \(V_3=2\pi^2A^3\) gives the stated density. The equation-of-state formula
\begin{equation}
w=-1-\frac{1}{3}\frac{\dd\ln(\epsD/\varepsilon_*)}{\dd\ln(A/A_*)}
\end{equation}
then gives \(w=-4/3\) for the bulk branch \(\epsD\propto A\), and \(w_j=-1+j/3\) for the boundary branch \(\varepsilon_{D,j}\propto A^{-j}\). Summing the branch pressures gives the stated expression for \(\pD(A)\).
\end{proof}



\subsection{Fixed-coefficient flat-limit status of the valuation branches}
The compactness stress test has a sharp consequence for the valuation interpretation.  Write the local curvature-measure boundary profile in terms of the normalized scale factor \(a=A/A_0\):
\begin{equation}
\varepsilon_{\partial}(a;A_0)
=\alpha_0+\frac{3\alpha_1}{A_0}a^{-1}
+\frac{3\alpha_2}{A_0^2}a^{-2}
+\frac{\alpha_3}{A_0^3}a^{-3}.
\end{equation}

\begin{proposition}[Fixed-coefficient flat-limit decoupling]
Let \(I\subset(0,\infty)\) be compact and let \(\alpha_j\) be independent of \(A_0\). Then, uniformly for \(a\in I\),
\begin{equation}
\lim_{A_0\to\infty}\varepsilon_{\partial}(a;A_0)=\alpha_0.
\end{equation}
More generally, a finite local limit retaining a nonzero \(a^{-j}\) amplitude requires and is characterized by the matched asymptotic condition
\begin{equation}
\frac{\alpha_j(A_0)}{A_0^j}\longrightarrow \bar\alpha_j,
\qquad
0<|\bar\alpha_j|<\infty,
\qquad j=1,2,3,
\end{equation}
or equivalently
\begin{equation}
\alpha_j(A_0)=\bar\alpha_j A_0^j+o(A_0^j).
\end{equation}
If \(\alpha_j(A_0)/A_0^j\to0\), the corresponding branch decouples; if the ratio diverges, the local amplitude diverges.
\end{proposition}
\begin{proof}
For \(a\in I\), the functions \(a^{-1},a^{-2},a^{-3}\) are bounded.  With fixed \(\alpha_j\), the nonconstant terms are bounded by constants times \(A_0^{-1}\), \(A_0^{-2}\), and \(A_0^{-3}\), respectively, and therefore converge uniformly to zero.  If \(\alpha_j=\alpha_j(A_0)\), the coefficient of \(a^{-j}\) in the displayed profile is \(c_j\alpha_j(A_0)/A_0^j\), with fixed nonzero numerical factor \(c_1=3\), \(c_2=3\), and \(c_3=1\).  Thus the branch has a finite nonzero limiting amplitude exactly when \(\alpha_j(A_0)/A_0^j\) has a finite nonzero limit, which is the stated asymptotic form.
\end{proof}

\begin{remark}[Local phenomenology versus compact valuation origin]
The decoupling result does not invalidate local scalar-field or dust phenomenology.  A canonical scalar with \(\rho\propto a^{-1}\), another with \(\rho\propto a^{-2}\), and a conserved dust current all have ordinary local descriptions. What fails in the pointed \(S^3_{A_0}\to\mathbb R^3\) limit at fixed \(\alpha_j\) is the claim that their compact valuation origin remains topology-independent. Nonconstant valuation amplitudes require matching coefficients tied to the compact curvature scale or to independent physics.
\end{remark}

\subsection{Finite dilation algebra, minimal annihilator, and canonical projectors}
The nontrivial closure of the round valuation profile is most sharply expressed through the action of homothetic dilations.  Let
\begin{equation}
a:=\frac{A}{A_0},
\qquad
\Dil:=A\frac{\dd}{\dd A}=a\frac{\dd}{\dd a}.
\end{equation}

\begin{definition}[Dimensionless logarithmic dilation coordinate]
Fix an arbitrary reference length \(A_*>0\) and define
\begin{equation}
x:=\ln\frac{A}{A_*}.
\end{equation}
Then
\begin{equation}
\Dil=A\frac{\dd}{\dd A}=\frac{\dd}{\dd x}.
\end{equation}
All dilation identities are independent of the arbitrary choice of \(A_*\).
\end{definition}
For a finite homothetic profile
\begin{equation}
\varepsilon(A)=\sum_{s\in S}\widehat C_s A^s
\end{equation}
define its dilation support by
\begin{equation}
\operatorname{supp}_{\Dil}(\varepsilon)
:=\{s\in S:\widehat C_s\neq0\}
\end{equation}
and its minimal dilation polynomial by
\begin{equation}
m_{\varepsilon}(z)
:=\prod_{s\in\operatorname{supp}_{\Dil}(\varepsilon)}(z-s).
\end{equation}

\begin{theorem}[Round valuation-profile ODE characterization]
Let \(J\subset(0,\infty)\) be an interval and let \(\varepsilon\in C^5(J)\). Then
\begin{equation}
\boxed{
P_{\rm val}(\Dil)\varepsilon=0,
\qquad
P_{\rm val}(z)=(z-1)z(z+1)(z+2)(z+3),
}
\end{equation}
if and only if
\begin{equation}
\boxed{
\varepsilon(A)
=C_1A+C_0+C_{-1}A^{-1}+C_{-2}A^{-2}+C_{-3}A^{-3}.
}
\end{equation}
Thus the kernel of \(P_{\rm val}(\Dil)\) on \(C^5(J)\) is exactly the round homothetic profile image of the full Euclidean valuation class.
\end{theorem}
\begin{proof}
Fix an arbitrary reference length \(A_*>0\) and put
\begin{equation}
x:=\ln\frac{A}{A_*}.
\end{equation}
Then \(\Dil=A\,\dd/\dd A=\dd/\dd x\), so the displayed equation becomes a constant-coefficient fifth-order ODE with five simple characteristic roots
\begin{equation}
1,\ 0,\ -1,\ -2,\ -3.
\end{equation}
Its solution space is therefore the span of \(e^x,1,e^{-x},e^{-2x},e^{-3x}\). Since \(A=A_*e^x\), the fixed powers of \(A_*\) are absorbed into the integration constants, giving precisely the stated power-law profile in \(A\). The converse is immediate because each basis element is annihilated by the corresponding factor of \(P_{\rm val}(\Dil)\).
\end{proof}

\begin{theorem}[Elementwise minimal dilation polynomial]
For a nonzero valuation profile
\begin{equation}
\varepsilon(A)=\sum_{s\in S_{\rm val}}\widehat C_sA^s,
\end{equation}
define its active dilation support and minimal polynomial by
\begin{equation}
\operatorname{supp}_{\Dil}(\varepsilon)
:=\{s:\widehat C_s\neq0\},
\qquad
m_\varepsilon(z)
:=\prod_{s\in\operatorname{supp}_{\Dil}(\varepsilon)}(z-s).
\end{equation}
Then
\begin{equation}
\boxed{
\operatorname{Ann}_{\Dil}(\varepsilon)
:=\{p\in\mathbb R[z]:p(\Dil)\varepsilon=0\}
=(m_\varepsilon),
}
\end{equation}
and
\begin{equation}
\boxed{m_\varepsilon(z)\mid P_{\rm val}(z).}
\end{equation}
\end{theorem}
\begin{proof}
The functions \(A^s\) with distinct real exponents are linearly independent on every interval in \((0,\infty)\), and \(p(\Dil)A^s=p(s)A^s\). Hence \(p(\Dil)\varepsilon=0\) exactly when \(p(s)=0\) on every active spectral value. The monic polynomial of least degree with those roots is \(m_\varepsilon\), which proves the ideal statement. Because the full round valuation support is
\begin{equation}
S_{\rm val}=\{1,0,-1,-2,-3\},
\end{equation}
its elementwise minimal polynomial divides \(P_{\rm val}\).
\end{proof}

For any subset \(S\subseteq S_{\rm val}\), define
\begin{equation}
\mathcal V_S
:=\ker\!\left[\prod_{s\in S}(\Dil-s)\right]
=\operatorname{span}\{A^s:s\in S\}.
\end{equation}
The full, local curvature-measure boundary, and minimal endpoint density spectra are
\begin{equation}
S_{\rm val}=\{1,0,-1,-2,-3\},
\qquad
S_{\partial}=\{0,-1,-2,-3\},
\qquad
S_{\rm end}=\{0,-3\}.
\end{equation}

\begin{corollary}[Spectral lattice]
For \(S,T\subseteq S_{\rm val}\),
\begin{equation}
\boxed{
\mathcal V_S\cap\mathcal V_T=\mathcal V_{S\cap T},
\qquad
\mathcal V_S+\mathcal V_T=\mathcal V_{S\cup T}.
}
\end{equation}
If \(T\subseteq S\), the canonical projector decomposition gives
\begin{equation}
\boxed{
\mathcal V_S/\mathcal V_T\simeq\mathcal V_{S\setminus T}.
}
\end{equation}
\end{corollary}

\begin{theorem}[Finite semisimple dilation-operator algebra]
Let
\begin{equation}
\mathcal V_{\rm val}:=\mathcal V_{S_{\rm val}}
=\operatorname{span}\{A,A^0,A^{-1},A^{-2},A^{-3}\}
\end{equation}
and define the polynomial representation
\begin{equation}
\varrho:\mathbb R[z]\longrightarrow\operatorname{End}(\mathcal V_{\rm val}),
\qquad
p\longmapsto p(\Dil)|_{\mathcal V_{\rm val}}.
\end{equation}
Then
\begin{equation}
\ker\varrho=(P_{\rm val}),
\end{equation}
so the image algebra is
\begin{equation}
\boxed{
\mathfrak A_{\Dil}
:=\varrho(\mathbb R[z])
\simeq\mathbb R[z]/(P_{\rm val})
\simeq\prod_{s\in S_{\rm val}}\mathbb R
\simeq\mathbb R^5.
}
\end{equation}
The profile space \(\mathcal V_{\rm val}\) is a faithful \(\mathfrak A_{\Dil}\)-module. It is not an algebra under pointwise multiplication.
\end{theorem}
\begin{proof}
The kernel consists of those polynomials vanishing at every dilation eigenvalue in \(S_{\rm val}\), hence is the principal ideal generated by \(P_{\rm val}\). The first isomorphism theorem gives the quotient representation. Since all roots of \(P_{\rm val}\) are simple, the Chinese remainder theorem yields the direct product of five copies of \(\mathbb R\). Faithfulness follows because the quotient has been taken by the full kernel of the action. Pointwise multiplication need not preserve \(\mathcal V_{\rm val}\); for example \(A\cdot A=A^2\notin\mathcal V_{\rm val}\).
\end{proof}

For each \(s\in S_{\rm val}\), define the primitive Lagrange idempotent
\begin{equation}
e_s(z)
:=\prod_{\substack{t\in S_{\rm val}\\t\neq s}}
\frac{z-t}{s-t},
\qquad
\mathsf P_s:=e_s(\Dil)|_{\mathcal V_{\rm val}}.
\end{equation}
Modulo \(P_{\rm val}\), these satisfy
\begin{equation}
e_s^2=e_s,
\qquad
e_se_t=0\ (s\neq t),
\qquad
\sum_s e_s=1,
\end{equation}
and therefore
\begin{equation}
\mathsf P_s\varepsilon=\widehat C_s A^s.
\end{equation}
The projectors are thus the primitive idempotents of the operator algebra \(\mathfrak A_{\Dil}\), not a multiplication algebra of the profiles themselves.

\begin{definition}[Reference-normalization-independent spectral coefficients and fractions]
The reference-normalization-independent coefficient of the \(s\)-branch is
\begin{equation}
\boxed{
\widehat{\mathcal I}_s[\varepsilon]
:=A^{-s}\mathsf P_s\varepsilon
=\widehat C_s.
}
\end{equation}
If instead one writes \(\varepsilon=\sum_s C_s a^s\) with \(a=A/A_0\), then
\begin{equation}
C_s=A_0^s\widehat{\mathcal I}_s.
\end{equation}
Thus \(C_s\) is a coordinate tied to the reference normalization \(A_0\), whereas \(\widehat{\mathcal I}_s\) is independent of that reference normalization. The latter is still dimensionful, so its numerical value changes under a change of physical units. The truly dimensionless spectral invariants are the fractions \(\pi_s\). Whenever \(\varepsilon\neq0\), define the dimensionless spectral fractions
\begin{equation}
\boxed{
\pi_s(A):=\frac{\mathsf P_s\varepsilon(A)}{\varepsilon(A)},
\qquad
\sum_s\pi_s=1.
}
\end{equation}
\end{definition}

\begin{corollary}[Local curvature-measure boundary obstruction]
A round profile lies in the local curvature-measure boundary subspace exactly when
\begin{equation}
\boxed{
\Dil(\Dil+1)(\Dil+2)(\Dil+3)\varepsilon=0.
}
\end{equation}
Inside the full valuation profile this is equivalent to
\begin{equation}
\widehat{\mathcal I}_1=0,
\end{equation}
which removes the bulk dilation eigenvalue \(s=1\).
\end{corollary}

\begin{corollary}[Minimal endpoint density obstruction]
A density has the minimal endpoint background form
\begin{equation}
\varepsilon(a)=\varepsilon_q+\varepsilon_w a^{-3}
\end{equation}
if and only if
\begin{equation}
\boxed{
\Dil(\Dil+3)\varepsilon=0.
}
\end{equation}
Inside the full round valuation profile this is equivalent to
\begin{equation}
\widehat{\mathcal I}_{1}=0,
\qquad
\widehat{\mathcal I}_{-1}=0,
\qquad
\widehat{\mathcal I}_{-2}=0.
\end{equation}
The identification of this density with a functional
\begin{equation}
E_{\rm end}=\sigma V_3+\mu\QU
\end{equation}
requires the separate minimal endpoint assumptions below; it does not follow from the dilation ODE alone.
\end{corollary}

\subsection{Positive valuation mixtures: moments, convexity, and rank tests}
Assume on an interval that every active branch has nonnegative energy,
\begin{equation}
\mathsf P_s\varepsilon\ge0,
\qquad
\varepsilon>0.
\end{equation}
Then the fractions \(\pi_s\) form a finite probability distribution on the dilation spectrum. Define its moments
\begin{equation}
\mathfrak m_n(A)
:=\frac{\Dil^n\varepsilon}{\varepsilon}
=\sum_s s^n\pi_s.
\end{equation}

\begin{theorem}[Log-convexity and monotone equation of state]
For a positive valuation mixture,
\begin{equation}
\boxed{
\Dil\mathfrak m_1
=\mathfrak m_2-\mathfrak m_1^2
=\operatorname{Var}_{\pi}(s)\ge0,
}
\end{equation}
and therefore
\begin{equation}
\boxed{
\Dil^2\ln(\varepsilon/\varepsilon_*)
=\operatorname{Var}_{\pi}(s)\ge0.
}
\end{equation}
Since
\begin{equation}
w=-1-\frac{\mathfrak m_1}{3},
\end{equation}
one also has
\begin{equation}
\boxed{
\Dil w=-\frac13\operatorname{Var}_{\pi}(s)\le0.
}
\end{equation}
Moreover,
\begin{equation}
\Dil\pi_s=(s-\mathfrak m_1)\pi_s.
\end{equation}
Hence a positive valuation mixture cannot possess an interior maximum of \(\varepsilon(A)\); it may possess a minimum when \(\mathfrak m_1\) crosses zero from negative to positive.
\end{theorem}
\begin{proof}
Differentiate \(\mathfrak m_1=(\Dil\varepsilon)/\varepsilon\) and use \(\mathfrak m_2=(\Dil^2\varepsilon)/\varepsilon\). The variance identity follows immediately. The equations for \(\ln(\varepsilon/\varepsilon_*)\), \(w\), and \(\pi_s\) are then algebraic consequences; the arbitrary fixed \(\varepsilon_*>0\) drops out of all dilation derivatives.
\end{proof}

\begin{theorem}[Spectral log-partition hierarchy and sharp diameter bound]
Fix \(A_*>0\) and a positive normalization \(\varepsilon_*>0\), and write
\begin{equation}
x=\ln\frac{A}{A_*},
\qquad
Z(x):=\frac{\varepsilon(A_*e^x)}{\varepsilon_*},
\qquad
K(x):=\ln Z(x).
\end{equation}
For a positive finite dilation mixture,
\begin{equation}
\varepsilon(A)=\sum_s\widehat C_sA^s,
\qquad
\widehat C_sA^s\ge0,
\end{equation}
define
\begin{equation}
\pi_s(x):=
\frac{\widehat C_sA_*^s e^{sx}}
{\sum_t\widehat C_tA_*^t e^{tx}}.
\end{equation}
Then \(K\) is the log-partition function of this exponentially tilted finite spectral measure and
\begin{equation}
\boxed{
K^{(n)}(x)=\kappa_n^{(x)}[s],
}
\end{equation}
where the right-hand side is the \(n\)-th cumulant with respect to \(\pi_s(x)\). In particular,
\begin{equation}
K'=\langle s\rangle_\pi,
\qquad
K''=\operatorname{Var}_\pi(s)\ge0.
\end{equation}
If the active spectral support lies in \([s_{\min},s_{\max}]\), then
\begin{equation}
0\le\operatorname{Var}_\pi(s)
\le\frac{(s_{\max}-s_{\min})^2}{4},
\end{equation}
and hence
\begin{equation}
-\frac{(s_{\max}-s_{\min})^2}{12}
\le\Dil w\le0.
\end{equation}
Consequently,
\begin{align}
-\frac43&\le\Dil w\le0
&&\text{for }S_{\rm val}\subset[-3,1],\\
-\frac34&\le\Dil w\le0
&&\text{for }S_{\partial}\subset[-3,0],\\
-\frac1{12}&\le\Dil w\le0
&&\text{for the positive }\Lambda+X\text{ spectrum }\{-1,0\}.
\end{align}
At \(a=1\), the last inequality is equivalent to \(0\le w_a\le1/12\).
\end{theorem}
\begin{proof}
Writing \(b_s=\widehat C_sA_*^s\), the weights \(\pi_s(x)\propto b_se^{sx}\) form the exponential tilt of a fixed finite spectral measure. Derivatives of the normalized log-partition function therefore generate the cumulants of the tilted measure at the same \(x\). The arbitrary factor \(\varepsilon_*\) adds only a constant to \(K\) and disappears from every derivative. The variance bound is Popoviciu's sharp bound for a distribution supported on a compact interval; substituting \(\Dil w=-\operatorname{Var}_\pi(s)/3\) gives the stated dilation bounds.
\end{proof}


\begin{theorem}[Bhatia--Davis dilation envelope]
Using the Bhatia--Davis variance bound \cite{BhatiaDavis2000}, let a positive dilation mixture have active spectrum contained in \([m,M]\), and let
\begin{equation}
\bar s:=\sum_s s\pi_s=-3(1+w).
\end{equation}
Then
\begin{equation}
\boxed{
0\le-\Dil w
=\frac13\operatorname{Var}_\pi(s)
\le
\frac13(M-\bar s)(\bar s-m).
}
\end{equation}
Equivalently,
\begin{equation}
\boxed{
0\le-\Dil w
\le
\frac13\bigl[M+3(1+w)\bigr]
\bigl[-3(1+w)-m\bigr].
}
\end{equation}
Equality in the upper bound holds if and only if all active spectral weight is supported on the two endpoints \(\{m,M\}\).
\end{theorem}
\begin{proof}
For every spectral value \(s\in[m,M]\),
\begin{equation}
(M-s)(s-m)\ge0.
\end{equation}
Taking the \(\pi\)-expectation gives
\begin{equation}
\langle s^2\rangle\le(M+m)\bar s-mM.
\end{equation}
Subtracting \(\bar s^2\) yields
\begin{equation}
\operatorname{Var}_\pi(s)\le(M-\bar s)(\bar s-m).
\end{equation}
Equality requires \((M-s)(s-m)=0\) on every active spectral point, hence support only at \(m\) and \(M\).
\end{proof}

\begin{corollary}[Endpoint saturation theorem]
For every positive local curvature-measure boundary mixture, \(S_\partial\subset[-3,0]\),
\begin{equation}
\boxed{
0\le-\Dil w\le-3w(1+w),
\qquad -1\le w\le0.
}
\end{equation}
The upper bound is saturated identically if and only if the active spectrum is contained in \(\{-3,0\}\), i.e. the positive profile belongs to the minimal endpoint density subspace. Equivalently,
\begin{equation}
\boxed{
\Dil w=3w(1+w)
}
\end{equation}
is the first-order equation-of-state closure of the positive endpoint sector.
\end{corollary}

\begin{corollary}[Exact positive \(\Lambda+X\) flow]
For the positive two-branch spectrum \(\{-1,0\}\),
\begin{equation}
\boxed{
\Dil w=(1+w)(2+3w),
\qquad
-1\le w\le-\frac23.
}
\end{equation}
At \(a=1\),
\begin{equation}
w_a=-\Dil w
=(1+w_0)(-2-3w_0),
\end{equation}
so the local CPL curve is the exact two-point spectral saturation law rather than an isolated Taylor-series coincidence.
\end{corollary}

\begin{corollary}[Single-transition ordering]
For any positive finite dilation mixture, \(\bar s=\Dil\ln(\varepsilon/\varepsilon_*)\) is nondecreasing and \(w\) is nonincreasing with \(x=\ln(A/A_*)\). Hence every fixed equation-of-state threshold can be crossed at most once unless the solution is identically on that threshold. In particular, the positive boundary class satisfies \(w\ge-1\) and never becomes phantom. In the full positive valuation class, a crossing of \(w=-1\), if it occurs, can only be from \(w>-1\) to \(w<-1\) as the universe expands.
\end{corollary}

\begin{theorem}[Projective equation-of-state closure]
Let \(\varepsilon>0\) and let
\begin{equation}
P(z)=\sum_{n=0}^{N}p_nz^n.
\end{equation}
Define
\begin{equation}
Y:=\Dil\ln(\varepsilon/\varepsilon_*)=-3(1+w),
\qquad
\mathcal B_0:=1,
\qquad
\mathcal B_{n+1}:=(\Dil+Y)\mathcal B_n.
\end{equation}
Then
\begin{equation}
\frac{\Dil^n\varepsilon}{\varepsilon}=\mathcal B_n
\end{equation}
and therefore
\begin{equation}
\boxed{
P(\Dil)\varepsilon=0
\quad\Longleftrightarrow\quad
\sum_{n=0}^{N}p_n\mathcal B_n=0.
}
\end{equation}
Thus every linear dilation closure of order \(N\) induces an amplitude-independent nonlinear equation-of-state closure of order \(N-1\). For the endpoint annihilator \(P(z)=z(z+3)\) this reduces to
\begin{equation}
\Dil w=3w(1+w),
\end{equation}
while for the \(\Lambda+X\) annihilator \(P(z)=z(z+1)\) it reduces to
\begin{equation}
\Dil w=(1+w)(2+3w).
\end{equation}
\end{theorem}
\begin{proof}
The claim follows by induction. For \(n=0\),
\(\Dil^0\varepsilon/\varepsilon=1=\mathcal B_0\). If
\(\Dil^n\varepsilon=\varepsilon\mathcal B_n\), then
\begin{equation}
\frac{\Dil^{n+1}\varepsilon}{\varepsilon}
=
\frac{\Dil(\varepsilon\mathcal B_n)}{\varepsilon}
=
\Dil\mathcal B_n+
\frac{\Dil\varepsilon}{\varepsilon}\mathcal B_n
=
(\Dil+Y)\mathcal B_n
=
\mathcal B_{n+1}.
\end{equation}
Hence
\begin{equation}
\frac{P(\Dil)\varepsilon}{\varepsilon}
=
\sum_{n=0}^{N}p_n\mathcal B_n,
\end{equation}
and \(\varepsilon>0\) gives the equivalence.
\end{proof}

\begin{corollary}[Positive-mixture CPL sign test]
For a local CPL tangent \(w(a)=w_0+w_a(1-a)+\cdots\),
\begin{equation}
\Dil w\big|_{a=1}=-w_a.
\end{equation}
Therefore every positive valuation mixture obeys
\begin{equation}
\boxed{w_a\ge0}
\end{equation}
at the present epoch. In particular, a phantom-to-nonphantom crossing with increasing \(a\), which would require \(\Dil w>0\) at the crossing, is impossible inside the positive-coefficient valuation class.
\end{corollary}

\begin{theorem}[Hankel positivity and active-branch rank]
For \(r\ge0\), define the dilation-moment matrix
\begin{equation}
H_r(A):=
\left[\Dil^{i+j}\varepsilon(A)\right]_{i,j=0}^{r}.
\end{equation}
For a positive valuation mixture,
\begin{equation}
\boxed{H_r(A)\succeq0.}
\end{equation}
If exactly \(k\) distinct dilation branches are active, then
\begin{equation}
\operatorname{rank}H_r=\min\{r+1,k\};
\end{equation}
in particular, for \(r\ge k-1\) the rank equals the number of active branches.
\end{theorem}
\begin{proof}
Writing \(v_s=(1,s,\ldots,s^r)^T\),
\begin{equation}
H_r=\sum_s (\widehat C_sA^s)\,v_sv_s^T,
\end{equation}
a positive sum of rank-one matrices. Distinct spectral values give a Vandermonde family, which yields the rank statement.
\end{proof}

\begin{remark}[Invariant content of the equation of state]
The equation-of-state values are images of the dilation eigenvalues,
\begin{equation}
s\longmapsto w_s=-1-\frac{s}{3},
\end{equation}
so that
\begin{equation}
s=1,0,-1,-2,-3
\quad\Longleftrightarrow\quad
w=-\frac43,-1,-\frac23,-\frac13,0.
\end{equation}
The valuation spectrum is therefore a finite spectral statement for \(\Dil\), not merely a list of fitted powers.
\end{remark}

\begin{closurecondition}[Local curvature-measure boundary restriction]
Assume that the physical filling energy is required to be expressible as a local curvature-measure functional of the boundary relative to the declared Euclidean auxiliary filling,
\begin{equation}
\partial\Bfour_A=\Sthree_A,
\end{equation}
built from the boundary curvature measures of the homothetic convex filling. The local curvature-measure boundary class used here means filling-relative local scalar densities built only from the induced boundary measure and the elementary symmetric curvatures \(e_j(K)\), with coefficients independent of \(A\), and not boundary integrals involving the support function \(x\cdot n\) or other ambient-position data. Then the admissible valuation sector excludes the bulk intrinsic-volume term:
\begin{equation}
\gamma=0.
\end{equation}
It leaves the four boundary curvature slots
\begin{equation}
w=-1,\ -\frac{2}{3},\ -\frac{1}{3},\ 0.
\end{equation}
\end{closurecondition}
\begin{proof}
In the full valuation class the term
\begin{equation}
\gamma V_4(A)=\frac{\pi^2}{2}\gamma A^4
\end{equation}
is the four-dimensional bulk intrinsic-volume valuation. The boundary curvature valuations have homothetic powers
\begin{equation}
\int_{\Sthree_A}e_j(K)\,\dd\Sigma\propto A^{3-j},
\qquad j=0,1,2,3.
\end{equation}
Therefore the bulk \(A^4\) branch is not a member of the boundary curvature valuation sector. Although the Minkowski formula can represent the bulk volume as a boundary integral involving the support function, \(V_4=\tfrac14\int_{\partial K}x\cdot n\,\dd\Sigma\), that ambient-position datum is excluded by the declared local curvature-measure class. Hence the \(A^4\) branch is absent and \(\gamma=0\) relative to this restriction, not from topology alone.
\end{proof}


\subsection{Expansion-rate quotient, geometric curvature splitting, and VDE coordinates}
The local curvature-measure boundary profile contains the four density powers
\begin{equation}
\mathcal V_{\partial}=\operatorname{span}\{1,a^{-1},a^{-2},a^{-3}\}.
\end{equation}
At the level of the expansion rate alone, the curved-\(\Lambda\)CDM nuisance subspace is
\begin{equation}
\mathcal V_{\Lambda Km}^{(E)}
:=\operatorname{span}\{1,a^{-2},a^{-3}\}.
\end{equation}

\begin{theorem}[Expansion-rate quotient of boundary valuation VDE]
For the function \(\mathcal H^2(a)\) alone,
\begin{equation}
\boxed{
\mathcal V_{\partial}/\mathcal V_{\Lambda Km}^{(E)}
\simeq\operatorname{span}\{[a^{-1}]\},
\qquad
\dim=1.
}
\end{equation}
Thus the unique non-standard dilation direction in the expansion-rate quotient is \(a^{-1}\), with
\begin{equation}
w_X=-\frac23.
\end{equation}
\end{theorem}
\begin{proof}
The powers \(a^0,a^{-2},a^{-3}\) coincide algebraically in \(\mathcal H^2(a)\) with vacuum, curvature, and dust scalings. The remaining basis vector \(a^{-1}\) is linearly independent of those directions on every interval in \((0,\infty)\).
\end{proof}

This theorem is deliberately restricted to the expansion-rate function. Let
\begin{equation}
\Omega_K^{\rm geom}
:=-\frac{c^2}{A_0^2H_0^2}<0
\end{equation}
denote the genuine geometric curvature parameter of the closed \(\Sthree\) metric, and let \(\Omega_{2,\rm val}\) denote a physical valuation component with \(w=-1/3\). The physically separated background equation is then
\begin{equation}
\mathcal H^2(a)=
\Omega_ra^{-4}
+\Omega_ma^{-3}
+\bigl(\Omega_K^{\rm geom}+\Omega_{2,\rm val}\bigr)a^{-2}
+\Omega_Xa^{-1}
+\Omega_\Lambda.
\end{equation}
The same coefficient of \(a^{-2}\) appears in \(\mathcal H^2\), but only \(\Omega_K^{\rm geom}\) determines the spatial geodesic-distance relation and the compact harmonic spectrum.

Define the dimensionless transverse comoving distance
\begin{equation}
d(z):=\frac{H_0D_M(z)}{c}.
\end{equation}
For an FLRW geometry the Clarkson--Bassett--Lu consistency relation gives \cite{ClarksonBassettLu2008}
\begin{equation}
\boxed{
\Omega_K^{\rm geom}(z)
=\frac{[\mathcal H(z)d'(z)]^2-1}{d(z)^2}
=\text{constant},
\qquad z>0,
}
\end{equation}
with the value at \(z=0\) understood through the regular \(z\to0\) limit. This relation separates geometric curvature from a matter component that merely shares the \(a^{-2}\) scaling in the Friedmann equation. In an observational analysis, \(d'(z)\) should not be estimated by raw finite differences; \(H(z)\), \(D_M(z)\), and \(\Omega_K^{\rm geom}\) should be fitted jointly through the full FLRW distance relation and its covariance.

The total spectral coefficient of the \(a^{-2}\) term in \(\mathcal H^2\) is extracted by
\begin{equation}
\boxed{
\mathcal C_{-2}[\mathcal H^2]
:=\frac{a^2}{4}\,
\Dil(\Dil+1)(\Dil+3)(\Dil+4)\mathcal H^2
=\Omega_K^{\rm geom}+\Omega_{2,\rm val}.
}
\end{equation}
Hence the curvature-splitting invariant is
\begin{equation}
\boxed{
\Omega_{2,\rm val}
=\mathcal C_{-2}[\mathcal H^2]-\Omega_K^{\rm geom}.
}
\end{equation}

\begin{corollary}[Physical parameter count]
Relative to curved \(\Lambda\)CDM, the general local curvature-measure boundary valuation sector is, at the level of background geometry plus expansion, a two-parameter extension generated by
\begin{equation}
(\Omega_X,\Omega_{2,\rm val}).
\end{equation}
The one-parameter VDE model is the nested subclass
\begin{equation}
\boxed{\Omega_{2,\rm val}=0,}
\end{equation}
or, equivalently, the one-dimensional quotient only when one works with \(\mathcal H^2(a)\) before separating geometric curvature from a physical \(w=-1/3\) component.
\end{corollary}



\subsection{Polynomial background trace and perturbative lifts}
\label{sec:polynomial-lifts}

The dilation module fixes background weights.  It does not, by itself, fix a unique perturbative stress tensor.  To avoid this category error, we distinguish the homogeneous polynomial branch from a chosen \emph{perturbative lift}.

\begin{theorem}[Homogeneous trace law for a reduced energy branch]
Let a reduced homogeneous energy branch scale as
\begin{equation}
E_q(A)=C_qA^q.
\end{equation}
Then
\begin{equation}
\rho_q(A)=\frac{E_q}{V_3}\propto A^{q-3},
\qquad
p_q(A)=-\frac{\dd E_q}{\dd V_3}=-\frac q3\rho_q(A),
\end{equation}
and hence
\begin{equation}
\boxed{w_q=-\frac q3.}
\end{equation}
\end{theorem}
\begin{proof}
Since \(V_3=2\pi^2A^3\), one has \(\rho_q=C_qA^q/V_3\propto A^{q-3}\).  Moreover
\begin{equation}
p_q=-\frac{\dd E_q/\dd A}{\dd V_3/\dd A}
=-\frac{qC_qA^{q-1}}{3V_3/A}
=-\frac q3\frac{C_qA^q}{V_3}
=-\frac q3\rho_q.
\end{equation}
\end{proof}

The reduced-energy exponent \(q\) and the density dilation exponent \(s\) are related by
\begin{equation}
\boxed{s=q-3.}
\end{equation}
Thus
\begin{equation}
\begin{array}{rclcrcl}
q=2 &\Longleftrightarrow& s=-1, &\qquad& w&=&-\frac23,\\[1mm]
q=1 &\Longleftrightarrow& s=-2, &\qquad& w&=&-\frac13,\\[1mm]
q=0 &\Longleftrightarrow& s=-3, &\qquad& w&=&0.
\end{array}
\end{equation}

This is a background theorem.  It is not a four-dimensional Ward identity and it does not determine a unique \(\delta T_{\mu\nu}\).  The same background branch may admit inequivalent lifts.  For example, the \(a^{-1}\) VDE background can be realized by the canonical scalar lift, with
\begin{equation}
c_{s,\rm rest}^2=1,
\qquad
\pi_{\mu\nu}=0,
\end{equation}
or by a geometric ADM-response prototype whose scalar response has anisotropic stress and higher spatial derivatives.  Thus
\begin{equation}
\boxed{
\begin{array}{c}
\mathfrak A_{\Dil}\ \text{fixes finite background weights,}\\[1mm]
\text{a chosen lift fixes }\delta T,\quad
\text{and numerical codes test the lift,}\\[1mm]
\text{not the valuation spectrum alone.}
\end{array}}
\end{equation}

The pressure-kernel branch \(q=0\) illustrates the distinction sharply.  The homogeneous energy \(E_0=C_0\) gives \(\rho_0=C_0/(2\pi^2A^3)\) and \(p_0=0\), but the minisuperspace action \(S_0=-C_0\int N\,\dd t\) does not supply a cold-dark-matter hierarchy.  The CDM-like perturbative lift is the conserved-current dust source, with
\begin{equation}
T_{\mu\nu}^{\rm dust}=\rho u_\mu u_\nu,
\qquad
\nabla_\mu(\rho u^\mu)=0,
\qquad
u^\mu\nabla_\mu u^\nu=0,
\qquad
\delta p=0,
\qquad
\pi^i{}_j=0,
\end{equation}
where the global rest-energy matching \(\mathcal E_{\rm rest}\overset{\rm match}{=}\mu\QU\) is imposed as matching data.  It should not be double-counted with baryons.

\subsection{ADM response prototype: intrinsic curvature branch}
\label{sec:ADM-curvature-response}

The physical \(s=-2\), \(w=-1/3\), branch can be represented at the
spatial-response level by an intrinsic ADM curvature functional. This subsection
is a response prototype and diagnostic, not a foundational part of the finite
dilation-algebra theorem. It does not construct the full valuation functional
\(S_{\rm val}\), does not replace the independent canonical-scalar realization,
and does not define an autonomous covariantly conserved matter fluid. It isolates
the part of the boundary valuation spectrum that is expressible by the intrinsic
three-curvature of the spatial slice.

Consider the ADM metric
\begin{equation}
\dd s^2=-N^2\dd\eta^2+h_{ij}\dd x^i\dd x^j,
\end{equation}
and define
\begin{equation}
\boxed{
S_R:=-\lambda_R\int \dd\eta\,\dd^3x\,N\sqrt h\,{}^{(3)}R[h].
}
\end{equation}
Here \(\lambda_R\) is a dimensional matching coefficient. This term is a
spatial curvature-response functional. By itself it is not a complete
four-dimensional diffeomorphism-invariant matter action.

\begin{theorem}[Background intrinsic-curvature branch]
On the round closed FLRW representative,
\begin{equation}
h_{ij}=A^2\gamma_{ij},
\qquad
{}^{(3)}R=\frac{6}{A^2},
\qquad
V_3=2\pi^2A^3,
\end{equation}
the ADM curvature-response energy is
\begin{equation}
E_R(A)=\lambda_R\int_{\Sthree_A}\sqrt h\,{}^{(3)}R
=12\pi^2\lambda_R A.
\end{equation}
Therefore
\begin{equation}
\boxed{
\rho_R=\frac{E_R}{V_3}=\frac{6\lambda_R}{A^2},
\qquad
p_R=-\frac{\dd E_R}{\dd V_3}
=-\frac{2\lambda_R}{A^2},
\qquad
w_R=-\frac13 .
}
\end{equation}
Thus the intrinsic ADM curvature response realizes the homogeneous
\(s=-2\) valuation branch.
\end{theorem}

\begin{proof}
The round three-curvature and volume give
\begin{equation}
\int_{\Sthree_A}\sqrt h\,{}^{(3)}R
=\frac{6}{A^2}\,2\pi^2A^3
=12\pi^2A .
\end{equation}
Division by \(V_3=2\pi^2A^3\) gives \(\rho_R=6\lambda_R/A^2\). Since
\(\dd V_3/\dd A=6\pi^2A^2\), one has
\begin{equation}
p_R=-\frac{\dd E_R/\dd A}{\dd V_3/\dd A}
=-\frac{12\pi^2\lambda_R}{6\pi^2A^2}
=-\frac{2\lambda_R}{A^2}.
\end{equation}
Hence \(p_R/\rho_R=-1/3\).
\end{proof}

The spatial metric variation gives
\begin{equation}
\delta\!\int N\sqrt h\,{}^{(3)}R
=
\int\sqrt h
\left[
N\,{}^{(3)}G_{ij}
+h_{ij}\Delta N
-D_iD_jN
\right]\delta h^{ij},
\end{equation}
and therefore the ADM spatial-response tensor is
\begin{equation}
\boxed{
\tau^{(R)}_{ij}
=
2\lambda_R
\left[
{}^{(3)}G_{ij}
+\frac{1}{N}
\left(
h_{ij}\Delta N-D_iD_jN
\right)
\right].
}
\end{equation}
For homogeneous \(N\), this reduces to
\begin{equation}
\tau^{(R)}_{ij}
=
2\lambda_R\,{}^{(3)}G_{ij}
=-\frac{2\lambda_R}{A^2}h_{ij},
\end{equation}
which reproduces the pressure found above.

\begin{theorem}[Scalar linear response of the intrinsic curvature branch]
Let the scalar perturbations be written in Newtonian form
\begin{equation}
N=A(1+\Psi),
\qquad
h_{ij}=A^2(1-2\Phi)\gamma_{ij},
\end{equation}
where \(\gamma_{ij}\) is the unit round \(\Sthree\) metric and
\(\bar D_i\), \(\bar\Delta\) are its covariant derivative and Laplacian. Then
\begin{equation}
{}^{(3)}R
=
\frac{1}{A^2}
\left(
6+12\Phi+4\bar\Delta\Phi
\right)
\end{equation}
to first order, and the curvature-response density perturbation is
\begin{equation}
\boxed{
\delta\rho_R
=
\frac{4\lambda_R}{A^2}
\left(3+\bar\Delta\right)\Phi .
}
\end{equation}
The mixed spatial response is
\begin{equation}
\boxed{
\delta\tau^{i}{}_{j,R}
=
\frac{2\lambda_R}{A^2}
\left[
\bar D^i\bar D_j(\Phi-\Psi)
-\delta^i{}_j\bar\Delta(\Phi-\Psi)
-2\Phi\,\delta^i{}_j
\right].
}
\end{equation}
Consequently the isotropic pressure perturbation is
\begin{equation}
\boxed{
\delta p_R
=
-\frac{4\lambda_R}{A^2}\Phi
-\frac{4\lambda_R}{3A^2}\bar\Delta(\Phi-\Psi),
}
\end{equation}
or equivalently
\begin{equation}
\boxed{
\delta p_R
=
-\frac13\delta\rho_R
+\frac{4\lambda_R}{3A^2}\bar\Delta\Psi .
}
\end{equation}
The traceless anisotropic response is
\begin{equation}
\boxed{
\Pi^{i}{}_{j,R}
=
\frac{2\lambda_R}{A^2}
\left(
\bar D^i\bar D_j
-\frac13\delta^i{}_j\bar\Delta
\right)(\Phi-\Psi).
}
\end{equation}
\end{theorem}

\begin{proof}
For a conformal perturbation
\begin{equation}
h_{ij}=A^2e^{-2\Phi}\gamma_{ij},
\end{equation}
the first-order scalar curvature is
\begin{equation}
{}^{(3)}R
=A^{-2}\left(6+12\Phi+4\bar\Delta\Phi\right).
\end{equation}
The density perturbation follows from the lapse variation of \(S_R\). The
spatial response follows from the variation formula above with
\(N=A(1+\Psi)\). Taking the trace of \(\delta\tau^i{}_{j,R}\) gives
\(\delta p_R=\frac13\delta\tau^i{}_{i,R}\). Subtracting this trace part gives the
displayed traceless response.
\end{proof}

For scalar harmonics on the unit \(\Sthree\),
\begin{equation}
\bar\Delta Y_{n\ell m}=-\kappa_nY_{n\ell m},
\qquad
\kappa_n=n(n+2),
\end{equation}
the response becomes
\begin{align}
\boxed{
\delta\rho_{R,n}
=
\frac{4\lambda_R}{A^2}
(3-\kappa_n)\Phi_n,
}
\end{align}
\begin{align}
\boxed{
\delta p_{R,n}
=
-\frac{4\lambda_R}{A^2}\Phi_n
+
\frac{4\lambda_R}{3A^2}\kappa_n(\Phi_n-\Psi_n),
}
\end{align}
\begin{align}
\boxed{
\Pi^{i}{}_{j,R,n}
=
\frac{2\lambda_R}{A^2}
\left(
\bar D^i\bar D_j
+\frac13\kappa_n\delta^i{}_j
\right)(\Phi_n-\Psi_n).
}
\end{align}

\begin{remark}[Response ratio is not a rest-frame sound speed]
The formal ratio
\begin{equation}
\frac{\delta p_{R,n}}{\delta\rho_{R,n}}
=
\frac{
-\Phi_n+\frac{\kappa_n}{3}(\Phi_n-\Psi_n)
}{
(3-\kappa_n)\Phi_n
}
\end{equation}
is a metric-response ratio in Newtonian gauge. It is not a rest-frame
microphysical sound speed. In particular, the \(n=1\) scalar sector has
\(\kappa_1=3\), so the density response vanishes in this expression and the
ratio is not a physical sound-speed diagnostic.
\end{remark}

\begin{remark}[ADM response versus conserved matter source]
The functional \(S_R\) is spatially diffeomorphism invariant and gives a
well-defined ADM curvature response. It does not by itself define a complete
four-dimensional covariantly conserved matter stress tensor. Therefore the
response may be used consistently in one of two ways:
\begin{enumerate}
\item as a modified-gravity or geometric-response term kept on the left-hand
side of the field equations; or
\item as an effective boundary response supplemented by whatever compensating
terms are required to restore the full four-dimensional Ward identity.
\end{enumerate}
It should not be inserted into a Boltzmann hierarchy as an ordinary conserved
fluid with an independently meaningful rest-frame \(c_s^2\) unless such a
four-dimensional completion is supplied.
\end{remark}

Since the displayed ADM potential contains no shift dependence, its momentum-constraint response vanishes in the preferred ADM slicing:
\begin{equation}
\boxed{\delta q^{(R)}_{\rm ADM}=0.}
\end{equation}
This is a preferred-foliation response statement, not a covariant matter-frame statement.

\begin{criterion}[Curvature-response implementation status]
A CLASS or CLASS-like implementation of the intrinsic curvature-response branch
must declare whether the above terms are treated as modified-gravity terms or
as an effective stress tensor with compensating conservation equations. The
closed-form kernels above are suitable for a transfer diagnostic of the
\(s=-2\) branch, but they do not constitute full compact-SVT closure for the
complete valuation sector.
\end{criterion}

\subsection[Square-root ADM response prototype for the a-to-minus-one branch]{Square-root ADM response prototype for the \(a^{-1}\) branch}
\label{sec:sqrtR-response}

The \(a^{-1}\) background branch has a clean canonical scalar lift in this paper.  There is also a distinct geometric ADM-response prototype with the same homogeneous scaling.  These lifts must not be mixed: the canonical scalar has \(c_{s,\rm rest}^2=1\) and \(\pi_{\mu\nu}=0\), while the geometric response below is a foliation-dependent spatial response kernel with anisotropic stress.

This prototype shares the \(a^{-1}\) homothetic scaling of the VDE branch, but it is not itself a Hadwiger valuation functional.  It is an optional foliation-dependent ADM response kernel.  Its domain is restricted to \({}^{(3)}R>0\), since \(F_R=1/(2\sqrt{{}^{(3)}R})\) is singular at \({}^{(3)}R=0\).

Consider
\begin{equation}
S_{\sqrt R}:=-\lambda_2\int \dd\eta\,\dd^3x\,N\sqrt h\,\sqrt{{}^{(3)}R[h]}.
\end{equation}
On the round background,
\begin{equation}
{}^{(3)}R=\frac{6}{A^2},
\qquad
\int_{\Sthree_A}\sqrt h\sqrt{{}^{(3)}R}=2\pi^2\sqrt6\,A^2,
\end{equation}
and therefore
\begin{equation}
\rho_{\sqrt R}=\frac{\sqrt6\lambda_2}{A},
\qquad
p_{\sqrt R}=-\frac23\rho_{\sqrt R}.
\end{equation}
This realizes the same homogeneous \(a^{-1}\), \(w=-2/3\), scaling as the nested VDE branch, but with a different perturbative lift.

Let
\begin{equation}
F=\sqrt{{}^{(3)}R},
\qquad
F_R=\frac{\partial F}{\partial {}^{(3)}R}=\frac1{2\sqrt{{}^{(3)}R}}.
\end{equation}
The spatial response is
\begin{equation}
\boxed{
\tau_{ij}^{(\sqrt R)}
=
2\lambda_2
\left[
F_R{}^{(3)}R_{ij}
-\frac12h_{ij}F
+\frac1N
\left(h_{ij}\Delta-D_iD_j\right)(NF_R)
\right].
}
\end{equation}
Since this ADM potential also contains no shift dependence, its momentum-constraint response vanishes in the preferred ADM slicing:
\begin{equation}
\boxed{\delta q^{(\sqrt R)}_{\rm ADM}=0.}
\end{equation}
This is a preferred-foliation response statement, not a covariant matter-frame statement.

For Newtonian scalar perturbations
\begin{equation}
N=A(1+\Psi),
\qquad
h_{ij}=A^2(1-2\Phi)\gamma_{ij},
\end{equation}
the density response is
\begin{equation}
\boxed{
\delta\rho_{\sqrt R}
=
\frac{\rho_{\sqrt R}}3(3+\bar\Delta)\Phi.
}
\end{equation}
The pressure response is
\begin{equation}
\boxed{
\delta p_{\sqrt R}
=
\frac{\rho_{\sqrt R}}6
\left[
\frac23\bar\Delta\Psi
-2\bar\Delta\Phi
-\frac29\bar\Delta^2\Phi
-4\Phi
\right].
}
\end{equation}
The traceless response is
\begin{equation}
\boxed{
\Pi^i{}_{j,\sqrt R}
=
\frac{\rho_{\sqrt R}}6
\left(
\bar D^i\bar D_j
-\frac13\delta^i{}_j\bar\Delta
\right)
\left(
2\Phi-\Psi+\frac13\bar\Delta\Phi
\right).
}
\end{equation}
In scalar harmonics, \(\bar\Delta Y_{n\ell m}=-\kappa_nY_{n\ell m}\), \(\kappa_n=n(n+2)\), this becomes
\begin{align}
\boxed{
\delta\rho_{\sqrt R,n}
=
\frac{\rho_{\sqrt R}}3(3-\kappa_n)\Phi_n,
}\end{align}
\begin{align}
\boxed{
\delta p_{\sqrt R,n}
=
\frac{\rho_{\sqrt R}}6
\left[
\left(2\kappa_n-\frac29\kappa_n^2-4\right)\Phi_n
-\frac23\kappa_n\Psi_n
\right],
}\end{align}
\begin{align}
\boxed{
\Pi^i{}_{j,\sqrt R,n}
=
\frac{\rho_{\sqrt R}}6
\left(
\bar D^i\bar D_j
+\frac13\kappa_n\delta^i{}_j
\right)
\left[
\left(2-\frac{\kappa_n}{3}\right)\Phi_n-\Psi_n
\right].
}\end{align}

The term \(\bar\Delta^2\Phi\) in \(\delta p_{\sqrt R}\) shows that this is a higher-spatial-derivative metric-response kernel.  It is not a rest-frame sound speed and it is not transfer-equivalent to the canonical VDE scalar.  In an Einstein-equation implementation, the shear equation must be solved with the response included algebraically,
\begin{equation}
\text{standard shear}+\text{ADM response shear}(\Phi,\Psi)
\quad\Rightarrow\quad
\text{solve for }\Psi,
\end{equation}
before computing the induced \(\delta\rho\), \(\delta p\), and \(\Pi^i{}_j\).  Thus a numerical test must first choose the lift:
\begin{equation}
\boxed{
\text{canonical scalar VDE}
\quad\text{or}\quad
\text{ADM geometric-response VDE}.}
\end{equation}



The \(a^{-1}\) coefficient itself has an exact differential extractor. Since radiation, dust, total \(a^{-2}\), and vacuum have dilation eigenvalues \(-4,-3,-2,0\), define
\begin{equation}
\mathcal R_X[\mathcal H^2]
:=\Dil(\Dil+2)(\Dil+3)(\Dil+4)\mathcal H^2.
\end{equation}
Then
\begin{equation}
\boxed{
\widehat\Omega_X(a)
:=-\frac a6\mathcal R_X[\mathcal H^2]
=\Omega_X
}
\end{equation}
for the exact VDE profile. Curved \(\Lambda\)CDM satisfies \(\mathcal R_X=0\), whereas the five-power background including \(a^{-1}\) satisfies
\begin{equation}
\Dil(\Dil+1)(\Dil+2)(\Dil+3)(\Dil+4)\mathcal H^2=0.
\end{equation}

\begin{criterion}[Covariance-aware quotient estimator]
High-order numerical derivatives are not required for data analysis. For a reconstructed data vector
\begin{equation}
y=S\beta+x\Omega_X+\eta,
\qquad
\operatorname{Cov}(\eta)=C,
\end{equation}
where the columns of \(S\) span the nuisance directions \(1,a^{-2},a^{-3},a^{-4}\), \(x=a^{-1}\), and \(W=C^{-1}\), define the weighted residual projector
\begin{equation}
M_S:=I-S(S^TWS)^{-1}S^TW.
\end{equation}
The generalized-least-squares quotient coordinate is
\begin{equation}
\boxed{
\widehat\Omega_X
=\frac{x^TWM_Sy}{x^TWM_Sx},
\qquad
\operatorname{Var}(\widehat\Omega_X)
=\frac{1}{x^TWM_Sx}.
}
\end{equation}
It is invariant under every nonsingular reparameterization \(S\mapsto SR\) of the nuisance basis. This is the covariance-optimal experimental analogue of the analytic quotient projector.
\end{criterion}

\begin{theorem}[Covariance-optimal spectral subspace test]
Let
\begin{equation}
y=X_Sc+\eta,
\qquad
\eta\sim\mathcal N(0,C),
\qquad
W=C^{-1},
\end{equation}
where \((X_S)_{is}=a_i^s\) is the design matrix for a declared spectral set \(S\) and \(X_S\) has full column rank \(k\). Then
\begin{equation}
\boxed{
\widehat c=(X_S^TWX_S)^{-1}X_S^TWy,
\qquad
\operatorname{Cov}(\widehat c)=(X_S^TWX_S)^{-1}.
}
\end{equation}
With
\begin{equation}
M_{X_S}:=I-X_S(X_S^TWX_S)^{-1}X_S^TW,
\end{equation}
the generalized least-squares closure residual
\begin{equation}
\boxed{Q_S:=y^TWM_{X_S}y}
\end{equation}
satisfies
\begin{equation}
Q_S\sim\chi^2_{N-k}
\end{equation}
under the exact Gaussian model with known \(C\) and \(N\) data points. Both the fitted subspace and \(Q_S\) are invariant under every nonsingular change of spectral basis \(X_S\mapsto X_SR\).
\end{theorem}
\begin{proof}
Whiten the data with \(C^{-1/2}\):
\(\widetilde y=C^{-1/2}y\) and
\(\widetilde X_S=C^{-1/2}X_S\).
The GLS estimator is the ordinary least-squares estimator in the whitened
variables, giving the stated mean and covariance. The matrix
\begin{equation}
P_{\widetilde X_S}
=
\widetilde X_S(\widetilde X_S^T\widetilde X_S)^{-1}\widetilde X_S^T
\end{equation}
is an orthogonal projector of rank \(k\), and
\begin{equation}
Q_S
=
\|(I-P_{\widetilde X_S})\widetilde y\|^2.
\end{equation}
Under the null, the projected noise is a standard Gaussian vector in an
\((N-k)\)-dimensional subspace, hence
\(Q_S\sim\chi^2_{N-k}\). A nonsingular basis change
\(X_S\mapsto X_SR\) leaves the column space, and therefore the projector and
\(Q_S\), unchanged.
\end{proof}

\begin{criterion}[Positive spectral-cone test]
For the positive valuation hypothesis, replace the unconstrained spectral fit by
\begin{equation}
\widehat c_+
=\arg\min_{c_s\ge0}(y-X_Sc)^TW(y-X_Sc).
\end{equation}
This projects the data onto the convex cone generated by the allowed spectral columns. Boundary-active null distributions are generally chi-bar-square rather than ordinary chi-square and should be calibrated by the corresponding cone geometry or by Monte Carlo simulation.
\end{criterion}

\begin{remark}[Discrete Prony annihilator]
For logarithmically equally spaced samples \(x_n=\ln a_n=x_0+nh\), a finite dilation spectrum can also be tested without numerical differentiation. If \(Lf_n=f_{n+1}\), then a sequence generated by spectral set \(S\) satisfies the exact recurrence
\begin{equation}
\boxed{
\prod_{s\in S}(L-e^{sh})f_n=0.
}
\end{equation}
This is the discrete Prony counterpart of the continuous dilation annihilator.
\end{remark}

\begin{criterion}[VDE null hypothesis]
For the nested positive one-parameter VDE subclass,
\begin{equation}
H_0:\Omega_X=0,
\qquad
H_1:\Omega_X\ge0.
\end{equation}
A detection corresponds to the profiled maximum lying at \(\widehat\Omega_X>0\). The analytic extractor, GLS quotient coordinate, or an equivalent direct likelihood parameterization must be statistically consistent with a single constant \(\Omega_X\) over the tested redshift range.
\end{criterion}

\subsubsection{Equation of state and local CPL tangent}
For a general boundary valuation profile
\begin{equation}
F(a)=r_0+r_1a^{-1}+r_2a^{-2}+r_3a^{-3},
\end{equation}
with \(F(1)=1\), equivalently \(r_0+r_1+r_2+r_3=1\),
\begin{equation}
w(a)=-1-\frac13\frac{\dd\ln F}{\dd\ln a}
=-1+\frac{r_1a^{-1}+2r_2a^{-2}+3r_3a^{-3}}{3F(a)}.
\end{equation}
Let
\begin{equation}
N_0=r_1+2r_2+3r_3,
\qquad
M_0=r_1+4r_2+9r_3.
\end{equation}
The local CPL expansion \(w(a)=w_0+w_a(1-a)+O((1-a)^2)\) gives
\begin{equation}
w_0=-1+\frac{N_0}{3},
\qquad
w_a=\frac{M_0-N_0^2}{3}.
\end{equation}
After the homogeneous quotient, the minimal dark-energy sector is
\begin{equation}
\rho_{\Lambda X}(a)=\rho_\Lambda+\rho_Xa^{-1}.
\end{equation}
Define its present \(X\)-fraction by
\begin{equation}
r_X:=\frac{\rho_X}{\rho_\Lambda+\rho_X}.
\end{equation}

\begin{corollary}[Local CPL tangent curve of positive minimal VDE]
If \(\rho_\Lambda,\rho_X\ge0\), then
\begin{equation}
w_0=-1+\frac{r_X}{3},
\qquad
w_a=\frac{r_X(1-r_X)}{3},
\end{equation}
and hence
\begin{equation}
\boxed{
w_a=(1+w_0)(-2-3w_0),
\qquad
-1\le w_0\le-\frac23,
\qquad
0\le w_a\le\frac1{12}.
}
\end{equation}
This is the tangent relation at \(a=1\); it is not identical to a finite-redshift CPL fit.
\end{corollary}
\begin{proof}
At \(a=1\), the fraction of the \(a^{-1}\) component is \(r_X\), which gives \(w_0=-1+r_X/3\). Differentiating the fractional contribution gives \(w_a=r_X(1-r_X)/3\). Eliminating \(r_X=3(1+w_0)\) gives the displayed curve. Positivity implies \(0\le r_X\le1\), from which the bounds follow.
\end{proof}

Thus the positive minimal VDE tangent cannot occupy the local CPL quadrant \(w_0>-1,\ w_a<0\). Negative valuation coefficients can evade this sign statement, but positivity is then abandoned as an additional physical admissibility condition.

\subsubsection{Canonical covariant realizations of positive negative-exponent valuation branches}
The valuation and quotient theorems determine homogeneous branches but not their covariant sources. Canonical scalar fields provide explicit stable existence realizations for positive branches with negative dilation exponent. Consider
\begin{equation}
S_s=\int_{\M}\sqrt{-g}\left[
-\frac12\nabla_\mu\phi_s\nabla^\mu\phi_s-V_s(\phi_s)
\right]\dd^4x.
\end{equation}

\begin{theorem}[Canonical-scalar realization of a positive branch with \(s<0\)]
Suppose
\begin{equation}
\rho_s(a)=\rho_{s0}a^s>0,
\qquad
w_s=-1-\frac{s}{3},
\qquad
s<0,
\end{equation}
on an expanding FLRW interval. In natural units \(c=\hbar=1\), with \(M_{\rm Pl}^{-2}=8\pi G\), a canonical scalar realization along the background trajectory is obtained from
\begin{equation}
\boxed{
\left(\frac{\dd\phi_s}{\dd\tau}\right)^2
=-\frac{s}{3}\rho_s,
\qquad
V_s=\frac{6+s}{6}\rho_s,
}
\end{equation}
or equivalently
\begin{equation}
\boxed{
\frac{\dd\phi_s}{\dd\ln a}
=\pm M_{\rm Pl}
\frac{\sqrt{-s\,\Omega_s a^s}}{\mathcal H(a)}.
}
\end{equation}
Whenever \(\phi_s(a)\) is monotone, the pair \((\phi_s(a),V_s(a))\) determines \(V_s(\phi_s)\) uniquely along the cosmological trajectory. The realization has
\begin{equation}
c_{s,\rm rest}^2=1,
\qquad
\pi_{\mu\nu}=0,
\end{equation}
and positive canonical kinetic energy for \(\Omega_s>0\).
\end{theorem}
\begin{proof}
For a canonical scalar,
\(\rho=\tfrac12(\dd\phi/\dd\tau)^2+V\) and
\(p=\tfrac12(\dd\phi/\dd\tau)^2-V\). Hence
\((\dd\phi/\dd\tau)^2=(1+w)\rho\) and \(V=(1-w)\rho/2\). Substitution of \(w_s=-1-s/3\) gives the displayed relations. Using \(\rho_s=3M_{\rm Pl}^2H_0^2\Omega_s a^s\) and
\(\dd\phi/\dd\ln a=(\dd\phi/\dd\tau)/H\) gives the reconstruction formula. Unit rest-frame sound speed and zero scalar anisotropic stress are standard properties of a canonical minimally coupled scalar.
\end{proof}

For the two positive physical boundary-valuation directions this gives
\begin{align}
s=-1:\quad &w_X=-\frac23,
&\left(\frac{\dd\phi_X}{\dd\tau}\right)^2&=\frac13\rho_X,
&V_X&=\frac56\rho_X,
\\
s=-2:\quad &w_2=-\frac13,
&\left(\frac{\dd\phi_2}{\dd\tau}\right)^2&=\frac23\rho_2,
&V_2&=\frac23\rho_2,
\end{align}
and
\begin{equation}
\boxed{
\frac{\dd\phi_2}{\dd\ln a}
=\pm\sqrt2\,M_{\rm Pl}
\frac{\sqrt{\Omega_{2,\rm val}}\,a^{-1}}{\mathcal H(a)}.
}
\end{equation}
Thus the positive two-parameter physical boundary extension \((\Omega_X,\Omega_{2,\rm val})\) admits an explicit two-canonical-scalar existence realization. For \(s=-3\), a canonical scalar can reproduce the homogeneous \(w=0\) scaling but still has \(c_{s,\rm rest}^2=1\); it is therefore not a cold-dust completion. The endpoint dust branch is instead retained as the conserved pressureless-current realization described below.


\begin{remark}[Existence, not uniqueness]
The canonical scalars prove that the stable closure \(c_{s,\rm rest}^2=1\), \(\pi_{\mu\nu}=0\) has an explicit covariant realization for positive \(s=-1\) and \(s=-2\) branches. These realizations are not selected by the valuation geometry and are not unique. A likelihood analysis must state its source realization explicitly.
\end{remark}


\subsubsection{Scalar-mode mass invariants for canonical valuation branches}
\label{sec:scalar-mode-mass-invariants}

The canonical scalar realization of a positive valuation branch fixes more than the homogeneous equation of state.  Once the source class is declared to be a minimally coupled canonical scalar, the second derivative of the reconstructed potential along the cosmological trajectory is not a free function.  It is an invariant of the branch exponent and the expansion history.

\begin{theorem}[Canonical scalar mass invariant]
\label{thm:canonical-mass-invariant-general}
Let a positive branch
\begin{equation}
\rho_s(a)=\rho_{s0}a^s,
\qquad s<0,
\qquad
w_s=-1-\frac{s}{3},
\end{equation}
be realized by a minimally coupled canonical scalar field in units \(c=\hbar=1\).  Along the reconstructed background trajectory,
\begin{equation}
\left(\frac{\dd\phi_s}{\dd\tau}\right)^2
=-\frac{s}{3}\rho_s,
\qquad
V_s=\frac{6+s}{6}\rho_s,
\end{equation}
and
\begin{equation}
\boxed{
V_{s,\phi\phi}
=
-\frac{6+s}{4}H^2
\left[
 s+\Dil\ln\left(\frac{\mathcal H^2}{\mathcal H_*^2}\right)
\right].
}
\label{eq:general-mass-invariant}
\end{equation}
The arbitrary positive normalization \(\mathcal H_*^2\) drops out of the derivative.  Thus the linear scalar mass coefficient is fixed by the dilation flow of the background once the canonical source class and the exponent \(s\) are specified.
\end{theorem}

\begin{proof}
Write \(N=\ln a\).  From the canonical reconstruction,
\begin{equation}
\frac{\dd\phi_s}{\dd N}
=
\pm M_{\rm Pl}\frac{\sqrt{-s\,\Omega_s a^s}}{\mathcal H(a)},
\qquad
V_s(a)=\frac{6+s}{2}M_{\rm Pl}^2H_0^2\Omega_s a^s.
\end{equation}
Therefore \(\dd/\dd\phi_s=(\dd N/\dd\phi_s)\Dil\).  Differentiating \(V_s\) once gives a factor \(s\).  Differentiating once more gives one additional contribution from \(\mathcal H(a)\) through \(\dd N/\dd\phi_s\).  Combining the terms yields
\begin{equation}
V_{s,\phi\phi}
=-\frac{6+s}{4}H_0^2\mathcal H^2
\left[s+\Dil\ln \mathcal H^2\right]
=-\frac{6+s}{4}H^2
\left[s+\Dil\ln\left(\frac{\mathcal H^2}{\mathcal H_*^2}\right)\right],
\end{equation}
as claimed.
\end{proof}

For the two positive non-dust boundary directions one obtains the explicit specializations
\begin{align}
s=-1:\qquad
&\boxed{
 m_X^2
 =V_{X,\phi\phi}
 =\frac54H^2
 \left[1-\Dil\ln\left(\frac{\mathcal H^2}{\mathcal H_*^2}\right)\right],
}
\label{eq:X-mass-invariant}\\[1mm]
s=-2:\qquad
&\boxed{
 m_2^2
 =V_{2,\phi\phi}
 =H^2
 \left[2-\Dil\ln\left(\frac{\mathcal H^2}{\mathcal H_*^2}\right)\right].
}
\label{eq:sminus2-mass-invariant}
\end{align}
The first is the mass-flow invariant used by the nested canonical VDE sector.  The second is the corresponding invariant for the physical \(a^{-2}\) scalar realization.  Neither invariant is a stability theorem by itself; each is the exact coefficient that must enter the scalar perturbation operator of the declared canonical source.

\begin{corollary}[Closed-\(\Sthree\) scalar-mode frequency invariant]
\label{cor:closed-s3-scalar-frequency}
On the round closed spatial representative,
\begin{equation}
-\nabla^2Y_{n\ell m}
=
\frac{n(n+2)}{A^2}Y_{n\ell m}.
\end{equation}
The free scalar part of the canonical branch mode operator therefore contains
\begin{equation}
\boxed{
\omega_{s,n}^2
=
\frac{n(n+2)}{A^2}+V_{s,\phi\phi}.
}
\end{equation}
Equivalently, restoring the geometric factor \(c^2\) in the Laplacian term,
\begin{equation}
\boxed{
\Xi_{s,n}(a)
:=
\frac{\omega_{s,n}^2}{H^2}
=
\frac{-\Omega_K^{\rm geom}\,n(n+2)a^{-2}}{\mathcal H^2}
-\frac{6+s}{4}
\left[
 s+\Dil\ln\left(\frac{\mathcal H^2}{\mathcal H_*^2}\right)
\right].
}
\label{eq:Xi-s-n}
\end{equation}
\end{corollary}

\begin{remark}[First source-to-spectrum bridge, not full transfer closure]
Theorem~\ref{thm:canonical-mass-invariant-general} and Corollary~\ref{cor:closed-s3-scalar-frequency} connect three layers that were previously separate: the valuation exponent \(s\), the declared canonical scalar source, and the compact \(\Sthree\) scalar spectrum.  This is a genuine source-to-compact-spectrum bridge.  It is still not a complete gauge-invariant scalar/vector/tensor transfer hierarchy: metric constraints, visible-sector Boltzmann equations, initial covariance, lensing, and likelihood evaluation remain separate open steps.
\end{remark}

% Insert after the present ``Canonical covariant realizations'' subsection.
% This section uses the macros and theorem environments already defined in the manuscript.


\subsubsection{Back-propagated minimal canonical VDE completion and compact scalar-transfer kernel}
\label{sec:canonical-completion}

The valuation theorem fixes a background direction but, by itself, does not fix a
covariant source or its perturbations.  We now make a separate, deliberately
minimal completion choice.  The reverse construction is not an inference of a
microscopic origin from the auxiliary filling: it starts from the required
homogeneous VDE trajectory and reconstructs the unique canonical potential that
realizes that trajectory in a declared one-field class.  Generic inverse
quintessence reconstruction has precedents \cite{SahlenLiddleParkinson2005};
the model-specific content here is the valuation-fixed $a^{-1}$ trajectory and
its no-free-transfer closure.

\begin{definition}[Minimal canonical VDE completion]
Fix the standard cosmological parameters
\begin{equation}
\Theta=(H_0,\Omega_r,\Omega_m,\Omega_K^{\rm geom},\Omega_\Lambda,\Omega_X),
\qquad
\Omega_K^{\rm geom}<0,\qquad \Omega_X\ge0,
\end{equation}
with present-day normalization
\begin{equation}
\Omega_r+\Omega_m+\Omega_K^{\rm geom}+\Omega_\Lambda+\Omega_X=1.
\end{equation}
Let
\begin{equation}
F_{\rm VDE}(a):=\mathcal H^2(a)=\Omega_ra^{-4}+\Omega_ma^{-3}
+\Omega_K^{\rm geom}a^{-2}+\Omega_\Lambda+\Omega_Xa^{-1}
\end{equation}
on a connected expanding interval $I\subset(0,\infty)$ on which
$F_{\rm VDE}>0$.  The
minimal canonical VDE completion is the Einstein--matter action
\begin{equation}
S_{\rm cVDE}=\int_{\M}\sqrt{-g}\left[
\frac{M_{\rm Pl}^2}{2}R
-\frac12\nabla_\mu\phi\nabla^\mu\phi-U_\Theta(\phi)
\right]\dd^4x+S_{\rm std},
\end{equation}
where $S_{\rm std}$ contains the standard radiation, matter, and neutrino
sectors.  The scalar is the entire $\Lambda+X$ residual sector; it is not an
extra field in addition to an independently specified $\Lambda+X$ fluid.
\end{definition}

\begin{theorem}[Inverse canonical completion theorem]
\label{thm:inverse-canonical-completion}
Suppose $\Omega_X>0$.  Within the class of one minimally coupled canonical
scalar with positive constant kinetic normalization, the required residual
trajectory
\begin{align}
\rho_\phi(a)&=3M_{\rm Pl}^2H_0^2
\left(\Omega_\Lambda+\Omega_Xa^{-1}\right),\\
p_\phi(a)&=3M_{\rm Pl}^2H_0^2
\left(-\Omega_\Lambda-\frac23\Omega_Xa^{-1}\right)
\end{align}
determines the field and potential uniquely on $I$, up to
$\phi\mapsto\pm\phi+\text{constant}$.  They are
\begin{equation}
\boxed{
\phi(a)=\phi_*\pm M_{\rm Pl}\sqrt{\Omega_X}
\int_1^a\frac{\dd\tilde a}
{\tilde a^{3/2}\sqrt{F_{\rm VDE}(\tilde a)}} ,
\qquad
U_\Theta(\phi(a))=3M_{\rm Pl}^2H_0^2\Omega_\Lambda
+\frac52M_{\rm Pl}^2H_0^2\Omega_Xa^{-1}.}
\end{equation}
The potential is a single-valued function on the field image $\phi(I)$, since
$\phi(a)$ is monotone on every such interval.  The construction has no
additional cosmological parameter relative to curved $\Lambda$CDM plus
$\Omega_X$; $\phi_*$ and the sign are field-coordinate redundancies.
\end{theorem}

\begin{proof}
For a canonical scalar,
\begin{equation}
\left(\frac{\dd\phi}{\dd\tau}\right)^2=\rho_\phi+p_\phi,
\qquad
U=\frac{\rho_\phi-p_\phi}{2}.
\end{equation}
The prescribed trajectory therefore gives
\begin{equation}
\left(\frac{\dd\phi}{\dd\tau}\right)^2
=M_{\rm Pl}^2H_0^2\Omega_Xa^{-1},
\qquad
U=3M_{\rm Pl}^2H_0^2\Omega_\Lambda
+\frac52M_{\rm Pl}^2H_0^2\Omega_Xa^{-1}.
\end{equation}
Using $\dd a/\dd\tau=aH=aH_0\sqrt{F_{\rm VDE}}$ gives the displayed integral for
$\phi(a)$.
The scalar continuity identity is equivalent, for $\dd\phi/\dd\tau\neq0$, to the
Klein--Gordon equation
\begin{equation}
\frac{\dd^2\phi}{\dd\tau^2}+3H\frac{\dd\phi}{\dd\tau}+U_{,\phi}=0.
\end{equation}
Hence the reconstructed pair is an exact covariant solution.  Conversely,
$\rho_\phi$ and $p_\phi$ fix $(\dd\phi/\dd\tau)^2$ and $U$ pointwise; integrating the
nonzero field velocity leaves only the stated sign and shift freedom.
\end{proof}

\begin{corollary}[Flat $\Lambda+X$ potential]
If $\Omega_r=\Omega_m=\Omega_K^{\rm geom}=0$ and $\Omega_\Lambda>0$, then, with
$\rho_\Lambda=3M_{\rm Pl}^2H_0^2\Omega_\Lambda$,
\begin{equation}
\boxed{
U(\phi)=\rho_\Lambda+
\frac{5\rho_\Lambda}{12}
\left[\cosh\!\left(\frac{\phi-\phi_*}{M_{\rm Pl}}\right)-1\right].}
\end{equation}
The scalar alone then realizes exactly
$\rho_\phi=\rho_\Lambda+\rho_Xa^{-1}$ and
$p_\phi=-\rho_\Lambda-2\rho_Xa^{-1}/3$.
If instead $\Omega_\Lambda=0$, the flat pure-$X$ limit has the exponential
potential
\begin{equation}
U(\phi)=\frac52M_{\rm Pl}^2H_0^2\Omega_X
\exp\!\left[\mp\frac{\phi-\phi_*}{M_{\rm Pl}}\right].
\end{equation}
\end{corollary}

\begin{theorem}[Back-propagated scalar mass-flow invariant]
\label{thm:mass-flow}
For the completion of \cref{thm:inverse-canonical-completion}, the potential
derivatives on the VDE trajectory are
\begin{align}
U_{,\phi}(\phi(a))
&=\mp\frac52M_{\rm Pl}H_0^2
\sqrt{\Omega_Xa^{-1}F_{\rm VDE}(a)},\\
\boxed{\quad
m_\phi^2(a):=U_{,\phi\phi}(\phi(a))
=\frac54H_0^2\bigl[F_{\rm VDE}(a)-\Dil F_{\rm VDE}(a)\bigr].\quad}
\end{align}
Equivalently,
\begin{equation}
\boxed{
\frac{m_\phi^2}{H^2}
=\frac54\left[1-\Dil\ln F_{\rm VDE}\right]
=\frac{5}{4F_{\rm VDE}}\left(
\Omega_\Lambda+2\Omega_Xa^{-1}+3\Omega_K^{\rm geom}a^{-2}
+4\Omega_ma^{-3}+5\Omega_ra^{-4}
\right).}
\end{equation}
Thus the scalar mass is not a new fitted function.  A sufficient no-tachyon
condition on an interval is $F_{\rm VDE}-\Dil F_{\rm VDE}\ge0$ there.  The action has a positive
canonical kinetic normalization by construction; reality and strict positivity
of the reconstructed kinetic energy require $\Omega_X>0$.  Conversely, a
target with $\Omega_X<0$ cannot be realized in the declared one-canonical-field
class.
\end{theorem}

\begin{proof}
Differentiate the reconstructed potential with $\Dil=a\dd/\dd a$ and use
\begin{equation}
\Dil\phi=\pm M_{\rm Pl}\frac{\sqrt{\Omega_Xa^{-1}}}{\sqrt{F_{\rm VDE}}}.
\end{equation}
The first derivative follows immediately.  A second derivative gives
\begin{equation}
U_{,\phi\phi}
=-\frac52H_0^2\frac{\sqrt{F_{\rm VDE}}}{\sqrt{\Omega_Xa^{-1}}}
\Dil\!\left(\sqrt{\Omega_Xa^{-1}F_{\rm VDE}}\right)
=\frac54H_0^2(F_{\rm VDE}-\Dil F_{\rm VDE}).
\end{equation}
Expanding the last expression in the five background powers yields the
displayed spectral form.
\end{proof}

\begin{corollary}[Early-time decoupling]
\label{cor:early-decoupling}
If $\Omega_r>0$, the scalar fractional background density obeys
\begin{equation}
\frac{\rho_\phi}{3M_{\rm Pl}^2H^2}
=\frac{\Omega_\Lambda+\Omega_Xa^{-1}}{F_{\rm VDE}(a)}
=\frac{\Omega_X}{\Omega_r}a^3+O(a^4)
\qquad(a\to0).
\end{equation}
Hence the completed $\Lambda+X$ scalar does not introduce an early-dark-energy
plateau.  With the explicitly declared adiabatic primordial mode, its leading
early-time contribution to the standard radiation transfer system vanishes.
An independent scalar-isocurvature amplitude would be an additional model
choice and is excluded from the minimal completion.
\end{corollary}

\begin{remark}[Early-time effective-field boundary]
In radiation domination,
\begin{equation}
\phi-\phi_i=
\pm\frac23M_{\rm Pl}\sqrt{\frac{\Omega_X}{\Omega_r}}a^{3/2}+o(a^{3/2}),
\end{equation}
so the reconstructed potential has the inverse-power endpoint behavior
\begin{equation}
U(\phi)-\rho_\Lambda\sim
\frac52M_{\rm Pl}^2H_0^2\Omega_X
\left[
\frac32\sqrt{\frac{\Omega_r}{\Omega_X}}
\frac{|\phi-\phi_i|}{M_{\rm Pl}}
\right]^{-2/3}.
\end{equation}
The canonical completion is therefore a trajectory-complete late- and
early-universe effective model, not a claim of a UV-regular potential selected
by the valuation theorem.  In a $\Lambda$-dominated future,
$m_\phi^2/H^2\to5/4$; in $X$-, matter-, and radiation-dominated regimes the
corresponding limits are $5/2$, $5$, and $25/4$.
\end{remark}

\subsection{\texorpdfstring{Closed-\(\Sthree\) scalar transfer kernel}{Closed-S3 scalar transfer kernel}}

Let $Q_n$ be scalar harmonics on the round spatial representative,
\begin{equation}
-\nabla^2_{\Sthree}Q_n=\frac{n(n+2)}{A^2}Q_n.
\end{equation}
The $n=1$ sector is treated with the usual closed-FLRW gauge and constraint
qualifications.  In the natural-unit source sector ($c=1$), define the dimensionless bare scalar
transfer kernel
\begin{equation}
\boxed{
\mathscr W_n(a):=
(-\Omega_K^{\rm geom})a^{-2}n(n+2)
+\frac54\left[F_{\rm VDE}(a)-\Dil F_{\rm VDE}(a)\right].}
\label{eq:VDE-transfer-kernel}
\end{equation}
It follows that
\begin{equation}
\frac{n(n+2)/A^2+m_\phi^2}{H_0^2}=\mathscr W_n(a).
\end{equation}
For the Newtonian-gauge convention
\begin{equation}
\dd s^2=-(1+2\Psi)\dd\tau^2
+A^2(1-2\Phi)\dd\Omega_3^2,
\end{equation}
the scalar supplies the explicit linear stress sources
\begin{align}
\delta\rho_{\phi,n}
&=\frac{\dd\phi}{\dd\tau}\frac{\dd\delta\phi_n}{\dd\tau}
-\left(\frac{\dd\phi}{\dd\tau}\right)^2\Psi_n
+U_{,\phi}\delta\phi_n,\\
\delta p_{\phi,n}
&=\frac{\dd\phi}{\dd\tau}\frac{\dd\delta\phi_n}{\dd\tau}
-\left(\frac{\dd\phi}{\dd\tau}\right)^2\Psi_n
-U_{,\phi}\delta\phi_n,\\
\delta q_{\phi,n}&=-\frac{\dd\phi}{\dd\tau}\delta\phi_n,
\qquad
\pi^{(\phi)}_{\mu\nu}=0,
\end{align}
where $\delta T^0{}_{i}=\partial_i\delta q_\phi$ defines the momentum-potential
convention.  Thus the VDE source for the closed-FLRW Einstein constraints is
fully specified at linear scalar order.
In Newtonian gauge, the VDE scalar perturbation obeys
\begin{equation}
\frac{\dd^2\delta\phi_n}{\dd\tau^2}
+3H\frac{\dd\delta\phi_n}{\dd\tau}
+H_0^2\mathscr W_n\delta\phi_n
=\frac{\dd\phi}{\dd\tau}
\left(\frac{\dd\Psi_n}{\dd\tau}+3\frac{\dd\Phi_n}{\dd\tau}\right)
-2U_{,\phi}\Psi_n,
\label{eq:VDE-delta-phi}
\end{equation}
supplemented by the standard closed-FLRW Einstein constraints and the standard
Boltzmann hierarchies for the visible sectors.  Equation
\eqref{eq:VDE-delta-phi} fixes the VDE contribution to the scalar transfer
system.  It is not, by itself, a gauge-invariant total-mode frequency because
the metric constraints remain coupled.

\begin{theorem}[Canonical VDE source-coefficient closure]
\label{thm:no-free-transfer}
For $\Omega_X>0$, the minimal canonical completion has
\begin{equation}
\boxed{
c_{s,\phi,\rm rest}^2=1,
\qquad
\pi^{(\phi)}_{\mu\nu}=0,
\qquad
m_\phi^2/H^2=\frac54\left(1-\Dil\ln F_{\rm VDE}\right).}
\end{equation}
Consequently the VDE scalar introduces no independent sound-speed,
anisotropic-stress, or mass-function parameter.  With the standard adiabatic
primordial mode imposed as the explicit initial-condition choice, curved
$\Lambda$CDM plus the canonical VDE sector has precisely one new cosmological
parameter, $\Omega_X$.  At $\Omega_X=0$, the scalar reduces to an inactive
constant-potential sector and the model continuously returns to curved
$\Lambda$CDM.
\end{theorem}

\begin{proof}
The first two statements are exact properties of a minimally coupled canonical
scalar.  The mass relation is \cref{thm:mass-flow}.  The action, background
reconstruction, and \cref{eq:VDE-delta-phi} then determine every VDE scalar
coefficient once the standard cosmological parameters and the primordial
adiabatic mode are fixed.  The field shift and sign do not change the stress
tensor or linear observables.  The limit $\Omega_X\to0$ removes the kinetic
trajectory and leaves the constant potential $\rho_\Lambda$.
\end{proof}

\begin{remark}[Background analogues are not transfer-equivalent]
An $a^{-1}$ component realized as a defect network can share the same
homogeneous scaling while having a different covariant stress tensor and hence
different perturbation transfer functions.  In particular, a likelihood or
constraint obtained for a domain-wall analogue cannot be imported as a
likelihood for the canonical action above.  The latter is the declared
$c_{s,\phi,\rm rest}^2=1$ completion and must be run through its own Boltzmann
transfer system \cite{AnHanZhang2025}.
\end{remark}

\begin{criterion}[Joint expansion--growth falsification test]
The minimal canonical VDE model is not tested by $H(a)$ alone.  It requires
simultaneous consistency of the constant quotient coordinate
\begin{equation}
\widehat\Omega_X(a)=-\frac a6\Dil(\Dil+2)(\Dil+3)(\Dil+4)\mathcal H^2(a)
=\Omega_X,
\end{equation}
the source-level identities in \cref{thm:no-free-transfer}, and the closed
harmonic transfer equation \cref{eq:VDE-delta-phi}.  In the late-time,
subhorizon pressureless-matter regime, where radiation and relativistic-species
corrections are negligible or are included separately in the standard hierarchy,
the leading GR growth null is
\begin{equation}
\boxed{
\mathcal G_{\rm VDE}[D_m,F_{\rm VDE}]
:=D_m''+\left(2+\frac12\Dil\ln F_{\rm VDE}\right)D_m'
-\frac{3\Omega_ma^{-3}}{2F_{\rm VDE}}D_m
=\mathcal O\!\left[(H/k_n)^2\right],}
\end{equation}
where primes denote $\dd/\dd\ln a$, $D_m$ is the growing-mode density contrast
of total pressureless matter, and $k_n^2=n(n+2)/A^2$ is the physical harmonic
wavenumber.  A statistically significant failure of this joint
background--growth--lensing system falsifies the minimal canonical VDE
completion within the declared GR and standard-sector assumptions, even if a
background-only $a^{-1}$ fit is acceptable.
\end{criterion}

\begin{remark}[What this completion does not prove]
The reverse construction upgrades the $a^{-1}$ branch to a fully specified
one-field effective cosmology, but it does not derive $U_\Theta$ from
$\Efour^{\rm val}$, determine $\Omega_X$ from topology, or establish a
compact-valuation origin.  The potential is a family of effective actions
labeled by the matched cosmological parameters $\Theta$.  A Boltzmann-code
likelihood evaluation, for example with a CLASS-type curved-FLRW implementation
\cite{BlasLesgourguesTram2011,LesgourguesTram2014}, is an empirical calculation
required to decide whether the completed one-parameter model is preferred or
excluded.
\end{remark}


\begin{remark}[Boundary null and likelihood calibration]
For \(\Omega_X\ge0\), the null point \(\Omega_X=0\) lies on the boundary of parameter space. A likelihood-ratio significance therefore should not automatically be calibrated with the ordinary interior-point Wilks law; boundary-aware asymptotics or simulation-based calibration are appropriate \cite{Chernoff1954,SelfLiang1987}.
\end{remark}

The local curvature-measure boundary restriction above is an independent geometric filter. Inside the explicitly stated minimal endpoint class, the following theorem gives a separate conditional endpoint classification. It is not a theorem from topology alone; it is conditionally unique only inside that class.

\begin{definition}[Endpoint locality convention]
In the minimal endpoint theorem, locality is germ-local with respect to the
metric dependence: the local metric density on an open set may depend only
on the local metric germ and is independent of global charge labels of the
connected component. Global topological data are permitted only through a
separate metric-independent additive functional \(f(Q)\). Thus a coupling
such as \(\sigma(Q)\operatorname{Vol}(\Sigma)\), in which a global charge
modulates the coefficient of a local metric density, is outside the declared
minimal endpoint class.
\end{definition}

\begin{theorem}[Finite-charge oriented minimal endpoint classification]
Let \(\Sigma\simeq S^3\) be connected and oriented. Assume:
\begin{enumerate}
\item locality and \({\rm Diff}^{+}(\Sigma)\)-naturality;
\item no local propagating endpoint degrees of freedom;
\item no intrinsic/extrinsic curvature terms, higher derivatives, embedding-position data, background connections, nonlocal spectral functionals, or additional geometric/dynamical scales in the functional dependence; the overall dimensional coefficients remain free;
\item additivity on finite disjoint unions;
\item the global internal topological data form a prescribed finite-rank additive charge vector
\begin{equation}
Q=(Q_1,\ldots,Q_r)\in\mathbb Z^r,
\end{equation}
are independent of the metric scale, and enter only through the separate additive metric-independent functional \(f(Q)\).
\end{enumerate}
The additivity assumption is understood on the extended class of finite disjoint unions of oriented endpoint components and is then restricted to the connected component \(\Sigma\simeq S^3\). Component-count constants are excluded as nonlocal topological inputs outside the declared local endpoint class. Then every endpoint energy in this class has the form
\begin{equation}
\boxed{
E_{\rm end}
=
\sigma\,\operatorname{Vol}(\Sigma)
+\sum_{A=1}^{r}\mu_AQ_A.
}
\end{equation}
\end{theorem}
\begin{proof}
By the germ-locality convention, the coefficient of the local metric density cannot depend on global charge labels. After the stated exclusions, the only natural zeroth-order metric density is therefore a universal constant multiple of the Riemannian volume density. The remaining internal dependence is the separate metric-independent function \(f:\mathbb Z^r\to\mathbb R\). Additivity makes \(f\) a group homomorphism, hence
\begin{equation}
f(Q)=\sum_AQ_Af(e_A)=\sum_A\mu_AQ_A.
\end{equation}
\end{proof}

\begin{corollary}[Single-degree endpoint sector]
If the declared internal topological datum is specifically the oriented degree
\begin{equation}
\QU=\deg\Phi,
\qquad
\Phi:S^3\to S^3,
\end{equation}
with normalized target volume form \(\omega_3\), then
\begin{equation}
\boxed{
E_{\rm end}=\sigma V_3+\mu\QU,
\qquad
\epsD=\sigma+\frac{\mu\QU}{2\pi^2A^3},
\qquad
\pD=-\sigma.
}
\end{equation}
Equivalently, for \(A=A_0a\),
\begin{equation}
\epsD(a)=\varepsilon_q+\varepsilon_w a^{-3},
\qquad
\varepsilon_q=\sigma,
\qquad
\varepsilon_w=\frac{\mu\QU}{2\pi^2A_0^3}.
\end{equation}
The degree-map choice is therefore an explicit model assumption, not a consequence of the phrase ``topological charge'' by itself.
\end{corollary}

\begin{remark}[Metric volume term versus topological charge]
The term \(\sigma V_3\) is not a topological invariant: it arises from the unique natural zeroth-order metric density. On connected \(\Sigma\), its coefficient is a constant, which is the only limited sense in which \(H^0(\Sigma)\) is relevant. By contrast, the charge \(\QU\) is genuinely topological and metric-scale independent.
\end{remark}

\subsection{Pressure Kernel and the Dust Branch}
The work-pressure map has a nontrivial kernel. This kernel is exactly the structural origin of the dust scaling in the minimal endpoint sector.

\begin{theorem}[Kernel of the work-pressure map]
Let \(E(A)\in C^1(J)\) on an interval \(J\subset(0,\infty)\), and let
\begin{equation}
\varepsilon(A)=\frac{E(A)}{V_3(A)},
\qquad
p(A)=-\frac{\dd E}{\dd V_3}.
\end{equation}
Then
\begin{equation}
p(A)=0
\end{equation}
on \(J\) if and only if
\begin{equation}
E(A)=E_0
\end{equation}
is constant on \(J\). In that case
\begin{equation}
\varepsilon(A)=\frac{E_0}{2\pi^2A^3}
=\varepsilon_w a^{-3}.
\end{equation}
Thus the homogeneous pressureless branch is precisely the image, under division by \(V_3\), of the constant-energy kernel of the work-pressure map.
\end{theorem}
\begin{proof}
Since
\begin{equation}
p=-\frac{\dd E}{\dd V_3}
=-\frac{\dd E/\dd A}{\dd V_3/\dd A},
\end{equation}
and since
\begin{equation}
\frac{\dd V_3}{\dd A}=6\pi^2A^2\neq0
\end{equation}
for \(A>0\), the condition \(p=0\) is equivalent to
\begin{equation}
\frac{\dd E}{\dd A}=0.
\end{equation}
Hence \(E(A)=E_0\). Division by \(V_3=2\pi^2A^3\) gives the dust scaling.
\end{proof}

\begin{corollary}[Endpoint normal form from volume plus pressure kernel]
The minimal endpoint energy
\begin{equation}
E_{\rm end}(A,\QU)=\sigma V_3(A)+\mu\QU
\end{equation}
is the direct sum of two distinguished work-map pieces:
\begin{equation}
\sigma V_3(A)
\quad\text{and}\quad
\mu\QU.
\end{equation}
The first is the volume branch and gives
\begin{equation}
\varepsilon_q=\sigma,
\qquad
p_q=-\sigma.
\end{equation}
At fixed topological charge, the second lies in the pressure kernel and gives
\begin{equation}
\varepsilon_w(a)=\frac{\mu\QU}{2\pi^2A_0^3}a^{-3},
\qquad
p_w=0.
\end{equation}
Therefore the endpoint vacuum-dust form is the normal form of the endpoint energy under the decomposition
\begin{equation}
E_{\rm end}
\in
\operatorname{span}\{V_3\}
\oplus
\ker\!\left(E\mapsto-\frac{\dd E}{\dd V_3}\right),
\end{equation}
with the second summand restricted by the endpoint classification to the additive topological charge term.
\end{corollary}

\begin{remark}[Scope of the endpoint classification]
This is not a theorem from topology alone. It is a conditional classification inside the declared minimal endpoint class: local \({\rm Diff}^{+}(\Sigma)\)-natural zeroth-order metric dependence, additivity, no local endpoint modes, no curvature or embedding duplication, and a prescribed finite-rank additive charge lattice. Enlarging the functional class by curvature, gradient, background-connection, nonlocal spectral, embedding-dependent, or additional dynamical structures invalidates the uniqueness statement. The single-degree sector \(\QU=\deg\Phi\) is a declared specialization of the finite-charge theorem, not its only possible topological realization.
\end{remark}

\begin{remark}[Relation to valuation coefficients]
The valuation endpoint notation
\begin{equation}
\alpha_0\int_{\Sthree_A}\dd\Sigma+
\alpha_3\int_{\Sthree_A}\det(K^a{}_b)\,\dd\Sigma
\end{equation}
contains the same two homothetic powers. In the normalized topological-charge convention used in the minimal endpoint classification, compare the two round \(A^0\) background amplitudes only after restricting to a fixed superselection sector
\begin{equation}
\QU=Q_*.
\end{equation}
Then
\begin{equation}
2\pi^2\alpha_0A^3+2\pi^2\alpha_3^{\rm(match)}
=2\pi^2\sigma A^3+\mu Q_*,
\end{equation}
so that
\begin{equation}
\alpha_0=\sigma,
\qquad
\boxed{
\alpha_3^{\rm(match)}=\frac{\mu Q_*}{2\pi^2}.
}
\end{equation}
This is an equality of the round \(A^0\) background amplitudes in that fixed sector, not an equality of the Gauss-curvature valuation functional with the internal matter-charge functional.
For the Gauss-map endpoint of the convex round filling, define \(\QG:=\deg(\nu)=1\). The valuation constant slot and the internal endpoint charge share the homothetic power \(A^0\), but their physical sources are distinct. The Gauss-map realization carries \(\QG=1\), whereas the internal matter field carries \(\QU\in\mathbb Z\). It is not an admission of an independent extrinsic-curvature action into the minimal endpoint class. More generally, an internal map \(\Phi:\Sthree\to\Sthree\) supplies the integer charge. The topological charge selects the dust background scaling; it is not by itself a cold-dark-matter perturbation theory.
\end{remark}

\begin{remark}[Metric volume term and the constant homothetic slot]
The volume contribution \(\sigma V_3\) comes from the unique natural zeroth-order metric density in the declared endpoint class. The topological term is separately supplied by the chosen additive charge lattice; in the single-degree specialization it is \(\mu\QU\). The two contributions therefore generate the \(A^3\) and \(A^0\) endpoint powers for different structural reasons. Neither \(H^0(S^3)\) nor the Gauss--Kronecker curvature term is identified with the internal matter charge functional.
\end{remark}

\begin{remark}[Gauss--Kronecker term and Gauss-map degree]
The endpoint invariant
\begin{equation}
\int_{\Sthree_A}\det(K^a{}_b)\,\dd\Sigma
\end{equation}
is not an Euler characteristic of \(\Sthree\). Since \(\Sthree\) is odd-dimensional,
\begin{equation}
\chi(\Sthree)=0.
\end{equation}
For an oriented closed hypersurface in the auxiliary Euclidean space \(\mathbb{R}^4\), the Gauss--Kronecker integral is instead the Jacobian degree of the Gauss map
\begin{equation}
\nu:\Sthree_A\longrightarrow \Sthree_{\rm unit}.
\end{equation}
With the chosen outward orientation,
\begin{equation}
\int_{\Sthree_A}\det(K^a{}_b)\,\dd\Sigma
=\deg(\nu)\,\vol(\Sthree_{\rm unit}).
\end{equation}
For the convex round filling, \(\QG:=\deg(\nu)=1\), and therefore
\begin{equation}
\int_{\Sthree_A}\det(K^a{}_b)\,\dd\Sigma=2\pi^2.
\end{equation}
Thus the top endpoint term is a Gauss-map degree invariant in the round convex case, not a Gauss--Bonnet--Chern Euler term. In the intrinsic-volume language, this top boundary curvature integral corresponds, up to the normalization absorbed into \(\alpha_3\), to the zeroth intrinsic volume of the convex body. The sign of the degree depends on the orientation convention for the normal; throughout this paper the outward normal of the convex filling is used.
\end{remark}

\begin{nogotheorem}[Scale-covariance no-go for topological amplitudes]
Within a construction whose inputs are purely topological or dimensionless and contain no metric scale, dimensional coupling, cutoff, or matching scale, no nonzero quantity with nonzero length dimension can be determined naturally from topology alone. In particular, pure topology cannot determine nonzero values of \(\sigma\), \(\mu\), \(H_0\), or any other dimensional dark-sector amplitude.
\end{nogotheorem}
\begin{proof}
Rescale all metric lengths by \(L\mapsto\lambda L\) while keeping the topological data and integer charges fixed. A quantity \(X\) constructed only from those topological data is unchanged. If instead \(X\) carries length dimension \(L^{-k}\), \(k\neq0\), dimensional covariance requires \(X\mapsto\lambda^{-k}X\). Compatibility for arbitrary \(\lambda>0\) implies \(X=0\). Therefore a nonzero dimensional amplitude requires an independent physical scale or matching coefficient.
\end{proof}

\begin{remark}[Amplitudes as matching data]
The coefficients \(\sigma\) and \(\mu\) are therefore not unresolved topological predictions inside the stated axioms. They must be treated as matching/coupling coefficients (Wilson coefficients only after embedding in an explicit EFT) supplied by microscopic physics, renormalized couplings, an EFT scale, or observations. The no-go theorem closes the topology-only amplitude-selection question negatively.
\end{remark}

\begin{corollary}[No amplitude from endpoint linearity]
The same no-go applies to every linear combination of the endpoint invariants. The minimal endpoint classification fixes the two powers \(A^3\) and \(A^0\), but the coefficients multiplying them remain free dimensional constants until an independent physical or observational input is supplied.
\end{corollary}

\begin{corollary}[Exact scope of predictivity]
The full valuation theorem gives the five-slot homothetic spectrum
\begin{equation}
\epsD(A)=\frac{\gamma A}{4}+\alpha_0+\frac{3\alpha_1}{A}+\frac{3\alpha_2}{A^2}+\frac{\alpha_3}{A^3}.
\end{equation}
The local curvature-measure boundary restriction \(\gamma=0\) removes the bulk branch relative to the explicitly stated local curvature-measure boundary restriction. Inside the stated minimal endpoint class one obtains the functional form
\begin{equation}
\epsD(a)=\varepsilon_q+\varepsilon_w a^{-3},
\qquad
\pD=-\varepsilon_q,
\end{equation}
but not the numerical values of \(\varepsilon_q\) and \(\varepsilon_w\). Thus the minimal endpoint class reduces the arbitrary background function \(\epsD(A)\) to two constants after the stated assumptions. It does not determine those constants by topology alone.
\end{corollary}

\section{Independent Covariant Source Realizations and Conditional Global Charge Matching}
The minimal endpoint classification selects the vacuum--dust background form, but a local stress tensor still requires a covariant matter model. The actions in this section are independent existence realizations of selected homogeneous branches. They are not derived by covariant variation of \(\Efour^{\rm val}\), and the valuation hypothesis does not uniquely select them. We keep the local dust current independent of the global topological label; this separation avoids identifying a global degree density with a local particle current.

Let
\begin{equation}
S_{\rm eff}=S_{\rm grav}+S_{\rm vis}+S_q+S_w,
\qquad
S_q=-\int_{\M}\sqrt{-g}\,\varepsilon_q\,\dd^4x,
\end{equation}
so that
\begin{equation}
T^q_{\mu\nu}=-\varepsilon_q g_{\mu\nu},
\qquad
\nabla_\mu\varepsilon_q=0.
\end{equation}
For dust we use a global label-space formulation of the standard pressureless-current construction \cite{Brown1993Fluid,BrownKuchar1995}.

\begin{definition}[Global dust-label space and conserved current]
Let \(\mathcal S\) be an oriented three-dimensional dust-label manifold with
a nowhere-vanishing volume form \(\omega_{\mathcal S}\), and let
\(Z:\M\to\mathcal S\) be the comoving-label map. Define
\begin{equation}
N_Z:=Z^*\omega_{\mathcal S},
\qquad
\dd N_Z=0.
\end{equation}
The associated densitized current is
\begin{equation}
\mathcal J^\mu
:=
\frac1{3!}\epsilon^{\mu\nu\rho\sigma}(N_Z)_{\nu\rho\sigma},
\qquad
\partial_\mu\mathcal J^\mu=0.
\end{equation}
On the timelike-current sector define
\begin{equation}
n
:=
\frac{\sqrt{-g_{\mu\nu}\mathcal J^\mu\mathcal J^\nu}}{\sqrt{-g}},
\qquad
u^\mu
:=
\frac{\mathcal J^\mu}{\sqrt{-g}\,n}.
\end{equation}
Then
\begin{equation}
\nabla_\mu(n u^\mu)=0.
\end{equation}
In a local oriented chart \(Z^A\), this reproduces the usual three-label
formula.
\end{definition}

\begin{nogotheorem}[No single global Cartesian dust chart on compact \(S^3\)]
There do not exist three globally defined real-valued functions \(Z^A:S^3\to\mathbb R\) such that
\begin{equation}
\dd Z^1\wedge\dd Z^2\wedge\dd Z^3
\end{equation}
is everywhere nonzero.
\end{nogotheorem}
\begin{proof}
Such a map \(Z:S^3\to\mathbb R^3\) would be a local diffeomorphism and therefore an open map. Its image would also be compact because \(S^3\) is compact. Hence the image would be a nonempty subset of connected \(\mathbb R^3\) that is both open and compact, which is impossible.
\end{proof}

\begin{theorem}[Global pressureless-current realization]
With the global densitized current above, the action
\begin{equation}
S_w=-m\int_{\M}\sqrt{-g}\,n\,\dd^4x
\end{equation}
has perfect-dust stress tensor
\begin{equation}
T_{\mu\nu}^{w}=\rho_wu_\mu u_\nu,
\qquad
\rho_w=mn.
\end{equation}
On shell,
\begin{equation}
\nabla_\mu(n u^\mu)=0,
\qquad
u^\mu\nabla_\mu u^\nu=0,
\qquad
p_w=0,
\end{equation}
and the minimal perturbative closure is
\begin{equation}
c_{s,\rm rest}^2=0,
\qquad
\pi_{\mu\nu}=0.
\end{equation}
For a homogeneous FLRW configuration,
\begin{equation}
n(a)=n_0a^{-3},
\qquad
\varepsilon_w=mn_0.
\end{equation}
The normalization of \(\omega_{\mathcal S}\) is matching data of the local covariant realization and is not fixed by the endpoint topology theorem.
\end{theorem}
\begin{proof}
The action has the barotropic pullback-fluid form
\begin{equation}
S_w=-\int_{\M}\sqrt{-g}\,\rho(n)\,\dd^4x,
\qquad
\rho(n)=mn.
\end{equation}
Metric variation at fixed pullback three-form gives the perfect-fluid stress tensor
\begin{equation}
T_{\mu\nu}
=(\rho+p)u_\mu u_\nu+pg_{\mu\nu},
\qquad
p=n\frac{\dd\rho}{\dd n}-\rho.
\end{equation}
For \(\rho(n)=mn\), one has \(p=0\), hence
\(T_{\mu\nu}=mn\,u_\mu u_\nu\). The pullback identity gives
\(\nabla_\mu(n u^\mu)=0\). On the label equations of motion,
diffeomorphism invariance gives \(\nabla_\mu T^{\mu\nu}=0\); combining this
with particle-number conservation yields
\(u^\mu\nabla_\mu u^\nu=0\). Since \(p\equiv0\),
\begin{equation}
c_{s,\rm rest}^2=0,
\qquad
\pi_{\mu\nu}=0.
\end{equation}
For a homogeneous FLRW flow, number conservation gives \(n\propto a^{-3}\).
\end{proof}

\begin{proposition}[Conditional global charge--local dust matching]
Let \(\QU\in\mathbb Z\) be an oriented global topological charge. After the additive endpoint classification theorem, restrict to a fixed positive-orientation superselection sector \(\QU\ge0\), and let \(\mu\) be a nonnegative energy-per-charge matching coefficient (equivalently mass per charge in the natural-unit convention \(c=1\)). Let \(\mathfrak n^\mu\) be the future-directed unit normal to the spacelike slice \(\Sigma_t\),
\begin{equation}
\mathfrak n_\mu\mathfrak n^\mu=-1.
\end{equation}
Define the conserved dust rest-energy/rest-mass charge
\begin{equation}
\mathcal E_{\rm rest}(t)
:=\int_{\Sigma_t}\rho_w(-u^\mu\mathfrak n_\mu)\,\dd\Sigma .
\end{equation}
If the matching condition
\begin{equation}
\boxed{\mathcal E_{\rm rest}(t_*)=\mu\QU}
\end{equation}
is imposed on one compact Cauchy slice \(\Sigma_{t_*}\), then current conservation implies
\begin{equation}
\boxed{\mathcal E_{\rm rest}(t)=\mu\QU}
\end{equation}
on every homologous compact Cauchy slice. On a comoving FLRW slice, where \(u^\mu=\mathfrak n^\mu\), this gives
\begin{equation}
\mathcal E_{\rm rest}=2\pi^2A_0^3\varepsilon_w,
\qquad
\boxed{\varepsilon_w=\frac{\mu\QU}{2\pi^2A_0^3}.}
\end{equation}
This is a matching condition propagated by the local conserved current, not a derivation of \(\mathcal E_{\rm rest}\overset{\rm match}{=}\mu\QU\) from topology alone.
\end{proposition}
\begin{proof}
For constant \(m\),
\begin{equation}
\nabla_\mu(\rho_wu^\mu)=m\nabla_\mu(nu^\mu)=0.
\end{equation}
Integrating this conserved current over the spacetime slab between two compact homologous Cauchy slices and applying Stokes' theorem gives equality of \(\int_{\Sigma_t}\rho_w(-u\cdot\mathfrak n)\dd\Sigma\) on the two slices. There is no external spatial boundary for \(\Sigma_t\simeq S^3\). Therefore a matching imposed on one slice is propagated to all homologous slices, but the value \(\mu\QU\) itself is external matching data.
\end{proof}

\begin{nogotheorem}[Metric-independent topological terms carry no stress tensor]
If a topological action \(S_{\rm top}\) is genuinely independent of the spacetime metric, then
\begin{equation}
\boxed{
T^{\rm top}_{\mu\nu}
=-\frac2{\sqrt{-g}}\frac{\delta S_{\rm top}}{\delta g^{\mu\nu}}
=0.
}
\end{equation}
Consequently an integer degree by itself may label or quantize a global sector, but it cannot by itself generate a pressureless stress tensor.
\end{nogotheorem}
\begin{proof}
Metric independence means \(\delta S_{\rm top}/\delta g^{\mu\nu}=0\) identically. The displayed result then follows directly from the variational definition of the stress tensor.
\end{proof}

\section{Topological Degree: Orientation and a Homogeneity Obstruction}
Let
\begin{equation}
U:\Sthree\longrightarrow SU(2)\simeq\Sthree
\end{equation}
be smooth. With normalized closed target volume form \(\omega_{SU(2)}\),
\begin{equation}
\QU[U]=\int_{\Sthree}U^*\omega_{SU(2)}\in\mathbb Z.
\end{equation}
The charge is oriented: an orientation reversal sends \(\QU\mapsto-\QU\). A positive particle-number interpretation may be imposed only after restricting to a fixed-orientation sector \(\QU\ge0\).

\begin{nogotheorem}[No arbitrary homogeneous degree from one smooth map]
Suppose a single smooth map \(U:\Sthree\to\Sthree\) is required to produce a strictly homogeneous nonzero topological density,
\begin{equation}
U^*\omega_{\Sthree}
=\QU\,\frac{\dd V_3}{V_3},
\qquad
\QU\neq0,
\end{equation}
where both sides are normalized so that integration gives \(\QU\). Then
\begin{equation}
\boxed{|\QU|=1.}
\end{equation}
Thus one smooth map cannot generate a strictly homogeneous degree density with arbitrary \(|\QU|>1\) on \(\Sthree\).
\end{nogotheorem}
\begin{proof}
The assumed pullback has everywhere nonzero Jacobian, so \(U\) is a local diffeomorphism. A local diffeomorphism from compact \(\Sthree\) to connected \(\Sthree\) is a covering map. Since the target \(\Sthree\) is simply connected, every connected covering is one-sheeted. Hence \(U\) is a diffeomorphism and its degree is \(\pm1\).
\end{proof}

\begin{remark}[Gauss degree and matter degree are distinct]
The convex-filling Gauss map carries
\begin{equation}
\QG:=\deg(\nu)=1,
\end{equation}
whereas an internal field carries \(\QU:=\deg(U)\). The equality of their homothetic \(A^0\) scaling does not identify them as the same physical charge. The valuation constant slot has a Gauss-map realization; the minimal endpoint matter term uses an independent oriented internal charge.
\end{remark}

\begin{remark}[Why the global matching is preferred]
The homogeneity obstruction shows why the local dust current should not be identified pointwise with the pullback density of a single arbitrary-degree map. The logically cleaner construction is
\begin{equation}
\text{global topology: }\QU
\quad\oplus\quad
\text{local dynamics: }\mathcal J^\mu,
\end{equation}
linked only by \(\mathcal E_{\rm rest}\overset{\rm match}{=}\mu\QU\). This keeps topological quantization and cold-dust perturbations conceptually separate.
\end{remark}

\section{Minisuperspace Contribution to the Total Action}
The full background action must have the form
\begin{equation}
S_{\rm tot}=S_{\rm grav}[A,N]+S_{\rm vis}[A,N,\psi]+S_{\rm fill}[A,N,\xi].
\end{equation}
The filling contribution is
\begin{equation}
S_{\rm fill}=-\int \dd t\,N(t)\Efour(A(t),\xi(t)).
\end{equation}
These reduced variational formulae hold in the sector where \(\Efour=\Efour(A,\xi)\) has no independent dependence on \(N\), \(\dot A\), \(\dot\xi\), curvature scalars, or extrinsic-curvature boundary data. Such dependences belong to the covariant completion problem.
Its variations are
\begin{equation}
\frac{\delta S_{\rm fill}}{\delta N(t)}=-\Efour,
\qquad
\frac{\delta S_{\rm fill}}{\delta A(t)}=-N\left.\frac{\partial\Efour}{\partial A}\right|_{\xi}.
\end{equation}
Hence
\begin{equation}
\epsD=-\frac{1}{V_3}\frac{\delta S_{\rm fill}}{\delta N},
\qquad
\pD=\frac{A}{3NV_3}\frac{\delta S_{\rm fill}}{\delta A}.
\end{equation}

\begin{remark}
The functional \(S_{\rm fill}\) determines the contribution of the filling sector to the full equations. It cannot be varied as a self-contained universe: equations of motion follow from \(S_{\rm tot}\), not from imposing \(\delta S_{\rm fill}=0\) separately. What has been constructed is a consistent minisuperspace contribution, not a complete local four-dimensional action. The distinction between a reduced action and a completed gravitational variational principle is standard in canonical and boundary-action treatments \cite{DeWitt1967,GibbonsHawking1977}.
\end{remark}

\section{Reconstructed, Source-Realized, and Physical Tensors}
For a specified metric and visible sector, define the reconstructed target residual
\begin{equation}
T^{\rm rec}_{\mu\nu}:=\frac{c^4}{8\pi G}G_{\mu\nu}-T^{\rm vis}_{\mu\nu}.
\end{equation}
This is a covariant reconstruction, not a microscopic derivation.

Let \(T^{\rm src}_{\mu\nu}\) denote the stress tensor of an independently supplied covariant source realization, such as the canonical-scalar or pressureless-current models described above. On the FLRW background, its mixed components must reduce to
\begin{equation}
(T^{\rm src})^{\mu}{}_{\nu}=\diag(-\epsD,\pD,\pD,\pD).
\end{equation}
The substantive matching condition for a source-realized residual is
\begin{equation}
T^{\rm src}_{\mu\nu}=T^{\rm rec}_{\mu\nu}.
\end{equation}
Only after this has been satisfied are both objects denoted by the physical dark tensor \(T^D_{\mu\nu}\). A stronger equality with a tensor obtained by covariantly varying a four-dimensional action derived from \(\Efour^{\rm val}\) is not constructed here; that is the open valuation-to-source bridge.

\begin{theorem}[Covariant algebraic perfect-fluid decomposition and differential vacuum--dust test]
Let \(T^\mu{}_{\nu}\) be a type-I stress tensor with one timelike eigenvector \(u^\mu\) and a triple-degenerate spatial eigenvalue, so algebraically it is an isotropic perfect-fluid tensor. Define
\begin{equation}
S^\mu{}_{\nu}
:=T^\mu{}_{\nu}-\frac14T\,\delta^\mu{}_{\nu}.
\end{equation}
Then the coordinate-invariant syzygies
\begin{equation}
\boxed{
3[\operatorname{tr}(S^3)]^2=[\operatorname{tr}(S^2)]^3,
\qquad
12\operatorname{tr}(S^4)=7[\operatorname{tr}(S^2)]^2
}
\end{equation}
hold for every isotropic perfect fluid, not only for vacuum plus dust. If \(\rho+p\ge0\), define the algebraic quantities
\begin{equation}
\boxed{
\rho_w:=\rho+p
=\sqrt{\frac43\operatorname{tr}(S^2)},
\qquad
\varepsilon_q:=-p=-\frac{T+\rho_w}{4}.
}
\end{equation}
Then pointwise
\begin{equation}
T^\mu{}_{\nu}
=-\varepsilon_q\delta^\mu{}_{\nu}+\rho_wu^\mu u_{\nu}.
\end{equation}
This algebraic rewriting becomes a genuine separately conserved vacuum--dust decomposition only if, in addition,
\begin{equation}
\boxed{
\nabla_\mu\varepsilon_q=0,
\qquad
\nabla_\mu(\rho_wu^\mu)=0,
\qquad
u^\mu\nabla_\mu u^\nu=0.
}
\end{equation}
Thus the algebraic syzygies test perfect-fluid eigenvalue structure; the differential conditions select the endpoint vacuum--dust sector.
\end{theorem}
\begin{proof}
For an isotropic perfect fluid with eigenvalues \((-\rho,p,p,p)\), the traceless eigenvalues are \(-3(\rho+p)/4\) and \((\rho+p)/4\) with multiplicity three. Direct evaluation gives the two polynomial identities and the reconstruction formula for \(\rho+p\). Setting \(\rho_w=\rho+p\) and \(\varepsilon_q=-p\) gives the pointwise decomposition. No constancy or separate conservation follows from this algebra alone; those properties are precisely the additional differential conditions displayed above.
\end{proof}


The Bianchi identities imply
\begin{equation}
\nabla_\mu(T^{\rm rec})^{\mu\nu}=-\nabla_\mu(T^{\rm vis})^{\mu\nu}.
\end{equation}
If \(\nabla_\mu(T^{\rm vis})^{\mu\nu}=J^\nu\), then \(\nabla_\mu(T^{\rm rec})^{\mu\nu}=-J^\nu\). For \(J^\nu=0\), the reconstructed residual is conserved. This condition does not replace the equations for the internal variables of the source sector.

\section{Symmetries and Noether Identities}
Suppose a four-dimensional diffeomorphism-invariant completion \(S^{\rm eff}_{\rm fill}[g,\xi]\) exists and the \(\xi^I\) are scalar fields. The distinctions used here follow the standard variational and covariant phase-space separation of identities, currents, constraints, and charges \cite{Noether1918,LeeWald1990,IyerWald1994,BarnichBrandt2002}. We use the stress-tensor convention
\begin{equation}
T_{\mu\nu}:=-\frac{2}{\sqrt{-g}}\frac{\delta S}{\delta g^{\mu\nu}}.
\end{equation}
With conventions
\begin{equation}
\delta_\zeta \xi^I=\zeta^\mu\nabla_\mu\xi^I,
\qquad
\mathcal{E}_I:=\frac{1}{\sqrt{-g}}\frac{\delta S^{\rm eff}_{\rm fill}}{\delta\xi^I},
\end{equation}
diffeomorphism invariance gives
\begin{equation}
\nabla_\mu(T^{\rm src})^{\mu}{}_{\nu}=\sum_I\mathcal{E}_I\nabla_\nu\xi^I.
\end{equation}
On shell, \(\mathcal{E}_I=0\), and therefore \(\nabla_\mu(T^{\rm src})^\mu{}_{\nu}=0\). For tensor or gauge fields, the right-hand side is replaced by the contraction of the Euler-Lagrange expressions with the corresponding diffeomorphism generators. This is a variational identity, not an independent global current.

Four objects must be kept distinct: global current, gauge identity, constraint, surface charge. A superpotential belongs to the decomposition of a Noether current or a surface charge; it is not appended arbitrarily to the divergence of an energy-momentum tensor. A possible covariant boundary completion must likewise distinguish bulk currents from quasilocal surface charges.

\section{Radius, Asymptotic Fate, and Dilation Sign}
\begin{nogotheorem}[No topological radius bound]
The topological class \(\Sthree\) alone provides neither a lower nor an upper bound on the round radius \(A\).
\end{nogotheorem}
\begin{proof}
For every \(A>0\), \(g_A=A^2\dd\Omega_3^2\) is a smooth metric on the same manifold \(\Sthree\).
\end{proof}

The dilation spectrum does, however, classify the asymptotic proper-time behavior of a positive branch once that branch is assumed to dominate the GR Friedmann equation.

\begin{theorem}[Finite-future-time criterion for a dominant dilation branch]
Suppose a positive component \(\rho\propto a^s\) dominates at large \(a\) in an ever-expanding GR solution. Then
\begin{equation}
H\propto a^{s/2},
\qquad
\boxed{
\Delta\tau=\int^{\infty}\frac{\dd a}{aH(a)}
}\ \propto\ \int^\infty a^{-1-s/2}\,\dd a.
\end{equation}
The future proper time to \(a=\infty\) is finite if and only if
\begin{equation}
\boxed{s>0}
\end{equation}
or, equivalently, \(w<-1\).
\end{theorem}

\begin{theorem}[Late-time ordering in the physical boundary model]
Consider
\begin{equation}
\mathcal H^2(a)=\Omega_\Lambda+\Omega_Xa^{-1}+C_2a^{-2}+\Omega_ma^{-3}+\Omega_ra^{-4},
\qquad
C_2:=\Omega_K^{\rm geom}+\Omega_{2,\rm val},
\end{equation}
with nonnegative \(\Omega_\Lambda,\Omega_X,\Omega_m,\Omega_r\) and \(\Omega_K^{\rm geom}<0\). Then:
\begin{enumerate}
\item if \(\Omega_\Lambda>0\), the ever-expanding branch is asymptotically de Sitter;
\item if \(\Omega_\Lambda=0\) and \(\Omega_X>0\), then \(a(\tau)\propto\tau^2\) asymptotically;
\item if \(\Omega_\Lambda=\Omega_X=0\) and \(C_2>0\), then \(a(\tau)\propto\tau\) asymptotically;
\item if \(\Omega_\Lambda=\Omega_X=0\), \(C_2=0\), and \(\Omega_m>0\), then \(a(\tau)\propto\tau^{2/3}\) asymptotically; if matter is also absent but \(\Omega_r>0\), then \(a(\tau)\propto\tau^{1/2}\);
\item if \(\Omega_\Lambda=\Omega_X=0\) and \(C_2<0\), an expanding solution with \(\mathcal H^2(1)>0\) cannot extend to arbitrarily large \(a\): continuity implies at least one future zero of \(\mathcal H^2\).
\end{enumerate}
Thus the physical \(s=-2\) valuation amplitude can alter the global fate by competing with, but not becoming identical to, the negative closed-curvature coefficient.
\end{theorem}

Among the five positive valuation slots, only \(s=1\), \(w=-4/3\), can produce a Big-Rip-type finite-time divergence under asymptotic domination. All local curvature-measure boundary modes have \(s\le0\) and do not produce such a divergence by themselves. For the positive \(\Lambda+X\) model, \(\Omega_\Lambda>0\) gives an asymptotically de Sitter future; if \(\Omega_\Lambda=0\) and \(\Omega_X>0\), then \(a(\tau)\propto \tau^2\), so the future proper time remains infinite.

\section{Perturbations and Transfer-Grade Completion}
The background map does not determine perturbations automatically. The perturbative variables and closure requirements follow the standard cosmological perturbation framework \cite{KodamaSasaki1984,HalliwellHawking1985,MukhanovFeldmanBrandenberger1992,MalikWands2009}. A minimal set includes
\begin{equation}
\delta\rhoFour,
\qquad
\delta\epsD,
\qquad
\delta\pD,
\qquad
\delta\QD,
\end{equation}
metric scalar modes on \(\Sthree\), vector and tensor harmonics, and
\begin{equation}
c_s^2,
\qquad
\Gamma_D,
\qquad
\pi_{\mu\nu}.
\end{equation}
Here \(c_s^2\) is the effective sound speed in the chosen physical frame, \(\Gamma_D\) is the nonadiabatic perturbation, and \(\pi_{\mu\nu}\) is the anisotropic stress. Their determination requires a covariant completion or independent equations for the internal variables. On compact \(\Sthree\), a complete treatment must additionally respect the discrete scalar, vector, and tensor harmonic spectra.

For a generic filling functional, realizing a dust slot as cold dark matter requires at least
\begin{equation}
c_s^2=0,
\qquad
\pi_{\mu\nu}=0,
\end{equation}
together with geodesic motion of the effective flow. The conserved-current realization of the endpoint-selected \(A^{-3}\) branch satisfies these conditions. Non-minimal realizations must prove them independently. The vacuum slot likewise requires a covariant condition eliminating an independent clustering perturbation; the background equality \(w=-1\) alone does not prove this.

\begin{closurecondition}[Perturbative completion beyond the minimal sector]
For extensions beyond the vacuum plus conserved-dust realization, construct the perturbed filling functional and derive the scalar, vector, and tensor equations on \(\Sthree\), including \(\delta\pD\), \(c_s^2\), \(\Gamma_D\), \(\pi_{\mu\nu}\), and the exchange law.
\end{closurecondition}



\begin{center}
\small
\begin{tabular}{@{}p{0.12\linewidth}p{0.10\linewidth}p{0.16\linewidth}p{0.12\linewidth}p{0.12\linewidth}p{0.22\linewidth}@{}}
\toprule
Branch & \(w\) & source realization & \(c_{s,\rm rest}^2\) & anisotropic stress & transfer status \\
\midrule
\(s=1\) & \(-4/3\) & none supplied; phantom-like algebraic slot & -- & -- & removed by local curvature-measure filter; not a stable source sector \\
\(s=0\) & \(-1\) & vacuum term \(T_{\mu\nu}=-\varepsilon_qg_{\mu\nu}\) & none & none & source-realized; no clustering degree \\
\(s=-1\) & \(-2/3\) & canonical scalar VDE completion & 1 & 0 & potential and scalar kernel fixed; full Boltzmann likelihood open \\
\(s=-2\) & \(-1/3\) & canonical scalar existence realization & 1 & 0 & source-realized; distinguishes stress from curvature at perturbative level; transfer open \\
\(s=-3\) & 0 & pressureless conserved current & 0 & 0 & source-realized as CDM-like current; compact transfer hierarchy open \\
\bottomrule
\end{tabular}
\end{center}

\begin{remark}[Transfer-grade status]
The ledger separates source realization from transfer prediction. A transfer-grade compact cosmology still requires the scalar, vector, and tensor mode operators on \(\Sthree\), gauge and constraint reduction, initial mode covariance, photon-baryon-neutrino evolution, scalar-field and dust transfer functions, and CMB/LSS/lensing likelihood maps. Thus compact harmonic bookkeeping is not yet compact perturbative closure.
\end{remark}


\section{Two Independent Mathematical Bases and Endpoint Classification}
The construction rests on two independent mathematical structures that restrict admissible residual-sector solutions at different levels: the valuation closure of the auxiliary filling and the compact harmonic spectrum of the spatial representative. The former constrains homogeneous background scaling laws, while the latter constrains the perturbative mode content. The minimal endpoint classification then selects the two endpoint powers inside a stated local boundary class. None of these structures fixes dimensional amplitudes or establishes observational superiority by itself.

\subsection{Valuation closure}
Under the Euclidean convex valuation hypothesis stated above, Hadwiger's classification restricts the auxiliary filling energy to the intrinsic-volume span \cite{Hadwiger1957,Schneider2014}. For the smooth round homothetic family this is represented by
\begin{equation}
\Efour^{\rm val}(A)
=\gamma V_4(A)
+\sum_{j=0}^{3}\alpha_j
\int_{\Sthree_A}e_j(K)\,\dd\Sigma .
\end{equation}
With the outward-normal convention,
\begin{equation}
K^a{}_b=A^{-1}\delta^a_b,
\qquad
e_j(K)=\binom{3}{j}A^{-j},
\end{equation}
and therefore
\begin{equation}
\epsD(A)
=\frac{\gamma A}{4}
+\alpha_0
+\frac{3\alpha_1}{A}
+\frac{3\alpha_2}{A^2}
+\frac{\alpha_3}{A^3}.
\end{equation}
The associated pure branches have
\begin{equation}
w=-\frac{4}{3},\ -1,\ -\frac{2}{3},\ -\frac{1}{3},\ 0.
\end{equation}
Thus arbitrary homogeneous functions \(\epsD(A)\) are excluded inside the stated valuation class. This is a finite-spectrum theorem relative to the Euclidean convex valuation hypothesis, not a prohibition on other dark-sector constructions outside that class.

\subsection{Compact harmonic spectrum}
The second structure is harmonic analysis on the compact round spatial representative. For curvature radius \(R_c\), the scalar Laplacian has the discrete spectrum \cite{HalliwellHawking1985,KodamaSasaki1984}
\begin{equation}
\operatorname{Spec}(-\nabla^2_{\Sthree})
=
\left\{
\frac{n(n+2)}{R_c^2}
\right\}_{n=0}^{\infty},
\qquad
d_n=(n+1)^2,
\end{equation}
where \(d_n\) is the degeneracy of the scalar eigenspace. The mode \(n=0\) is homogeneous. Non-constant scalar harmonics begin at \(n=1\), subject to the usual gauge and constraint qualifications of cosmological perturbation theory. The harmonic spectrum is geometric; the physical perturbation spectrum is obtained only after constraints and gauge removal. Consequently, compact \(\Sthree\) does not possess a continuum of arbitrarily small nonzero scalar wavenumbers, unlike non-compact \(\mathbb{R}^3\).

This discreteness supplies a possible infrared spectral cutoff when the curvature scale, primordial spectrum, and transfer functions make it observationally accessible. It does not by itself imply suppression of a particular CMB multipole. Any observable infrared signature depends on \(R_c\), the allowed scalar, vector, and tensor harmonics, transfer functions, initial conditions, and any quotient topology \(\Sthree/\Gamma\).

\subsection{Synthesis}
The two structures constrain different layers of the model:
\begin{equation}
\boxed{
\text{valuation closure on }\Bfour_A
\quad\Longrightarrow\quad
\text{finite background equation-of-state spectrum}
}
\end{equation}
and
\begin{equation}
\boxed{
\begin{gathered}
\text{compact harmonics on }\Sthree_A\\
\Longrightarrow\ \text{discrete perturbation spectrum and possible IR signatures}
\end{gathered}
}.
\end{equation}
The minimal endpoint classification adds a separate endpoint selection:
\begin{equation}
\boxed{
\begin{gathered}
\text{minimal endpoint class}\\
\Longrightarrow\ E_{\rm end}=\sigma V_3+\mu\QU\\
\Longrightarrow\ \epsD(a)=\varepsilon_q+\varepsilon_w a^{-3},\quad \pD=-\varepsilon_q .
\end{gathered}
}
\end{equation}
Thus valuation closure constrains the background spectrum, the stated minimal endpoint class selects the vacuum-dust endpoint functional, and compact harmonics constrain the perturbative mode spectrum. Pure topology cannot determine the dimensional amplitudes and does not establish statistical preference over \(\Lambda\)CDM. The amplitudes are therefore matching coefficients supplied by independent physics or observation, while a predictive cosmological model additionally requires an explicit perturbative realization and direct likelihood-level comparison with current astronomical data.

\section{Observational Map, Empirical Pruning, and Falsification}
At background level, the minimal endpoint sector fixes the functional form of the residual sector within the stated endpoint class, but \(H(a)\) is determined only after the amplitudes and the visible sector have been specified. With radiation, baryons, the endpoint dust coefficient, and the vacuum coefficient, the minimal effective realization gives
\begin{equation}
H^2(a)=H_0^2\left[
\Omega_r a^{-4}
+\Omega_b a^{-3}
+\Omega_w a^{-3}
+\Omega_q
+\Omega_K^{\rm geom} a^{-2}
\right],
\qquad
\Omega_K^{\rm geom}=-\frac{c^2}{A_0^2H_0^2}<0
\end{equation}
for the closed convention used here. The density parameters satisfy
\begin{equation}
\Omega_r+\Omega_b+\Omega_w+\Omega_q+\Omega_K^{\rm geom}=1,
\qquad
\Omega_m=\Omega_b+\Omega_w.
\end{equation}
Thus, at the minimal effective level, the observational background is not a new expansion law beyond curved \(\Lambda\)CDM with an effective cold component. The novelty is the full valuation classification, its finite dilation-spectrum closure, and the minimal endpoint classification of the residual functional form.

\subsection{Empirical pruning of the valuation branches}
The analytical classification contains more directions than are currently useful as physical dark-energy candidates. Existing data should therefore be used as a pruning layer, not as a claim of detection.

The bulk branch \(s=1\), \(w=-4/3\), is retained only as a mathematical member of the full Euclidean valuation spectrum. It is phantom-like, lacks a stable source realization here, is removed by the local curvature-measure boundary filter, and is not used as a physical dark-energy branch. A pure \(s=-1\), \(w=-2/3\), component replacing \(\Lambda\) is likewise not viable as the whole dark energy. The full DES Y6 multi-probe combination gives a constant-\(w\) constraint close to \(\Lambda\), \(w=-0.981^{+0.021}_{-0.022}\), with no support for either pure \(w=-4/3\) or pure \(w=-2/3\) as the total dark-energy law \cite{DESY6ThreeByTwo2026}. The phenomenologically meaningful version is therefore the nested small-amplitude \(\Lambda+X\) model with \(\Omega_X\ge0\), tested against \(\Omega_X=0\).

A nearly direct published background analogue of the \(s=-1\) branch adds an effective domain-wall component with \(w=-2/3\), \(\rho\propto a^{-1}\), to \(\Lambda\)CDM and compares DESI DR2 BAO, compressed CMB information, and Pantheon+/Union3/DESY5 supernova compilations. It finds at most a few-percent component, with the largest DESY5 improvement \(\Delta\chi^2_{\rm MAP}=-1.72\), \(\Delta {\rm DIC}=-0.94\), and \(\ln B=0.32\pm0.32\), while Pantheon+ and Union3 give essentially no significant gain \cite{AnHanZhang2025}. This is a useful background benchmark, but it is not the canonical-scalar \(c_{s,\rm rest}^2=1\) realization used here and it does not include the exact compact-\(\Sthree\) transfer likelihood. It supports the status of \(\Lambda+X\) as a small nested null-test branch, not as a detection.

The physical \(s=-2\), \(w=-1/3\), branch is retained mainly as a curvature-splitting and null-test parameter. It does not by itself accelerate the universe, and current background comparisons do not require an additional positive \(a^{-2}\) stress component. The same topological-defect background study finds that the cosmic-string analogue with \(w=-1/3\) does not improve the fit, with \(\Delta\chi^2_{\rm MAP}\simeq0.19,0.21,0.037\) for Pantheon+, Union3, and DESY5 and negative Bayes factors \cite{AnHanZhang2025}. The observational baseline used below is therefore \(\Omega_{2,\rm val}=0\), with \(\Omega_{2,\rm val}\) released only in a separate curvature-splitting test.

The positive valuation flow condition \(\Dil w\le0\) implies a local tangent sign \(w_a\ge0\). Current DESI/DES dynamical-dark-energy fits prefer the opposite CPL quadrant \(w_0>-1\), \(w_a<0\). This means that positive valuation mixtures should be treated as sharp null tests rather than as explanations of the present CPL phantom-crossing trend. This sign comparison is not a direct exclusion of the exact \(\Lambda+X\) likelihood model, because CPL is a finite-redshift parameterization and the exact VDE trajectory must be fitted directly.

Finally, the endpoint vacuum--dust sector is observationally degenerate with ordinary \(\Lambda\)CDM once its amplitudes are free. Standard evidence for a CDM-like stress tensor is not evidence for the topological label \(\QU\); the latter becomes testable only after an independent amplitude or sector-matching prediction is supplied.

\subsection{Dilation-Closure Null Tests}
The differential characterization supplies sharp null tests once a residual density \(\epsD(a)\) has been reconstructed from the background under a declared visible-sector and curvature model. Define
\begin{align}
\mathcal R_{\rm val}[\epsD]
&:=(\Dil-1)\Dil(\Dil+1)(\Dil+2)(\Dil+3)\epsD,
\\
\mathcal R_{\rm bdry}[\epsD]
&:=\Dil(\Dil+1)(\Dil+2)(\Dil+3)\epsD,
\\
\mathcal R_{\rm end}[\epsD]
&:=\Dil(\Dil+3)\epsD.
\end{align}
The exact full valuation, local curvature-measure, and minimal endpoint classes satisfy respectively
\begin{equation}
\mathcal R_{\rm val}=0,
\qquad
\mathcal R_{\rm bdry}=0,
\qquad
\mathcal R_{\rm end}=0.
\end{equation}
Inside the full valuation class, the projected amplitudes must also satisfy
\begin{equation}
\frac{\dd\widehat{\mathcal I}_s}{\dd\ln a}=0,
\qquad
s\in\{1,0,-1,-2,-3\}.
\end{equation}
For the minimal endpoint sector the sharper conditions are
\begin{equation}
\widehat{\mathcal I}_1
=\widehat{\mathcal I}_{-1}
=\widehat{\mathcal I}_{-2}=0,
\qquad
\frac{\dd\widehat{\mathcal I}_0}{\dd\ln a}=0,
\qquad
\frac{\dd\widehat{\mathcal I}_{-3}}{\dd\ln a}=0.
\end{equation}

\begin{criterion}[Dilation-spectrum closure test]
A reconstructed residual history is consistent with a declared homothetic closure class only if the corresponding annihilator residual is statistically consistent with zero and the surviving projector amplitudes are consistent with constants over the tested interval, with uncertainties propagated through the full covariance of the reconstruction. Because high-order numerical differentiation is noise-sensitive, a practical implementation should infer these quantities from a covariance-aware basis fit, Gaussian-process reconstruction, or direct likelihood parameterization rather than from unregularized finite differences.
\end{criterion}

These null tests determine membership in the declared functional class; they do not by themselves identify a microscopic source or establish observational superiority.


\subsection{Algorithmic closure protocols for lifts, curvature, and null tests}
\label{sec:algorithmic-closure-protocols}
The preceding invariants can be translated into a closed algorithmic specification for CLASS-like or likelihood-level implementations.  This is not a claim that the full compact Boltzmann hierarchy has already been computed.  It is a declaration of the minimal mathematical interface between the homogeneous valuation alphabet, the chosen perturbative lift, and the statistical test.

\begin{criterion}[Lift-interface closure protocol]
The implementation of a valuation branch must be specified in three logically separate steps.

\emph{Step 1: background alphabet.}  The general physical local curvature-measure boundary background is
\begin{equation}
\boxed{
\mathcal H^2_{\partial}(a)=
\Omega_r a^{-4}+\Omega_m a^{-3}
+\bigl(\Omega_K^{\rm geom}+\Omega_{2,\rm val}\bigr)a^{-2}
+\Omega_Xa^{-1}
+\Omega_\Lambda.
}
\end{equation}
The two \(a^{-2}\) contributions are not separable from \(\mathcal H^2(a)\) alone; they are separated only after transverse-distance information fixes \(\Omega_K^{\rm geom}\).
The nested one-parameter canonical VDE test imposes
\begin{equation}
\boxed{\Omega_{2,\rm val}=0,\qquad \Omega_X\ge0.}
\end{equation}
Thus the background alphabet is fixed by the dilation spectrum, but the perturbations are not.

\emph{Step 2: lift declaration.}  A lift must be declared before perturbations are evolved.  The admissible examples constructed in this paper include the canonical scalar lift for the positive \(s=-1\) VDE branch, the canonical scalar lift for a positive physical \(s=-2\) branch, the conserved-current dust lift for \(s=-3\), and the ADM geometric-response prototypes for the \(R^{(3)}\) and \(\sqrt{R^{(3)}}\) branches.  The same background power with two different lifts is not the same perturbative model.

\emph{Step 3: perturbation grammar.}  For the canonical scalar lift one evolves the scalar perturbation variables, for example \(\delta\phi_{s,n}\) and its time derivative, with no free sound speed, anisotropic stress, or independent mass function.  The free scalar operator is fixed by
\begin{equation}
\boxed{
\omega_{s,n}^2(a)=\frac{n(n+2)}{A^2}+V_{s,\phi\phi},
\qquad
V_{s,\phi\phi}
=-\frac{6+s}{4}H^2
\left[s+\Dil\ln\left(\frac{\mathcal H^2}{\mathcal H_*^2}\right)\right].
}
\end{equation}
For the ADM-response lifts no autonomous fluid hierarchy is implied.  The response is algebraic in the metric potentials and must be inserted as a modified constraint/geometric-response kernel rather than as an ordinary conserved CLASS species.
\end{criterion}

Equivalently,
\begin{equation}
\boxed{
\begin{array}{c}
\text{valuation spectrum}\Rightarrow\text{allowed background weights},\\[1mm]
\text{lift choice}\Rightarrow\delta T_{\mu\nu}\ \text{or ADM response kernel},\\[1mm]
\text{Boltzmann solver}\Rightarrow\text{transfer functions and likelihood}.
\end{array}}
\end{equation}

The valuation algebra supplies the finite background alphabet; the lift supplies the grammar of perturbations.

\subsubsection{Convention-safe closed-\texorpdfstring{\(\Sthree\)}{S3} scalar implementation}
\label{sec:convention-safe-numerical-domain}
For the nested canonical VDE lift, the source equations and their coefficients are fully specified at linear scalar order, but a numerical implementation still requires a declared closed-harmonic convention and the complete closed-FLRW constraint system.  In the convention used in this paper,
\begin{equation}
\boxed{
-\nabla^2_{\Sthree_A}Y_{n\ell m}
=\frac{n(n+2)}{A^2}Y_{n\ell m},
\qquad
n=0,1,2,\ldots,
\quad
\ell=0,\ldots,n,
\quad
m=-\ell,\ldots,\ell .}
\label{eq:closed-harmonic-convention-final}
\end{equation}
The degeneracy is \(d_n=(n+1)^2\).  The mode \(n=0\) is homogeneous, while the lowest non-constant sectors, especially \(n=1\), must be handled with the usual closed-FLRW gauge and constraint qualifications.  Alternative conventions that relabel the principal mode by one unit are allowed only if the eigenvalue and degeneracy formulas are relabelled consistently.

Let \(\eta\) denote conformal time and define the conformal Hubble rate by \(\mathscr H_c:=aH\), to avoid confusing it with the normalized expansion function \(\mathcal H^2=H^2/H_0^2\).  For the one-parameter canonical VDE lift,
\begin{equation}
F_{\rm VDE}(a)=
\Omega_r a^{-4}+\Omega_m a^{-3}+\Omega_K^{\rm geom}a^{-2}
+\Omega_\Lambda+\Omega_Xa^{-1},
\end{equation}
\begin{equation}
\mathscr W_n(a)=
(-\Omega_K^{\rm geom})a^{-2}n(n+2)
+\frac54\left[F_{\rm VDE}(a)-\Dil F_{\rm VDE}(a)\right].
\end{equation}
The second-order scalar equation \cref{eq:VDE-delta-phi} is equivalently the first-order system
\begin{align}
\frac{\dd\delta\phi_n}{\dd\eta}&=v_n,\label{eq:VDE-first-order-1}\\
\frac{\dd v_n}{\dd\eta}
&=-2\mathscr H_c v_n
-a^2H_0^2\mathscr W_n\delta\phi_n
+\phi'\left(\Psi_n'+3\Phi_n'\right)
-2a^2U_{,\phi}\Psi_n,
\label{eq:VDE-first-order-2}
\end{align}
where primes denote \(\dd/\dd\eta\).  The background source coefficients are fixed by the inverse reconstruction,
\begin{equation}
\boxed{
(\phi')^2=a^2M_{\rm Pl}^2H_0^2\Omega_Xa^{-1},
\qquad
U_{,\phi}
=\mp\frac52M_{\rm Pl}H_0^2
\sqrt{\Omega_Xa^{-1}F_{\rm VDE}(a)}.}
\label{eq:VDE-conformal-source-coefficients}
\end{equation}
The sign corresponds to the two monotone canonical branches and must be chosen consistently with the homogeneous trajectory.  No independent sound-speed, anisotropic-stress, or scalar-mass function is present:
\begin{equation}
c_{s,\phi,{\rm rest}}^2=1,
\qquad
\pi^{(\phi)}_{\mu\nu}=0,
\qquad
m_\phi^2/H^2=\frac54(1-\Dil\ln F_{\rm VDE}).
\end{equation}

The scalar stress sources entering the closed-FLRW Einstein constraints are those in \cref{eq:VDE-delta-phi}: \(\delta\rho_{\phi,n}\), \(\delta p_{\phi,n}\), \(\delta q_{\phi,n}\), and \(\pi_{\mu\nu}^{(\phi)}=0\).  The schematic replacement \(k^2\mapsto n(n+2)/A^2\) identifies the compact scalar spectrum only.  A transfer code must use the standard closed-FLRW Einstein--Boltzmann constraint equations with their curvature-dependent coefficients, gauge-mode exclusions, and visible-sector hierarchies.  Thus \cref{eq:VDE-first-order-1,eq:VDE-first-order-2} are source equations ready for implementation after a closed-FLRW harmonic convention and constraint system are fixed; they are not a substitute for the full closed-FLRW Einstein--Boltzmann system.

\begin{remark}[Bianchi IX matrix elements are convention dependent]
For the fixed-shape Bianchi IX diagnostic, the authorized first-order matrix is
\begin{equation}
\mathcal M^{\rm IX}_{IJ}=-2\sum_{i=1}^3\beta_i\langle I|L_i^2|J\rangle .
\end{equation}
The matrix elements are to be computed in a declared orthonormal \(\Sthree\) harmonic basis.  They are block-diagonal in the principal round Laplacian eigenspaces and real symmetric, but their explicit sparse form depends on whether one uses a Wigner-\(D\), coupled \(SO(4)\), or hyperspherical \((n,\ell,m)\) basis.  They should not be replaced by flat-space or ordinary rigid-rotor formulas unless the chosen basis and representation map justify that identification.
\end{remark}

\begin{criterion}[Numerical-domain freeze]
The homogeneous plus canonical linear-scalar sector is theoretically closed once the amplitudes, lift, harmonic convention, and initial-condition class are declared.  The remaining work is numerical implementation in a chosen closed-FLRW harmonic convention.  In particular,
\begin{equation}
\boxed{
\begin{array}{c}
\text{valuation algebra}\Rightarrow\text{background weights},\\[1mm]
\text{canonical lift}\Rightarrow\delta\phi_n\text{ system},\\[1mm]
\text{closed-FLRW code}\Rightarrow\text{observables}.
\end{array}}
\end{equation}
The Bianchi IX matrix kernel is a robustness diagnostic whose isotropic limit must reduce to \(\mathscr W_n\); it is not a nonlinear Mixmaster transfer theory.
\end{criterion}

\begin{criterion}[Statistical closure protocol]
For the nested positive one-parameter VDE test
\begin{equation}
H_0:\Omega_X=0,
\qquad
H_1:\Omega_X\ge0,
\end{equation}
the likelihood-ratio statistic
\begin{equation}
q_X=2\left[\ln L(\widehat\Omega_X)-\ln L(0)\right]
\end{equation}
has, in the exact linear Gaussian benchmark with known covariance and full-rank nuisance space,
\begin{equation}
\boxed{
q_X\sim\frac12\delta_0+\frac12\chi_1^2,
\qquad
p_{\rm one\text{-}sided}=1-\Phi_{\rm N}(\sqrt{q_X})\quad(q_X>0).
}
\end{equation}
The covariance-optimal quotient coordinate is obtained by profiling out the nuisance subspace.  If
\begin{equation}
y=S\beta+x\Omega_X+\eta,
\qquad
W=C^{-1},
\qquad
M_S=I-S(S^TWS)^{-1}S^TW,
\end{equation}
then
\begin{equation}
\boxed{
\widehat\Omega_X^{\rm GLS}=
\frac{x^TWM_Sy}{x^TWM_Sx},
\qquad
\operatorname{Var}(\widehat\Omega_X^{\rm GLS})=
\frac{1}{x^TWM_Sx}.
}
\end{equation}
The corresponding Gaussian closure residual for a declared spectral subspace with design matrix \(X_S\) is
\begin{equation}
\boxed{
Q_S=(y-X_S\widehat\beta_S)^TW(y-X_S\widehat\beta_S),
\qquad
\widehat\beta_S=(X_S^TWX_S)^{-1}X_S^TWy,
}
\end{equation}
with the usual chi-square calibration only under the stated Gaussian, known-covariance, full-rank assumptions.  Full nonlinear cosmological likelihoods require boundary asymptotics or simulation calibration.
\end{criterion}

\begin{criterion}[Geometric-curvature splitting protocol]
The physical \(a^{-2}\) valuation stress and the spatial curvature of the FLRW metric are separated by expansion and transverse-distance data.  With
\begin{equation}
d(z):=\frac{H_0D_M(z)}{c},
\end{equation}
the Clarkson--Bassett--Lu FLRW consistency relation gives, for \(z>0\) or by the regular \(z\to0\) limit,
\begin{equation}
\boxed{
\Omega_K^{\rm geom}(z)=
\frac{[\mathcal H(z)d'(z)]^2-1}{d(z)^2}=\text{constant}.
}
\end{equation}
The total coefficient of the \(a^{-2}\) term in \(\mathcal H^2\) is extracted by
\begin{equation}
\boxed{
\mathcal C_{-2}[\mathcal H^2]
=\frac{a^2}{4}\Dil(\Dil+1)(\Dil+3)(\Dil+4)\mathcal H^2,
}
\end{equation}
and the physical valuation contribution is
\begin{equation}
\boxed{
\Omega_{2,\rm val}=\mathcal C_{-2}[\mathcal H^2]-\Omega_K^{\rm geom}.
}
\end{equation}
The compact scalar Laplacian then supplies the spectral-geometric invariant
\begin{equation}
\boxed{
\frac{c^2\lambda_n}{H_0^2}=(-\Omega_K^{\rm geom})n(n+2),
\qquad
\lambda_n=\frac{n(n+2)}{A_0^2}.
}
\end{equation}
Thus geometric curvature is not inferred from the \(a^{-2}\) coefficient of \(\mathcal H^2\) alone.
\end{criterion}

\begin{criterion}[Dilation-annihilator closure protocol]
At the density-profile level, membership in the declared valuation classes is tested by
\begin{align}
\boxed{\mathcal R_{\rm val}[\epsD]
=(\Dil-1)\Dil(\Dil+1)(\Dil+2)(\Dil+3)\epsD=0,}\\
\boxed{\mathcal R_{\rm bdry}[\epsD]
=\Dil(\Dil+1)(\Dil+2)(\Dil+3)\epsD=0,}\\
\boxed{\mathcal R_{\rm end}[\epsD]
=\Dil(\Dil+3)\epsD=0.}
\end{align}
For the normalized expansion function, the nested VDE coefficient is extracted by
\begin{equation}
\boxed{
\widehat\Omega_X(a)=
-\frac a6\Dil(\Dil+2)(\Dil+3)(\Dil+4)\mathcal H^2(a),
}
\end{equation}
and must be statistically consistent with a constant.  Curved \(\Lambda\textrm{CDM}\) is the sharper null \(\widehat\Omega_X(a)=0\).  These operator criteria define closure of the homogeneous functional class, not a microscopic source theory.
\end{criterion}

\begin{nogotheorem}[Exact homogeneous-background degeneracy]
For the minimal endpoint realization, after identifying
\begin{equation}
\Omega_m=\Omega_b+\Omega_w,
\qquad
\Omega_\Lambda=\Omega_q,
\end{equation}
and keeping the same geometric curvature \(\Omega_K^{\rm geom}\), every homogeneous FLRW observable constructed from the same geometric curvature, the same expansion history, and the standard curved-FLRW distance-redshift relation is exactly degenerate with curved \(\Lambda\)CDM. Any statistical preference must therefore come from an additional prior, amplitude principle, compact spectral observable, perturbative observable, or non-minimal branch.
\end{nogotheorem}
\begin{proof}
The endpoint-sector expansion law is exactly that of curved \(\Lambda\)CDM after the parameter identifications above, with the same geometric curvature parameter. Radial distances depend on \(H^{-1}(z)\), while transverse distances are obtained from the same radial integral through the identical curved-FLRW sine/sinh geometry fixed by \(\Omega_K^{\rm geom}\). Therefore all homogeneous FLRW observables built from the expansion history and the standard curved distance-redshift map coincide after parameter relabelling. Distinguishability requires information beyond this homogeneous sector.
\end{proof}


\subsection{Final invariant dictionary and scope ledger}
\label{sec:final-invariant-dictionary}

We now collect the invariant dictionary in a form that is independent of any
particular observational data vector.  The entries below are either
reference-normalization independent, dimensionless, or explicitly marked as
external matching amplitudes.  This dictionary is the summary algebraic
interface of the framework: it does not close the microscopic origin of the
couplings and it does not replace a Boltzmann likelihood.

\begin{criterion}[Complete homogeneous invariant dictionary within the declared valuation problem]
The following quantities are the finite invariant data of the homogeneous
valuation profile and of the declared source lifts.
\begin{center}
\small
\begin{tabular}{p{0.27\linewidth}p{0.64\linewidth}}
\toprule
Layer & Invariant content \\
\midrule
Spectral support &
\(S_{\rm val}=\{1,0,-1,-2,-3\}\),
\(P_{\rm val}(z)=(z-1)z(z+1)(z+2)(z+3)\), and
\(\mathfrak A_{\Dil}\simeq\mathbb R[z]/(P_{\rm val})\simeq\mathbb R^5\). \\
Projectors and coefficients &
Primitive idempotents
\(\mathsf P_s=\prod_{t\neq s}(\Dil-t)/(s-t)\),
reference-normalization-independent coefficients
\(\widehat{\mathcal I}_s=A^{-s}\mathsf P_s\epsD\), and dimensionless fractions
\(\pi_s=\mathsf P_s\epsD/\epsD\), \(\sum_s\pi_s=1\). \\
Moments and convexity &
\(\Dil^n\epsD/\epsD=\sum_s s^n\pi_s\), cumulants of
\(K=\ln(\epsD/\varepsilon_*)\),
\(\Dil^2K=\operatorname{Var}_\pi(s)\ge0\), and
\(\Dil w_D=-\operatorname{Var}_\pi(s)/3\le0\) under nonnegative active branches. \\
Diameter bounds &
Full positive spectrum:
\(-4/3\le\Dil w\le0\); boundary positive spectrum:
\(-3/4\le\Dil w\le0\); positive \(\Lambda+X\):
\(-1/12\le\Dil w\le0\), equivalently \(0\le w_a\le1/12\) at \(a=1\). \\
Rank and Prony tests &
Hankel matrices
\(H_r=[\Dil^{i+j}\epsD]_{i,j=0}^r\succeq0\) under positivity, with rank equal to the number of active branches; the discrete Prony recurrence is the sampled counterpart of the same annihilator. \\
Annihilators &
Full valuation:
\(P_{\rm val}(\Dil)\epsD=0\); boundary valuation:
\(\Dil(\Dil+1)(\Dil+2)(\Dil+3)\epsD=0\), equivalently
\(\widehat{\mathcal I}_1=0\); minimal endpoint:
\(\Dil(\Dil+3)\epsD=0\), equivalently
\(\widehat{\mathcal I}_1=\widehat{\mathcal I}_{-1}=\widehat{\mathcal I}_{-2}=0\). \\
Curvature split &
\(\mathcal C_{-2}[\mathcal H^2]
=\frac{a^2}{4}\Dil(\Dil+1)(\Dil+3)(\Dil+4)\mathcal H^2\),
\(\Omega_{2,\rm val}=\mathcal C_{-2}[\mathcal H^2]-\Omega_K^{\rm geom}\), and
\(\Omega_K^{\rm geom}=([\mathcal H d']^2-1)/d^2\). \\
Compact spectral geometry &
\(\lambda_n=n(n+2)/A_0^2\) and
\(c^2\lambda_n/H_0^2=(-\Omega_K^{\rm geom})n(n+2)\). \\
Canonical source lifts &
For \(s<0\),
\((\dd\phi_s/\dd\tau)^2=-s\rho_s/3\),
\(V_s=(6+s)\rho_s/6\),
\(c_{s,\rm rest}^2=1\), \(\pi_{\mu\nu}=0\), and
\(V_{s,\phi\phi}=-(6+s)H^2[s+\Dil\ln(\mathcal H^2/\mathcal H_*^2)]/4\). \\
General canonical branch kernel &
For \(F=H^2/H_0^2\) and \(s<0\),
\(\mathscr W_{s,n}=(-\Omega_K^{\rm geom})a^{-2}n(n+2)
-(6+s)[sF+\Dil F]/4\).  In particular,
\(m_{-2}^2/H_0^2=2F-\Dil F\). \\
Amplitude ledger &
The five valuation coefficients map exactly to normalized profile amplitudes
once \(H_0\) and \(A_0\) are supplied; this is a matching dictionary, not an
amplitude selector. \\
Canonical VDE kernel &
For \(s=-1\),
\(m_X^2/H^2=\frac54[1-\Dil\ln F_{\rm VDE}]\) and
\(\mathscr W_n=(-\Omega_K^{\rm geom})a^{-2}n(n+2)+\frac54(F_{\rm VDE}-\Dil F_{\rm VDE})\). \\
Endpoint charge &
Inside the declared minimal endpoint class,
\(E_{\rm end}=\sigma V_3+\mu\QU\), with matching
\(\mathcal E_{\rm rest}\overset{\rm match}{=}\mu\QU\); \(\sigma\) and \(\mu\) are dimensional matching data, not topological predictions. \\
\bottomrule
\end{tabular}
\end{center}
These entries constitute the homogeneous invariant dictionary.  They do not
assert that a unique microscopic source action follows from
\(\Efour^{\rm val}\).
\end{criterion}

\subsection{Elevated rigidity theorems}
\label{sec:elevated-rigidity-theorems}

The preceding dictionary can be sharpened into a finite list of rigidity
statements.  Each theorem below stays inside the Euclidean convex valuation
hypothesis, the homogeneous dilation module, and the separately declared source
lifts.  None of them derives a covariant source action from \(\Efour^{\rm val}\),
selects dimensional amplitudes, or replaces the closed-FLRW Boltzmann
likelihood.

\begin{theorem}[Complete spectral rigidity of the round valuation module]
Let
\begin{equation}
\mathcal V_{\rm val}=\operatorname{span}\{A,1,A^{-1},A^{-2},A^{-3}\}.
\end{equation}
The polynomial action
\begin{equation}
\varrho:\mathbb R[z]\longrightarrow\operatorname{End}(\mathcal V_{\rm val}),
\qquad p\longmapsto p(\Dil)|_{\mathcal V_{\rm val}},
\end{equation}
factors faithfully through
\begin{equation}
\mathbb R[z]/(P_{\rm val})\simeq\mathbb R^5,
\qquad
P_{\rm val}(z)=(z-1)z(z+1)(z+2)(z+3).
\end{equation}
Consequently the Lagrange idempotents \(\mathsf P_s=e_s(\Dil)\) are the unique
complete set of primitive mutually orthogonal spectral projectors on
\(\mathcal V_{\rm val}\), up to a permutation of the five labelled roots.  A
profile with all five branch coefficients nonzero is a cyclic vector for this
representation; profiles with smaller support are cyclic for their corresponding
support submodules.
\end{theorem}
\begin{proof}
The ODE characterization gives the exact minimal polynomial of \(\Dil\) on
\(\mathcal V_{\rm val}\), namely \(P_{\rm val}\).  Since the five roots are real
and distinct, the quotient is semisimple and decomposes by the Chinese
remainder theorem into the direct product of five one-dimensional factors.  The
primitive central idempotents of this product are the Lagrange factors
\(\mathsf P_s\), and their images give the one-dimensional eigenspaces
\(\mathbb R A^s\).  Any complete primitive orthogonal resolution of the identity
must select these simple factors, up to relabelling.
\end{proof}

\begin{theorem}[Closed positive branch cone and dilation-flow invariance]
Let
\begin{equation}
\overline{\mathcal C}_+=\left\{\epsD=\sum_{s\in S_{\rm val}}C_sA^s:
C_s\ge0\right\}
\end{equation}
be the closed branch-coefficient cone, including the zero profile.  Its nonzero
part
\begin{equation}
\mathcal C_+^\ast=\overline{\mathcal C}_+\setminus\{0\}
\end{equation}
is the nonzero positive part on which the normalized fractions
\(\pi_s=\mathsf P_s\epsD/\epsD\) are defined.  The finite dilation flow
\(\exp(t\Dil)\) preserves \(\overline{\mathcal C}_+\), \(\mathcal C_+^\ast\),
every support face, and the relative interior of every support face.  The local
curvature-measure boundary cone is the invariant face obtained by setting
\(C_1=0\).
\end{theorem}
\begin{proof}
In branch coordinates the inequalities \(C_s\ge0\) are linear and include the
origin, so \(\overline{\mathcal C}_+\) is a closed convex cone.  Removing only the
origin gives the nonzero cone on which \(\epsD(A)>0\) for all \(A>0\) and the
fractions \(\pi_s\) are well-defined.  The finite dilation flow acts as
\begin{equation}
\exp(t\Dil)\epsD(A)=\epsD(e^tA)=\sum_s e^{st}C_sA^s.
\end{equation}
For all real \(t\), the factors \(e^{st}\) are strictly positive.  Hence the sign
pattern of the branch coefficients is preserved, every vanishing coefficient
remains zero, and every positive coefficient remains positive.  Therefore the
closed cone, its nonzero part, each support face, and each relative interior are
invariant.  The boundary filter is exactly the support restriction excluding the
\(s=1\) ray.
\end{proof}

\begin{theorem}[Unique mass-flow invariant for canonical negative branches]
Fix \(s<0\) and a positive density branch \(\rho_s=\rho_{s0}a^s\).  Inside the
class of minimally coupled canonical scalars with fixed positive kinetic
normalization, the reconstructed trajectory is unique up to field sign and
constant shift, and the second derivative of the potential along it is fixed by
\begin{equation}
\boxed{
V_{s,\phi\phi}
=-\frac{6+s}{4}H^2
\left[
 s+\Dil\ln\left(\frac{\mathcal H^2}{\mathcal H_*^2}\right)
\right].
}
\end{equation}
No independent function of \(a\) is compatible with the same homogeneous branch
and the declared canonical normalization.
\end{theorem}
\begin{proof}
For a canonical scalar,
\(\rho_\phi=\frac12\dot\phi^2+V\) and
\(p_\phi=\frac12\dot\phi^2-V\).  With
\(w_s=-1-s/3\), the algebraic relations give
\begin{equation}
\dot\phi_s^2=(1+w_s)\rho_s=-\frac{s}{3}\rho_s,
\qquad
V_s=\frac{1-w_s}{2}\rho_s=\frac{6+s}{6}\rho_s.
\end{equation}
Since \(\rho_s>0\) and \(s<0\), \(\phi_s\) is monotone on any interval where
\(H>0\).  Differentiating \(V_s(a)\) twice with respect to this field and using
\(\mathcal H^2=H^2/H_0^2\) yields the displayed expression.  The only residual
freedom is the sign and origin of \(\phi\), neither of which changes
\(V_{\phi\phi}\).
\end{proof}

\begin{theorem}[Principal annihilator ideals and the spectral lattice]
For every nonempty subset \(S\subseteq S_{\rm val}\), define
\begin{equation}
\mathcal V_S=\operatorname{span}\{A^s:s\in S\},
\qquad
m_S(z)=\prod_{s\in S}(z-s).
\end{equation}
The annihilator ideal of \(\mathcal V_S\) is the principal ideal
\((m_S)\subset\mathbb R[z]\).  Under inclusion of subspaces the correspondence
\(\mathcal V_S\mapsto(m_S)\) is order reversing; equivalently, the spectral
subspace lattice is anti-isomorphic to the divisor lattice of \(P_{\rm val}\).
\end{theorem}
\begin{proof}
The distinct powers \(A^s\) are linearly independent on every open interval in
\((0,\infty)\).  A polynomial \(p(\Dil)\) annihilates \(\mathcal V_S\) exactly
when \(p(s)=0\) for every \(s\in S\).  Therefore every such polynomial is a
multiple of \(m_S\), and \(m_S\) itself annihilates \(\mathcal V_S\).  If
\(S\subseteq T\), then \(\mathcal V_S\subseteq\mathcal V_T\) and
\(m_S\) divides \(m_T\), so \((m_T)\subseteq(m_S)\).
\end{proof}

\begin{theorem}[Endpoint saturation of the boundary Bhatia--Davis bound]
On the positive local curvature-measure boundary cone with
\(S_\partial=\{0,-1,-2,-3\}\subset[-3,0]\), one has
\begin{equation}
\boxed{
0\le-\Dil w\le -3w(1+w).
}
\end{equation}
Equality in the upper bound at every scale occurs if and only if the active
support is contained in the endpoint pair \(\{-3,0\}\).
\end{theorem}
\begin{proof}
Let \(\mu_1=\sum_s s\pi_s\).  The Bhatia--Davis bound for a probability measure
supported on \([-3,0]\) gives
\begin{equation}
\operatorname{Var}_\pi(s)\le (0-\mu_1)(\mu_1+3).
\end{equation}
Because \(w=-1-\mu_1/3\) and \(-\Dil w=\operatorname{Var}_\pi(s)/3\), this is
exactly the displayed inequality.  The equality case of Bhatia--Davis is a
measure supported on the two endpoints of the interval; here those endpoints
are \(-3\) and \(0\).  Requiring equality along the whole flow therefore keeps
the profile in the minimal endpoint subspace.
\end{proof}

\begin{theorem}[First-order fixed-shape anisotropic deformation of the VDE kernel]
Within the fixed-shape Bianchi-IX metric-Laplacian lift considered here, the
first-order deformation of the round canonical VDE free kernel is
\begin{equation}
\boxed{
\mathscr W_{IJ}^{\rm IX}(a,\beta)=
(-\Omega_K^{\rm geom})a^{-2}
\left[n(n+2)\delta_{IJ}+\mathcal M_{IJ}^{\rm IX}\right]
+\frac54\left(F_{\rm VDE}-\Dil F_{\rm VDE}\right)\delta_{IJ}
+O(\beta^2),
}
\end{equation}
where
\begin{equation}
\mathcal M_{IJ}^{\rm IX}=-2\sum_{i=1}^3\beta_i\langle I|L_i^2|J\rangle
\end{equation}
is computed in the declared orthonormal \(\Sthree\) harmonic basis.  The
isotropic limit \(\beta_i\to0\) recovers the diagonal round kernel.
\end{theorem}
\begin{proof}
At fixed shape the grading identities give \([\Dil,L_i]=0\) and
\([\Dil,A^2\Delta_{\rm IX}]=0\), so the first-order deformation of the Laplacian
is scale-homogeneous and is exhausted, within the chosen metric-Laplacian lift,
by the displayed matrix element of the shape perturbation.  The mass-flow term
is a homogeneous background scalar and therefore remains diagonal in the
harmonic labels at this order.  Setting \(\beta_i=0\) removes
\(\mathcal M^{\rm IX}\) and gives the round \(\mathscr W_n\).
\end{proof}

\begin{theorem}[Sharp scale-covariance no-go]
No rule using only topological and dimensionless data of an oriented spatial
three-manifold can assign a nonzero quantity of nonzero length dimension while
remaining covariant under global metric rescalings.  Thus dimensional
amplitudes such as \(\sigma\), \(\mu\), \(H_0\), unnormalized branch densities,
mass scales, or cutoffs require external dimensional matching data, a
microscopic scale, or observational input.
\end{theorem}
\begin{proof}
A global metric rescaling by \(\lambda>0\) leaves topological and dimensionless
data unchanged.  A nonzero quantity of length dimension \(L^k\), with
\(k\neq0\), rescales by \(\lambda^k\).  If the rule depends only on the
unchanged data and is covariant under all rescalings, the assigned number must
be equal to \(\lambda^k\) times itself for every \(\lambda>0\).  Hence it is zero.
\end{proof}
\begin{remark}[Dimensionless density parameters]
Dimensionless parameters such as \(\Omega_X\) and \(\Omega_{2,\rm val}\) are not
fixed by this dimensional argument alone.  They become meaningful only after a
normalization scale, critical density, or independent amplitude ledger has been
declared, and are then constrained by matching or data rather than by topology
alone.
\end{remark}

\begin{theorem}[Unique closures inside declared lift classes]
Within the hypotheses stated in this paper, the following uniqueness statements
are closed.
\begin{enumerate}
\item Positive valuation mixtures have monotone equation-of-state flow
\(\Dil w\le0\), so every fixed threshold of \(w\) is crossed at most once as
long as all active branch energies are nonnegative.
\item The nested canonical VDE branch with \(\Omega_X>0\) has a unique
canonical trajectory and potential up to field sign and additive field shift:
\begin{equation}
\phi(a)=\phi_*\pm M_{\rm Pl}\sqrt{\Omega_X}
\int_1^a\frac{\dd\tilde a}{\tilde a^{3/2}\sqrt{F_{\rm VDE}(\tilde a)}},
\qquad
U(\phi(a))=3M_{\rm Pl}^2H_0^2\Omega_\Lambda
+\frac52M_{\rm Pl}^2H_0^2\Omega_Xa^{-1}.
\end{equation}
Its scalar mass, rest-frame sound speed, anisotropic stress, early-time
fraction, and compact scalar free kernel are then fixed.
\item Inside the oriented minimal endpoint class, the pressureless source
completion that gives CDM perturbative behavior is a conserved current lift,
with matching \(\mathcal E_{\rm rest}\overset{\rm match}{=}\mu\QU\).  The
matching is conserved once imposed, but it is not an amplitude derivation.
\item The ADM \(R^{(3)}\) and \(\sqrt{R^{(3)}}\) prototypes give unique
spatial response tensors for their declared functionals.  They are geometric
response diagnostics, inequivalent to canonical scalar or current matter
lifts, and not autonomous conserved fluid hierarchies.
\end{enumerate}
No additional nonzero dimensional amplitude is selected by the homogeneous
valuation algebra or by topology alone.
\end{theorem}
\begin{proof}
The first statement is the variance identity for positive branch fractions.
The second is the standard inverse canonical-scalar reconstruction applied to
the already fixed \(s=-1\) density branch, with the field sign and field-origin
freedom left unfixed.  The third is the Brown-type conserved-current
realization of pressureless dust together with the separately imposed global
matching condition.  The fourth follows by direct variation of the displayed
ADM spatial functionals.  The final sentence is precisely the scale-covariance
no-go: dimensionless topological data are invariant under global metric
rescaling, whereas any nonzero physical amplitude is not.
\end{proof}

\begin{criterion}[Sign-constrained observational falsifier]
For nonnegative active valuation branches, the sign theorem
\(\Dil w\le0\) implies the local CPL-tangent sign \(w_a\ge0\).  Present DESI
DR2 and DES combinations prefer the opposite CPL quadrant
\(w_0>-1,\ w_a<0\) at dataset-dependent significance, while current curvature
constraints remain consistent with flat geometry \cite{DESI2025DR2BAO,DESY6DynamicalDE2026,Planck2018Parameters,ACTDR6Extended2025}.  This is a
sharp sign-level tension for the positive-coefficient boundary valuation class.
It is not a direct exclusion of the framework, because the exact VDE background
and the chosen perturbative lift must be fitted directly.  Negative
coefficients, released five-slot spectra, or different declared lifts remain
separate null-test hypotheses.
\end{criterion}

No additional identity is required for the finite-module null tests developed
here.  This statement is restricted to the declared valuation module and does
not assert completeness of microscopic source theory, compact transfer, or
likelihood inference.  The remaining work is either numerical---closed-SVT
Boltzmann transfer, initial covariance, lensing map, and likelihood---or
microscopic, namely a genuine valuation-to-source action bridge and an
amplitude-selection mechanism.




\subsection{Uniqueness envelope of the homogeneous invariant structures}
\label{sec:uniqueness-envelope}

The word ``unique'' is used here only inside a declared invariant problem.  No
claim is made that all dark-sector theories, nonlocal functionals, effective
field theories, geometric completions, or source models reduce to the
structures below.  The statement is instead conditional: once the round
homothetic Euclidean valuation problem and its dilation action are fixed, the
homogeneous invariant normal form is forced.

\begin{definition}[Declared homogeneous invariant problem]
The declared homogeneous invariant problem consists of
\begin{enumerate}
\item the round homothetic family \(\Bfour_A\) with boundary \(\Sthree_A\);
\item the scale action generated by \(\Dil=A\dd/\dd A\);
\item continuous, Euclidean-invariant, additive valuations on convex bodies in
\(\mathbb R^4\), restricted to the round homothetic family;
\item homogeneous density profiles obtained from the work map
\(\Efour(A)\mapsto\epsD(A)=\Efour(A)/V_3(A)\);
\item linear and dimensionless invariants constructed from the resulting finite
\(\Dil\)-module, without adding new fields, nonlocal spectral functionals,
support-function boundary data, external length scales, microscopic source
dynamics, or initial-covariance data.
\end{enumerate}
All uniqueness statements in this subsection refer only to this declared
problem.
\end{definition}

\begin{theorem}[Conditional completeness and unique invariant normal form]
\label{thm:uniqueness-envelope}
Inside the declared homogeneous invariant problem, the following structures are
forced.
\begin{enumerate}
\item The valuation-density module is
\begin{equation}
\boxed{
\mathcal V_{\rm val}=\operatorname{span}\{A,1,A^{-1},A^{-2},A^{-3}\},
\qquad
S_{\rm val}=\{1,0,-1,-2,-3\}.}
\end{equation}
Every admitted profile has a unique expansion
\(\epsD=\sum_{s\in S_{\rm val}}C_sA^s\).

\item The unique monic minimal polynomial of \(\Dil\) on the full module is
\begin{equation}
\boxed{P_{\rm val}(z)=(z-1)z(z+1)(z+2)(z+3),}
\end{equation}
so the polynomial action factors faithfully through
\begin{equation}
\boxed{
\mathfrak A_{\Dil}\simeq\mathbb R[z]/(P_{\rm val})\simeq\mathbb R^5.}
\end{equation}

\item Because the spectrum of \(\Dil\) is simple, its full commutant on
\(\mathcal V_{\rm val}\) is exactly the polynomial algebra generated by
\(\Dil\):
\begin{equation}
\boxed{
\operatorname{Comm}(\Dil)
=\{T\in\operatorname{End}(\mathcal V_{\rm val}):T\Dil=\Dil T\}
=\mathfrak A_{\Dil}.}
\end{equation}
Hence every linear endomorphism preserving all dilation eigenspaces is uniquely
represented by a polynomial in \(\Dil\) modulo \(P_{\rm val}\).

\item The primitive projectors
\begin{equation}
\boxed{
\mathsf P_s=
\prod_{\substack{t\in S_{\rm val}\\t\neq s}}
\frac{\Dil-t}{s-t}}
\end{equation}
are the unique complete set of primitive mutually orthogonal spectral
projectors, up to relabelling of the five roots.  The coefficients
\(\widehat{\mathcal I}_s=A^{-s}\mathsf P_s\epsD\) are therefore the canonical
linear spectral coordinates.

\item For every \(S\subseteq S_{\rm val}\),
\(\mathcal V_S=\operatorname{span}\{A^s:s\in S\}\) is determined only by its
support set, with
\begin{equation}
\mathcal V_S\cap\mathcal V_T=\mathcal V_{S\cap T},
\qquad
\mathcal V_S+\mathcal V_T=\mathcal V_{S\cup T},
\end{equation}
and
\begin{equation}
\boxed{
\operatorname{Ann}_{\Dil}(\mathcal V_S)
=\left(\prod_{s\in S}(z-s)\right).}
\end{equation}
Thus the support lattice is the subset lattice of the five forced eigenvalues,
while the annihilator lattice is its order-reversing divisor-ideal image.

\item On the nonzero positive cone, the normalized branch composition
\begin{equation}
\boxed{\pi_s=\frac{\mathsf P_s\epsD}{\epsD},\qquad \sum_s\pi_s=1}
\end{equation}
is unique.  Equivalently, the four nontrivial moments
\begin{equation}
\mathfrak m_n=\frac{\Dil^n\epsD}{\epsD}=\sum_s s^n\pi_s,
\qquad n=1,\ldots,4,
\end{equation}
together with \(\mathfrak m_0=1\) determine the same five fractions uniquely,
because the Vandermonde matrix on the five distinct spectral values is
invertible.  The variance, cumulant, Hankel, diameter, and equation-of-state
flow identities are therefore forced consequences of this finite positive
spectral measure.

\item Inside the declared local curvature-measure boundary class, the unique
support face is obtained by setting \(C_1=0\):
\begin{equation}
\boxed{S_\partial=\{0,-1,-2,-3\},\qquad
\Dil(\Dil+1)(\Dil+2)(\Dil+3)\epsD=0.}
\end{equation}
At the level of the background expansion-rate quotient by
\(\operatorname{span}\{1,a^{-2},a^{-3}\}\), the unique remaining class is
\begin{equation}
\boxed{[a^{-1}].}
\end{equation}
This quotient statement does not identify physical \(a^{-2}\) stress with
geometric curvature; the split by \(\Omega_K^{\rm geom}\) remains necessary.
\end{enumerate}
\end{theorem}

\begin{proof}
The Euclidean convex valuation hypothesis restricts the filling energy on the
round homothetic family to intrinsic-volume degrees \(4,3,2,1,0\).  Division by
\(V_3\propto A^3\) gives the five density exponents in item 1, which are
linearly independent on every interval \(A>0\).  The minimal-polynomial and
semisimple-algebra statements then follow from the five distinct eigenvalues of
\(\Dil\).  For a linear operator with simple spectrum, commutation with
\(\Dil\) forces diagonality in the eigenbasis \(\{A^s\}\); Lagrange
interpolation produces the unique polynomial of degree at most four with those
five diagonal values, proving item 3 and the projector statement.  The support
lattice and principal annihilators follow from linear independence of the power
basis.  On the positive cone, the primitive projectors recover each branch
contribution uniquely; normalization gives \(\pi_s\), and the moment map is
invertible because its Vandermonde determinant is nonzero.  Finally, the
boundary-local restriction excludes the bulk intrinsic-volume branch and no
other branch inside the declared boundary class, while quotienting the
remaining background powers by the curved-\(\Lambda\)CDM nuisance span leaves
only \([a^{-1}]\).
\end{proof}

\begin{corollary}[Relative uniqueness of canonical scalar lifts]
For a positive branch \(\rho_s\propto a^s\) with \(s<0\), once the minimally
coupled canonical scalar source class with fixed positive kinetic normalization
is declared, its homogeneous realization is unique up to field sign and
additive field shift:
\begin{equation}
\left(\frac{\dd\phi_s}{\dd\tau}\right)^2=-\frac{s}{3}\rho_s,
\qquad
V_s=\frac{6+s}{6}\rho_s.
\end{equation}
Consequently its mass-flow invariant and compact scalar free kernel are fixed by
the background history.  This is uniqueness inside the declared canonical
source class, not uniqueness among all stress tensors realizing the same
homogeneous scaling.
\end{corollary}

\begin{proposition}[Scope boundary: global uniqueness is not implied]
\label{prop:no-global-uniqueness}
The assumptions of the declared homogeneous invariant problem do not imply
uniqueness after the admissible object class is enlarged.  New invariant data
can arise from, for example, nonlocal spectral functionals, support-function or
ambient-position boundary terms, additional fields, EFT operators and Wilson
coefficients, external length scales, quotient-topology data, or initial
covariance data.  Therefore the maximal defensible uniqueness statement is
\begin{equation}
\boxed{
\text{the listed invariants are uniquely forced inside the declared homogeneous valuation category.}}
\end{equation}
No absolute uniqueness claim over all dark-sector theory space is made.
\end{proposition}
\begin{proof}
Each listed enlargement introduces admissible input data that are absent from
the declared problem and on which new invariants may depend.  The manuscript
already provides explicit examples of this non-entailment: canonical scalar and
ADM-response lifts can share a homogeneous scaling while having inequivalent
perturbative responses.  Hence the valuation-module assumptions cannot imply a
global uniqueness theorem for a larger category that contains such additional
objects.
\end{proof}

\begin{criterion}[Final conditional uniqueness statement]
The complete invariant normal form of the declared homogeneous problem consists
of the five dilation eigenvalues, the semisimple operator algebra and its
commutant, primitive projectors, support/annihilator lattice, normalized positive
spectral fractions and their moment hierarchy, the boundary-local support face,
and the one-dimensional expansion-rate quotient.  Declared source-class kernels
are then fixed only conditional on the chosen lift.  Outside this category,
uniqueness is not claimed.
\end{criterion}

\subsection{Derived branch-to-source and amplitude ledgers}
\label{sec:derived-ledgers}

The preceding theorems determine a finite background alphabet but do not select
a unique microscopic grammar.  It is nevertheless useful to collect the exact
maps that are already derivable once a source class and a normalization ledger
are declared.  These maps are consequences of the existing theorems; they add
no new physical postulate.

\begin{criterion}[Declared branch-to-source lift map]
For the boundary spectrum, a convenient minimal lift assignment is
\begin{equation}
\boxed{
\begin{array}{ccl}
s=0   &\longmapsto& T_{\mu\nu}^{(\Lambda)}=-\rho_\Lambda g_{\mu\nu},\\[1mm]
s=-1  &\longmapsto& \text{canonical scalar }\phi_X,\\[1mm]
s=-2  &\longmapsto& \text{canonical scalar }\phi_2,\\[1mm]
s=-3  &\longmapsto& \text{conserved pressureless current for CDM-like transfer}.
\end{array}}
\end{equation}
This is a declared lift map, not a derivation from \(\Efour^{\rm val}\), and it
is not a functor unless an explicit source category and morphism map are added.
A canonical scalar also exists as a homogeneous realization for any positive
branch with \(s<0\), including \(s=-3\), but at \(s=-3\) it has
\(c_{s,{\rm rest}}^2=1\) and therefore is not the CDM-like lift adopted here.
\end{criterion}

\begin{corollary}[Canonical branch reconstruction and general compact scalar kernel]
\label{cor:general-canonical-branch-kernel}
Let \(s<0\), \(\Omega_s>0\), and
\begin{equation}
\rho_s(a)=3M_{\rm Pl}^2H_0^2\Omega_s a^s,
\qquad
F(a):=\frac{H^2(a)}{H_0^2}>0.
\end{equation}
Inside the minimally coupled canonical scalar class,
\begin{align}
\left(\frac{\dd\phi_s}{\dd\tau}\right)^2
&=-sM_{\rm Pl}^2H_0^2\Omega_s a^s,\\
V_s(a)&=\frac{6+s}{2}M_{\rm Pl}^2H_0^2\Omega_s a^s,\\
\Dil\phi_s
&=\pm M_{\rm Pl}\sqrt{\frac{-s\,\Omega_s a^s}{F(a)}},\\
V_{s,\phi}
&=\mp\frac{6+s}{2}M_{\rm Pl}H_0^2
\sqrt{-s\,\Omega_s a^sF(a)},\\
\frac{V_{s,\phi\phi}}{H_0^2}
&=-\frac{6+s}{4}\left[sF(a)+\Dil F(a)\right].
\end{align}
Equivalently,
\begin{equation}
\boxed{
\frac{V_{s,\phi\phi}}{H_0^2}
=-\frac{6+s}{4}\left[sF(a)+\Dil F(a)\right].}
\end{equation}
Hence the dimensionless free scalar kernel on the round closed representative is
\begin{equation}
\boxed{
\mathscr W_{s,n}(a)=
(-\Omega_K^{\rm geom})a^{-2}n(n+2)
-\frac{6+s}{4}\left[sF(a)+\Dil F(a)\right].}
\end{equation}
For \(s=-1\) this reduces to
\(\frac54[F-\Dil F]\); for \(s=-2\) it reduces to
\begin{equation}
\boxed{
\frac{m_2^2}{H_0^2}=2F-\Dil F.}
\end{equation}
The last expression is the correct \(s=-2\) specialization; a formula
\(F-\tfrac12\Dil F\) would differ by an erroneous factor of two.
This corollary closes the declared scalar source equation for each branch; it
does not supply the full metric-coupled compact SVT hierarchy.
\end{corollary}

\begin{proof}
The first two relations are the canonical identities
\(\dot\phi_s^2=(1+w_s)\rho_s\) and
\(V_s=(1-w_s)\rho_s/2\) with \(w_s=-1-s/3\).  Dividing the field velocity by
\(H\) gives the logarithmic trajectory derivative.  Differentiating \(V_s(a)\)
through this monotone trajectory gives \(V_{s,\phi}\) and then the mass-flow
invariant already proved above.  Adding the closed-\(\Sthree\) Laplacian
eigenvalue gives the displayed kernel.
\end{proof}

\begin{proposition}[Valuation-coefficient to cosmological-amplitude ledger]
\label{prop:valuation-amplitude-ledger}
Write
\begin{equation}
\epsD(A)=\frac{\gamma A}{4}+\alpha_0+
\frac{3\alpha_1}{A}+\frac{3\alpha_2}{A^2}+\frac{\alpha_3}{A^3},
\qquad A=A_0a,
\end{equation}
and define bookkeeping density parameters by
\begin{equation}
\rho_{s0}:=3M_{\rm Pl}^2H_0^2\Omega_s^{\rm prof}
\end{equation}
for the five profile slots before background quotienting.  Then
\begin{align}
\rho_{1,0}&=\frac{\gamma A_0}{4}, &
\rho_{0,0}&=\alpha_0, &
\rho_{-1,0}&=\frac{3\alpha_1}{A_0},\\
\rho_{-2,0}&=\frac{3\alpha_2}{A_0^2}, &
\rho_{-3,0}&=\frac{\alpha_3}{A_0^3},
\end{align}
and therefore
\begin{align}
\gamma&=\frac{12M_{\rm Pl}^2H_0^2}{A_0}\Omega_1^{\rm prof},\\
\alpha_0&=3M_{\rm Pl}^2H_0^2\Omega_0^{\rm prof},\\
\alpha_1&=M_{\rm Pl}^2H_0^2A_0\Omega_{-1}^{\rm prof},\\
\alpha_2&=M_{\rm Pl}^2H_0^2A_0^2\Omega_{-2}^{\rm prof},\\
\alpha_3&=3M_{\rm Pl}^2H_0^2A_0^3\Omega_{-3}^{\rm prof}.
\end{align}
For a closed round geometry, \(\Omega_K^{\rm geom}<0\) and
\begin{equation}
A_0=\frac{c}{H_0\sqrt{-\Omega_K^{\rm geom}}}.
\end{equation}
Thus the ledger converts valuation coefficients into cosmological amplitudes
after \(H_0\), \(A_0\), and the normalized branch amplitudes are supplied.  It
is an exact matching dictionary, not an amplitude-selection theorem.  After the
boundary filter and background quotient one identifies
\(\Omega_{-1}^{\rm prof}=\Omega_X\) and
\(\Omega_{-2}^{\rm prof}=\Omega_{2,{\rm val}}\), while the constant and dustlike
slots can be absorbed into the corresponding effective vacuum and matter
parameters.
\end{proposition}

\begin{proposition}[Endpoint charge-current matching ledger]
\label{prop:endpoint-matching-ledger}
For
\begin{equation}
E_{\rm end}=\sigma V_3+\mu\QU,
\qquad V_3(A_0)=2\pi^2A_0^3,
\end{equation}
the present-day vacuum and dustlike densities satisfy
\begin{equation}
\boxed{\sigma=3M_{\rm Pl}^2H_0^2\Omega_\Lambda,}
\qquad
\boxed{\mu\QU=6\pi^2M_{\rm Pl}^2H_0^2A_0^3\Omega_w.}
\end{equation}
For closed geometry this may be rewritten as
\begin{equation}
\boxed{
\mu\QU=
6\pi^2M_{\rm Pl}^2H_0^2
\left(\frac{c}{H_0\sqrt{-\Omega_K^{\rm geom}}}\right)^3\Omega_w.}
\end{equation}
The equality
\(\mathcal E_{\rm rest}\overset{\rm match}{=}\mu\QU\) is therefore a conserved
matching equation once imposed; it does not determine \(\mu\), \(\QU\), or
\(\Omega_w\) from topology alone.
\end{proposition}

\begin{criterion}[Physical boundary parameter ledger]
The background model is
\begin{equation}
\boxed{
F_\partial(a)=
\Omega_ra^{-4}+\Omega_ma^{-3}
+\bigl(\Omega_K^{\rm geom}+\Omega_{2,{\rm val}}\bigr)a^{-2}
+\Omega_Xa^{-1}+\Omega_\Lambda,}
\end{equation}
with present-day normalization
\begin{equation}
\Omega_r+\Omega_m+\Omega_K^{\rm geom}+\Omega_{2,{\rm val}}
+\Omega_X+\Omega_\Lambda=1.
\end{equation}
The two \(a^{-2}\) terms are rank-degenerate in \(F_\partial\) alone and are
split only after independent transverse-distance information determines
\(\Omega_K^{\rm geom}\).  The one-parameter canonical VDE subclass imposes
\(\Omega_{2,{\rm val}}=0\) and \(\Omega_X\ge0\).  The statistical machinery
already developed above---GLS projection, one-sided boundary likelihood ratios,
and chi-bar-square cone calibration---tests these declared parameter spaces; it
does not itself supply the cosmological likelihood.
\end{criterion}

\begin{remark}[Derived interface status]
The chain established by the present manuscript is therefore
\begin{equation}
\boxed{
\begin{gathered}
\text{Euclidean valuation closure}
\Longrightarrow \text{ finite homogeneous spectrum}\\
\Longrightarrow \text{ boundary/background quotient}\\
\Longrightarrow \text{ declared source lift}\\
\Longrightarrow \text{ branchwise scalar kernel or current response}\\
\Longrightarrow \text{ curvature-split and likelihood interface}.
\end{gathered}}
\end{equation}
Each arrow is to be read at the scope stated in its theorem or criterion.  In
particular, the chain is not a derivation of the source action from the
valuation functional and not an observational preference statement.
\end{remark}

\subsection{Geometric two-parameter extension and nested one-parameter VDE comparison}
The endpoint vacuum--dust completion discussed above is exactly background-degenerate with curved \(\Lambda\)CDM. The general physical local curvature-measure boundary sector is different: after geometric curvature is separated from a physical \(w=-1/3\) valuation component, its background geometry plus expansion contains
\begin{equation}
\mathcal H^2_{\partial}(a)=
\Omega_r a^{-4}+\Omega_m a^{-3}
+\bigl(\Omega_K^{\rm geom}+\Omega_{2,\rm val}\bigr)a^{-2}
+\Omega_\Lambda+\Omega_Xa^{-1}.
\end{equation}
Relative to curved \(\Lambda\)CDM the general physical boundary extension therefore has two new coordinates, \(\Omega_X\) and \(\Omega_{2,\rm val}\). The deliberately nested one-parameter VDE test imposes \(\Omega_{2,\rm val}=0\), leaving
\begin{equation}
\mathcal H^2_{\rm VDE}(a)=
\Omega_r a^{-4}+\Omega_m a^{-3}+\Omega_K^{\rm geom}a^{-2}
+\Omega_\Lambda+\Omega_Xa^{-1}.
\end{equation}
For that nested subclass the direct model-comparison question is one-dimensional:
\begin{equation}
H_0:\Omega_X=0
\qquad\text{versus}\qquad
H_1:\Omega_X\ge0,
\end{equation}
with detection corresponding to a profiled maximum at \(\widehat\Omega_X>0\), and with all common cosmological and nuisance parameters treated identically. A test of the full physical boundary sector must instead fit \((\Omega_X,\Omega_{2,\rm val})\) jointly with \(\Omega_K^{\rm geom}\).

The differential quotient invariant gives an independent background diagnostic,
\begin{equation}
\widehat\Omega_X(a)
=-\frac a6\,\Dil(\Dil+2)(\Dil+3)(\Dil+4)\mathcal H^2(a).
\end{equation}
A valid VDE background requires \(\widehat\Omega_X(a)\) to be constant. Curved \(\Lambda\)CDM is the sharper null \(\widehat\Omega_X(a)=0\).

Current data do not justify treating a closed \(\Sthree\) VDE cosmology as established. Planck 2018 combined with BAO found curvature consistent with flatness, and ACT DR6 combined CMB+lensing+BAO constraints likewise remain fully consistent with flat geometry \cite{Planck2018Parameters,ACTDR6Extended2025}. The DESI DR2 BAO analysis is well described by flat \(\Lambda\)CDM and, in combinations with CMB and supernova data, reports preferred evolving-dark-energy solutions in the quadrant \(w_0>-1,\ w_a<0\) \cite{DESI2025DR2BAO}. The positive minimal VDE tangent instead obeys \(w_a\ge0\). This is a useful sign-level diagnostic but not a substitute for fitting the exact VDE background and perturbation model directly.

A credible VDE claim therefore requires the one-parameter improvement in likelihood or Bayesian evidence to survive the appropriate boundary-null calibration and robustness tests. For the nested positive model the primary frequentist summary should report the likelihood-ratio statistic \(q_X\), its boundary-calibrated one-sided tail probability, and profile-likelihood constraints on \(\Omega_X\). A Bayesian analysis should report the evidence ratio with explicit priors. AIC and BIC may be supplied as supplementary diagnostics rather than as simultaneous necessary conditions.

\begin{criterion}[Two-level evidential logic for VDE]
Two claims must be distinguished.
\begin{enumerate}
\item A \emph{phenomenological VDE detection} requires evidence for the nested \(\Omega_X\) direction under the declared perturbative closure, with the one-parameter gain surviving the statistical penalty and data/likelihood robustness tests.
\item A \emph{compact-valuation origin claim} requires, in addition, a geometrically consistent closed branch \(\Omega_K^{\rm geom}<0\), the curvature split separating \(\Omega_{2,\rm val}\) from \(\Omega_K^{\rm geom}\), compatibility with the valuation closure, and compact-\(\Sthree\) spectral-geometric information at the inferred curvature scale.
\end{enumerate}
Thus a nonzero best-fit \(\Omega_X\) would support the phenomenological quotient branch first; it would not by itself prove the auxiliary filling interpretation or the compact-topological origin.
\end{criterion}

\subsection{Background integrability}

Let
\begin{equation}
\mathcal H^2(z):=\frac{H^2(z)}{H_0^2}.
\end{equation}
For fixed visible sector, fixed \(\Omega_{r0}\), fixed \(\Omega_{b0}\), fixed geometric \(\Omega_K^{\rm geom}\), and a separately conserved dark-energy branch with profile \(f_q(z)\), write
\begin{equation}
\mathcal H^2(z)=
\Omega_{r0}(1+z)^4
+\Omega_{b0}(1+z)^3
+\Omega_K^{\rm geom}(1+z)^2
+\Omega_{q0}f_q(z)
+\Omega_{w0}f_w(z).
\end{equation}
The reconstructed amplitude of the \(w\)-branch is
\begin{equation}
\widehat{\Omega}_{w0}(z)
=
\frac{
\mathcal H^2(z)
-\Omega_{r0}(1+z)^4
-\Omega_{b0}(1+z)^3
-\Omega_K^{\rm geom}(1+z)^2
-\Omega_{q0}f_q(z)
}{
f_w(z)
}.
\end{equation}

\begin{criterion}[Background integrability]
A single FRW background with one separately integrable \(w\)-branch requires
\begin{equation}
\frac{\dd}{\dd z}\widehat{\Omega}_{w0}(z)=0
\end{equation}
on the redshift interval on which the branch is defined and \(f_w(z)\neq0\). In a statistical implementation, the reconstructed values \(\widehat{\Omega}_{w0}(z_i)\) must be consistent with a single constant \(\Omega_{w0}\) under the full covariance matrix of the data.
\end{criterion}

For the endpoint dust branch,
\begin{equation}
f_w(z)=(1+z)^3.
\end{equation}
Failure of the constancy condition means that the proposed component is not one separately integrable dust-like branch of the assumed FRW background.

\subsection{Independent spectral-geometric invariant}

For a compact simply connected spatial section \(\Sthree\) with curvature radius \(R_c\), the standard scalar harmonic spectrum is \cite{HalliwellHawking1985,KodamaSasaki1984}
\begin{equation}
\operatorname{Spec}(-\nabla^2_{\Sthree})
=
\left\{
\frac{n(n+2)}{R_c^2}
\right\}_{n=0}^{\infty},
\qquad
 d_n=(n+1)^2,
\end{equation}
where \(d_n\) is the scalar-mode degeneracy.
The same curvature radius gives the dimensionless spectral-geometric invariant
\begin{equation}
\boxed{
\frac{c^2\lambda_n}{H_0^2}
=(-\Omega_K^{\rm geom})\,n(n+2),
}
\end{equation}
which directly links the discrete compact harmonic spectrum to the geometric curvature parameter inferred from distances. With the closed-sign convention \(\Omega_K^{\rm geom}<0\),
\begin{equation}
R_c=\frac{c}{H_0\sqrt{-\Omega_K^{\rm geom}}}.
\end{equation}

A smooth \(w(z)\)-parameterization may change the inferred distance-redshift relation, growth history, or preferred curvature radius. It does not replace the discrete compact-\(\Sthree\) eigenvalue structure by the continuous flat-space spectrum.

\begin{criterion}[Spectral-geometric independence]
A compact-\(\Sthree\) branch is independently supported only if the inferred curvature scale and perturbation analysis are compatible with the discrete scalar, vector, and tensor harmonic spectra of \(\Sthree\). In the simply connected case this concerns curvature scale, mode discreteness, and large-angle correlations. Matched circles or multiple imaging require an additional nontrivial quotient \(\Sthree/\Gamma\) and are not consequences of simply connected \(\Sthree\) alone.
\end{criterion}

Cosmic-topology searches and matched-circle tests belong to this quotient-space sector rather than to simply connected \(\Sthree\) alone \cite{CornishSpergelStarkman1998,LuminetEtAl2003,Planck2016Topology,CosmicTopologyReview2026}. Current topology reviews find no convincing detection of nontrivial cosmic topology; existing searches exclude only parts of the topology, scale, and observer-position parameter space \cite{CosmicTopologyReview2026}. For spherical quotients such as lens spaces \(L(p,q)=\Sthree/\mathbb Z_p\), matched-circle non-detections place curvature-scale-dependent limits, for example approximately \(p\le9\) at \(|\Omega_K|\simeq0.05\) and up to \(p\simeq24\) at \(|\Omega_K|\simeq0.02\) \cite{COMPACTLensSpaces2025}. These quotient-space constraints are not direct constraints on simply connected \(\Sthree\).

Absence of a detected spectral imprint does not by itself disprove compactness if the curvature radius is too large for the data to resolve. It does, however, remove independent confirmation of the compact-\(\Sthree\) branch unless the completed model predicts no observable spectral effect in the tested range.

\begin{remark}[What the compact spectrum can and cannot distinguish]
The compact \(\Sthree\) spectrum is a geometric and topological discriminator. It can distinguish a compact curved spatial section from a flat or non-compact parametrization, provided the curvature scale is observationally accessible. It does not by itself distinguish the endpoint matter law from a closed curved \(\Lambda\)CDM model placed on the same \(\Sthree\) geometry. Endpoint-sector distinguishability requires additional information: amplitude priors, perturbative observables, spectral weights, allowed initial spectra, non-minimal valuation branches, or penalized statistical preference against competitors with the same spatial topology included.
\end{remark}

\subsection{Likelihood tests for nested VDE branches}
The one-parameter positive VDE submodel is naturally tested as a nested boundary problem,
\begin{equation}
H_0:\ \Omega_X=0,
\qquad
H_1:\ \Omega_X\ge0.
\end{equation}
Define the likelihood-ratio statistic
\begin{equation}
q_X:=2\left[
\ln L(\widehat\Omega_X)-\ln L(0)
\right],
\qquad
\widehat\Omega_X\ge0.
\end{equation}

\begin{theorem}[Exact Gaussian one-sided VDE likelihood-ratio law]
Consider
\begin{equation}
y=S\beta+x\Omega_X+\eta,
\qquad
\eta\sim\mathcal N(0,C),
\qquad
\Omega_X\ge0,
\end{equation}
with known \(C\), full-rank nuisance matrix \(S\), and no nuisance parameter that becomes unidentified at \(\Omega_X=0\). Let
\begin{equation}
W=C^{-1},
\qquad
M_S=I-S(S^TWS)^{-1}S^TW,
\qquad
I_X=x^TWM_Sx.
\end{equation}
The unconstrained GLS coordinate is
\begin{equation}
\widetilde\Omega_X=\frac{x^TWM_Sy}{I_X},
\qquad
\operatorname{Var}(\widetilde\Omega_X)=I_X^{-1},
\end{equation}
and the constrained maximum-likelihood estimate is
\begin{equation}
\boxed{\widehat\Omega_X=\max(0,\widetilde\Omega_X).}
\end{equation}
For the null \(H_0:\Omega_X=0\), the likelihood-ratio statistic
\begin{equation}
q_X:=2\left[\ln L(\widehat\Omega_X)-\ln L(0)\right]
\end{equation}
obeys exactly in this Gaussian linear problem
\begin{equation}
\boxed{q_X\sim\frac12\delta_0+\frac12\chi_1^2.}
\end{equation}
For \(q_X>0\),
\begin{equation}
\boxed{p_{\rm one\text{-}sided}=1-\Phi_{\rm N}(\sqrt{q_X}),}
\end{equation}
where \(\Phi_{\rm N}\) is the standard-normal cumulative distribution function.
\end{theorem}
\begin{proof}
Under \(H_0\),
\begin{equation}
Z:=\frac{x^TWM_Sy}{\sqrt{I_X}}\sim\mathcal N(0,1).
\end{equation}
Profiling the Gaussian likelihood gives \(q_X=[\max(0,Z)]^2\). Half of the Gaussian probability lies at \(Z\le0\), producing the atom at zero; conditional on \(Z>0\), \(Z^2\) gives the positive half of the \(\chi_1^2\) law. This is the exact Gaussian realization of the standard boundary-null structure \cite{Chernoff1954,SelfLiang1987}.
\end{proof}

\begin{remark}[Two positive boundary amplitudes]
For the jointly constrained Gaussian linearized model
\begin{equation}
\Omega_X\ge0,
\qquad
\Omega_{2,\rm val}\ge0,
\end{equation}
the null lies at the corner of a two-dimensional cone. If \(r\) is the correlation of the two unconstrained, nuisance-profiled amplitude estimators, the whitened cone has opening angle
\begin{equation}
\alpha=\arccos(-r),
\end{equation}
with chi-bar-square weights
\begin{equation}
w_2=\frac{\alpha}{2\pi},
\qquad
w_0=\frac{\pi-\alpha}{2\pi},
\qquad
w_1=\frac12.
\end{equation}
For a full nonlinear Boltzmann likelihood, simulation-based calibration is preferable.
\end{remark}

\begin{criterion}[Model-comparison hierarchy]
For the nested positive one-parameter VDE test, the primary frequentist calibration is the boundary-aware likelihood-ratio law above in the Gaussian limit, with asymptotic boundary theory or simulation calibration used for the full likelihood. A Bayesian analysis should report the evidence or Bayes factor together with the stated prior on \(\Omega_X\). AIC and BIC may be reported as supplementary information criteria, but simultaneous inequalities involving AIC, BIC, and Bayesian evidence are not treated as a theorem-level necessary condition for model superiority.
\end{criterion}

Planck 2018 combined with BAO found spatial curvature consistent with flatness, \(\Omega_K^{\rm geom}=0.001\pm0.002\), so that combined constraint provides no evidence for nonzero closed curvature \cite{Planck2018Parameters}. ACT DR6, combined with Planck primary CMB, CMB lensing, and DESI Year-1 BAO in its geometry analysis, finds \(\Omega_k=0.0019\pm0.0015\) and full consistency with flatness; conditional on a closed branch it reports \(R_k>105\,\mathrm{Gpc}\) at 95\% confidence \cite{ACTDR6Extended2025}. A more recent compressed full-CMB combination using SPT-3G D1, ACT DR6, Planck primary CMB, and Planck NPIPE lensing, when combined with DESI DR2 in \(o\Lambda\)CDM, gives \(\Omega_k=0.0026\pm0.0011\), i.e. a simple-model posterior shifted toward the open side rather than toward closed \(\Sthree\) \cite{CMBComp2026}. This is not a model-independent exclusion of closed geometry: when evolving dark energy and curvature are released simultaneously, Planck+ACT DR6+DESI DR2+SNe analyses find much broader curvature constraints, for example \(10^3\Omega_k=0.9\pm1.4\) in a \(w_0w_a+\Omega_k\) setting \cite{PulidoCervantes2026}. Thus closed \(\Sthree\) is disfavored in simple \(o\Lambda\)CDM fits but is not model-independently excluded.

DESI DR2 uses more than 14 million galaxies and quasars from three years of observations; its BAO distances are well described by flat \(\Lambda\)CDM, while combinations with CMB and supernovae prefer a CPL region with \(w_0>-1\), \(w_a<0\) at dataset-dependent significance \cite{DESI2025DR2BAO}. The extended DESI DR2 dark-energy analysis using parametric and nonparametric reconstructions reports a robust preference for models featuring phantom crossing; non-crossing alternatives are disfavored but not yet ruled out \cite{DESIExtendedDE2025}. The full six-year Dark Energy Survey multi-probe analysis finds, with DESI DR2 BAO, \(w_0=-0.84^{+0.06}_{-0.07}\), \(w_a=-0.53^{+0.33}_{-0.28}\), and with primary CMB information \(w_0=-0.82\pm0.05\), \(w_a=-0.63^{+0.21}_{-0.18}\); the reported departure from a cosmological constant varies from about \(2.3\sigma\) to \(3.2\sigma\) under probe removal, with the all-data combination at \(3.0\sigma\), and is described as a weak preference for evolving dark energy \cite{DESY6DynamicalDE2026}. DES Y6 3\(\times\)2pt also gives strong growth constraints and, in a joint CMB plus low-redshift constant-\(w\) fit, returns \(S_8=0.806^{+0.006}_{-0.007}\), \(\Omega_m=0.302\pm0.003\), and \(w=-0.981^{+0.021}_{-0.022}\) \cite{DESY6ThreeByTwo2026}. ACT DR6 finds a moderate \(w_0w_a\) hint in combinations containing DESI BAO and supernova data, while its overall extended-model assessment finds no model significantly preferred over baseline \(\Lambda\)CDM; the late-time hint is largely driven by the external BAO and supernova information \cite{ACTDR6Extended2025}. Euclid's public timeline places DR1 in 2026/2027, and the Euclid Consortium states that the first main cosmology results are expected in 2027; no final Euclid dark-energy likelihood is therefore used here \cite{EuclidTimeline2026,EuclidCosmology2026}.

The positive valuation theorem predicts \(\Dil w\le0\), equivalently \(w_a\ge0\) for the local CPL tangent, which is opposite to the currently preferred CPL direction. The sign comparison is meaningful only after the visible-matter, geometric-curvature, and residual-sector decompositions used in the observational reconstruction are matched to those of the valuation model. It is therefore a model-specific null test, not a parameterization-independent statement about arbitrary reconstructed dark energy.

\subsection{Separated evidential criteria}

\begin{criterion}[Phenomenological VDE preference]
A phenomenological VDE preference is established only if the declared VDE source realization yields a statistically nonzero \(\Omega_X\) (or a nonzero two-parameter boundary direction in the general model), the improvement survives the appropriate parameter penalty and boundary-null calibration, and the conclusion is robust to reasonable changes of data combination, nuisance treatment, and likelihood implementation. Independent detection of compact spatial topology is not required for this phenomenological claim.
\end{criterion}

\begin{criterion}[Compact-valuation origin support]
A stronger compact-valuation origin claim requires all phenomenological conditions above and, in addition:
\begin{enumerate}
\item a geometrically consistent closed FLRW branch with \(\Omega_K^{\rm geom}<0\), inferred jointly from expansion and transverse-distance information rather than from the \(a^{-2}\) coefficient of \(\mathcal H^2\) alone;
\item a consistent curvature split
\begin{equation}
\Omega_{2,\rm val}=\mathcal C_{-2}[\mathcal H^2]-\Omega_K^{\rm geom};
\end{equation}
\item compatibility with the local curvature-measure valuation closure and, where observationally accessible, the discrete scalar/vector/tensor harmonic structure of compact \(\Sthree\).
\end{enumerate}
Failure of the compact-geometric conditions does not by itself falsify a phenomenological \(w=-2/3\) component; it falsifies or weakens the proposed compact-valuation origin of that component.
\end{criterion}

\subsection{Falsification logic}

The construction is falsified or restricted at the corresponding level if at least one of the following holds:
\begin{enumerate}
\item the endpoint-selected or independently selected \(\Efour(A,\xi)\) fails to reproduce the reconstructed residual in the matching condition;
\item the internal equations are incompatible with the exchange law;
\item the dust slot necessarily has unacceptable \(c_s^2\), \(\pi_{\mu\nu}\), or perturbation growth;
\item the compact radius required by an independently completed model is excluded by curvature, spectral, or topology tests;
\item the gluing conditions of the compatible boundary-data condition are incompatible with the proposed \(\Sfour\) functional;
\item no covariant completion exists, or the completion violates the stated symmetries;
\item the reconstructed branch amplitude fails the integrability condition;
\item the completed compact branch predicts a resolvable compact-spectral or topological signature in the tested regime and that predicted signature is excluded by the data;
\item the completed branch fails the declared boundary-calibrated likelihood-ratio test and/or the independently specified Bayesian model comparison against the nearest competitors.
\end{enumerate}
Absence of an observable compact signature when the curvature radius is too large for the tested regime removes independent support but is not itself a falsification.

Thus background integrability, spectral-geometric discreteness, and robust statistical performance are necessary filters for turning a geometric branch into an observationally superior physical cosmology.



\section{Supplementary Closed-\texorpdfstring{$\Sthree$}{S3} Transfer Implementation Status}
\label{sec:numerical-transfer-package}

A supplementary numerical package is supplied for the \emph{nested canonical
VDE lift} only.  The implementation does not identify the valuation functional
with a Boltzmann-fluid action.  It uses the independently declared canonical
source class at the level of its linear stress-energy data,
\begin{equation}
 w_X=-\frac23,
 \qquad
 c_{s,{\rm rest}}^2=1,
 \qquad
 \pi_{\mu\nu}=0,
\end{equation}
and maps that source to the true-fluid perturbation backend of a pinned
non-flat CLASS revision.  In the supplied parameter ledger the map is
\begin{equation}
\boxed{
\begin{gathered}
\Omega_X\mapsto\texttt{Omega\_fld},\qquad
w_X\mapsto\texttt{w0\_fld}=-\frac23,\\
\texttt{wa\_fld}=0,\qquad
\texttt{cs2\_fld}=1,\qquad
\texttt{use\_ppf}=\texttt{no}.
\end{gathered}}
\end{equation}
The geometric curvature remains a separate input through
\(\texttt{Omega\_k}=\Omega_K^{\rm geom}<0\) for the closed model.  The
cosmological-constant density is left to the CLASS background budget after
\(\Omega_X\) and the remaining components have been declared.

For the harmonic convention used in this paper,
\begin{equation}
 -\nabla^2Y_n=n(n+2)K\,Y_n,
 \qquad K=A_0^{-2}>0,
\end{equation}
the CLASS closed-scalar variable satisfies
\begin{equation}
 q_n=(n+1)\sqrt K,
 \qquad
 k_n^2=q_n^2-K=n(n+2)K.
\end{equation}
Thus the first physical scalar transfer mode is \(n=2\), corresponding to
\(q=3\sqrt K\) and \(k^2=8K\); the \(n=0\) and \(n=1\) sectors retain the
homogeneous and scalar gauge/constraint qualifications stated earlier.

The package contains parameter files for the nested canonical VDE model and a
curved-\(\Lambda\)CDM control, a Python runner requesting background,
perturbation-transfer, raw-CMB, lensed-CMB, lensing, and matter-power outputs,
a discrete \(\Sthree\) mode ledger, a VDE-versus-control comparison script,
unit tests for the analytic invariants and mode map, and a pinned-upstream
installer.  The pinned CLASS revision used for the interface audit is
\texttt{e85808324f51fc694d12e3ed7439552a3c3f9540}.

\begin{criterion}[Numerical implementation status]
\label{crit:numerical-implementation-status}
The existence of this package upgrades the nested canonical VDE transfer layer
from ``not supplied'' to ``implementation supplied, execution validation
pending.''  The analytic/unit tests supplied with the package pass in the
present build environment.  The integration smoke test requiring the
\texttt{classy} module is included but was not executed here because that
module is absent and the present container cannot fetch the pinned upstream
repository.  Consequently no CMB spectrum, matter spectrum, transfer table,
likelihood value, or numerical model preference is claimed in this paper.
The two-parameter \((\Omega_X,\Omega_{2,{\rm val}})\) source implementation,
ADM-response prototypes, Bianchi-IX transfer, blind validation, and exact
likelihood remain outside the supplied numerical closure.
\end{criterion}

This numerical status preserves the hierarchy
\begin{equation}
\boxed{
\text{valuation closure}\neq\text{declared source lift}
\neq\text{validated transfer output}\neq\text{likelihood evidence}.}
\end{equation}

\section{Remaining Numerical and Model-Completion Program}
\label{sec:remaining-program}

The remaining programme now begins after the supplied nested-canonical-VDE
transfer interface.  The homogeneous dilation classification is complete under
its stated assumptions, and a reproducible interface to the pinned non-flat
CLASS backend is supplied for the one-parameter canonical VDE lift.  What
remains open is the covariant valuation-to-source bridge, validated execution
of the supplied transfer package, the general two-parameter boundary-source
implementation, the Bianchi-IX/ADM numerical extensions, and likelihood
analysis:
\begin{enumerate}
\item execute and validate the supplied nested canonical VDE transfer package against its curved-\(\Lambda\)CDM control, including the \(\Omega_X\to0\) null limit, closed-mode normalization, gauge checks, source-column checks, and convergence/precision scans; then implement the general two-parameter \((\Omega_X,\Omega_{2,\rm val})\) source realization and the Bianchi IX fixed-shape coupling diagnostic;
\item fit the scalar completions and geometric curvature split jointly to BAO, supernova, CMB, growth, and lensing likelihoods, using the full FLRW distance relation rather than differentiated distance data;
\item test the annihilator/GLS/Prony closure diagnostics and the positivity constraints with full observational covariance;
\item treat \(\sigma\) and \(\mu\) as dimensional matching/coupling coefficients; if a microscopic theory or EFT is proposed, match them to its independent dimensional scales, operator basis, or renormalized couplings, since topology-only determination is excluded by the scale-covariance no-go theorem;
\item test any stronger compact-origin claim against the discrete \(\Sthree\) harmonic spectrum and topology observables at the inferred \(\Omega_K^{\rm geom}\).
\end{enumerate}

\section{Claim Boundary and Layer Status}
\label{sec:claim-boundary-status}

The manuscript separates theorem closure from model assumptions and from objects
that have not yet been supplied numerically.  The strongest supported reading is
\begin{equation}
\boxed{
\begin{gathered}
\text{finite homogeneous residual-profile classification}\\
+\ \text{declared source-existence lifts}\\
+\ \text{null-test protocols}
\end{gathered}}
\end{equation}
The Euclidean convex valuation hypothesis, the filling-relative boundary filter,
positivity of physical branch coefficients, the restricted endpoint class, and
the initial charge-current matching are declared assumptions or model choices;
they are not promoted to consequences of topology alone.

\begin{center}
\small
\begin{tabular}{@{}p{0.12\linewidth}p{0.28\linewidth}p{0.50\linewidth}@{}}
\toprule
Layer & Object & Scoped status \\
\midrule
R & Residual reconstruction & Closed for the homogeneous round-FLRW work map and inverse representation; not a microscopic source inverse. \\
H & Homogeneous valuation closure & Closed under the declared Euclidean convex valuation hypothesis; the hypothesis itself is a physical/model assumption. \\
S & Source realization & Closed as existence realizations inside the declared canonical-scalar/current source classes; a valuation-derived covariant source action is not supplied. \\
A & Amplitude layer & Coefficient and charge ledgers are closed matching identities; no nonzero amplitude selector is supplied, and dimensional topology-only selection is excluded by scale covariance. \\
T & Transfer evolution & A pinned non-flat CLASS implementation package is supplied for the nested canonical VDE source lift; its integration smoke test and numerical spectra are not executed in the present environment. The general two-parameter, ADM, Bianchi-IX, and likelihood transfer closures are not supplied. \\
C & Consistency/covariance & Curvature splitting, harmonic conventions, and conservation distinctions are supplied at analytic/protocol level; a complete compact covariance/observer implementation is not supplied. \\
L & Likelihood & GLS, one-sided boundary-null, and chi-bar-square protocols are supplied under their statistical assumptions; no exact cosmological likelihood run is supplied. \\
B & Blind distinguishability & Not supplied: no blind injections or synthetic recovery suite is claimed. \\
M & Reproducibility & The analytic package and a version-pinned numerical interface package are supplied with source files, parameter fixtures, tests, build/install instructions, status metadata, and checksums. Executed transfer outputs, likelihood chains, and numerical-validation evidence are not supplied. \\
X & External empirical container & Literature comparisons are contextual benchmarks; no compact-VDE detection or model preference is established by this manuscript. \\
\bottomrule
\end{tabular}
\end{center}

Accordingly, ``complete'' and ``unique'' in this paper mean complete or unique
\emph{inside the declared homogeneous valuation problem}.  They do not assert
global uniqueness over nonlocal functionals, additional fields, EFT operators,
quotient-topology data, initial covariance choices, or alternative source lifts.
Likewise, the canonical VDE equations are not promoted to a valuation-derived
Boltzmann theory.  A source-level implementation package is now supplied for
the nested canonical lift using a pinned non-flat CLASS backend and the harmonic
convention stated above.  The package has not produced validated spectra in the
present manuscript environment, so implementation availability is kept distinct
from executed numerical evidence.

The remaining open chain is therefore explicit:
\begin{equation}
\boxed{
\begin{gathered}
\Efour^{\rm val}\not\Rightarrow S_{\rm src}
\not\Rightarrow \text{amplitude selector},\\
\text{validated compact transfer output}
\not\Rightarrow \text{likelihood evidence}.
\end{gathered}}
\end{equation}
This is a scope map, not a ranking of scientific value.

\section{Conclusion}
\label{sec:conclusion}

This paper proves a finite homogeneous classification theorem, not a topology-only derivation of the observed dark sector.  Starting from the work map
\(\Efour(A)\mapsto(\epsD,\pD)\), homogeneous reconstruction is exact but non-predictive.  Predictivity begins only after imposing the Euclidean convex valuation hypothesis on the auxiliary round \(\Bfour_A\) filling.  Under that hypothesis the residual density belongs to
\begin{equation}
\mathcal V_{\rm val}=\operatorname{span}\{A,A^0,A^{-1},A^{-2},A^{-3}\},
\qquad
P_{\rm val}(\Dil)\epsD=0,
\end{equation}
with
\begin{equation}
P_{\rm val}(z)=(z-1)z(z+1)(z+2)(z+3),
\qquad
\mathfrak A_{\Dil}\simeq\mathbb R[z]/(P_{\rm val})\simeq\mathbb R^5.
\end{equation}
The corresponding primitive projectors, minimal annihilators, spectral lattice,
Hankel/Prony tests, and branch fractions give exact finite-dimensional null
tests for reconstructed homogeneous residual histories.

For positive branch energies the dilation spectrum carries a sharp convex
structure:
\begin{equation}
\Dil^2\ln(\epsD/\varepsilon_*)=\operatorname{Var}_\pi(s)\ge0,
\qquad
\Dil w_D=-\frac13\operatorname{Var}_\pi(s)\le0.
\end{equation}
The associated Popoviciu and Bhatia--Davis bounds, the endpoint saturation
criterion, and the local CPL tangent sign \(w_a\ge0\) are consequences of the
finite positive spectral measure.  These are falsification handles; they are not observational detections.

The filling-relative local curvature-measure boundary filter removes the bulk
\(s=1\) branch.  In \(\mathcal H^2(a)\), the quotient by the background-degenerate
\(a^0\), \(a^{-2}\), and \(a^{-3}\) directions leaves the nonstandard
\(a^{-1}\) branch.  The physical \(a^{-2}\) stress must nevertheless be separated
from geometric curvature by transverse-distance information:
\begin{equation}
\Omega_{2,{\rm val}}=\mathcal C_{-2}[\mathcal H^2]-\Omega_K^{\rm geom},
\qquad
\Omega_K^{\rm geom}=\frac{[\mathcal H d']^2-1}{d^2}.
\end{equation}
Thus the general boundary-profile extension is \((\Omega_X,\Omega_{2,{\rm val}})\), while the one-parameter canonical VDE model is the nested subclass \(\Omega_{2,{\rm val}}=0\).

The source layer is deliberately separate.  Positive negative-exponent branches
admit independent canonical-scalar existence realizations, with fixed
\(c_{s,{\rm rest}}^2=1\), \(\pi_{\mu\nu}=0\), reconstructed trajectory, and mass-flow invariant.  The \(s=-3\) dust branch requires a conserved current to reproduce CDM perturbations; topology supplies at most a global label matched by \(\mathcal E_{\rm rest}\overset{\rm match}{=}\mu\QU\).  The ADM \(R^{(3)}\) and \(\sqrt{R^{(3)}}\) constructions are response prototypes and diagnostics, not ordinary conserved fluids and not Hadwiger valuation functionals.  The fixed-shape Bianchi IX section supplies only a convention-dependent first mode-coupling diagnostic; nonlinear Mixmaster dynamics remain outside scope.

The amplitude problem is closed only negatively for dimensional quantities: scale covariance forbids pure topology or dimensionless data from determining nonzero dimensional coefficients such as \(\sigma\), \(\mu\), \(H_0\), unnormalized branch densities, mass scales, or cutoffs.  Dimensionless ratios such as \(\Omega_X\) and \(\Omega_{2,{\rm val}}\) are not selected by this no-go theorem alone; they require a declared normalization, critical density, amplitude ledger, matching condition, or observational input.  For the nested canonical VDE source class, a pinned non-flat CLASS transfer interface is supplied as supplementary material, including closed-mode bookkeeping, parameter fixtures, source/output checks, and validation tests.  No executed transfer spectra are reported from the present environment, and the general two-parameter boundary model, ADM/Bianchi-IX numerical extensions, full compact covariance validation, blind tests, and common likelihood comparison remain unsupplied numerical tasks.

Within the declared homogeneous invariant problem, these structures form a unique invariant normal form: the spectrum, commutant, primitive projectors, support/annihilator lattice, positive branch coordinates, boundary-local face, and background quotient are forced.  This conditional uniqueness does not extend to enlarged source, EFT, nonlocal, transfer, or likelihood categories.

The scope of the result is therefore
\begin{equation}
\boxed{
\begin{array}{c}
\text{finite homogeneous invariant classification}\\
\text{with independent source-existence examples}\\
\text{and sharp null tests.}
\end{array}}
\end{equation}
It is not a completed cosmological model, not a topology-only derivation of dark energy, and not an observational preference claim.  In particular,
\begin{equation}
\begin{array}{c}
\text{reconstruction}\neq\text{source}\neq\text{amplitude}\\
\neq\text{transfer}\neq\text{likelihood}.
\end{array}
\end{equation}

\section*{Data and Code Availability}
This theoretical work generated no new observational data sets, and no numerical output is used as evidence for the analytic theorems.  The accompanying analytic package includes the LaTeX source and compiled PDF together with build metadata, SHA-256 checksums, and machine-readable claim/scope ledgers.  A separate supplementary numerical package is also supplied for the nested canonical VDE lift.  It contains a version-pinned interface to unmodified non-flat CLASS, parameter fixtures, an \(\Sthree\) mode mapper, transfer/output runners, a curved-\(\Lambda\)CDM control, analytic/unit tests, an integration smoke test, and a validation plan.  The analytic/unit tests pass in the present environment; the \texttt{classy}-dependent integration test was not executed here because the dependency is unavailable.  No executed Boltzmann spectra, likelihood chain, observational data product, ADM/Bianchi-IX transfer implementation, or numerical model-preference claim is supplied.


\appendix
\section{Bianchi IX Volume--Shape Robustness and Anisotropic Mode Coupling}
\label{sec:bianchi-ix}

The preceding valuation and VDE theorems are round and homogeneous.  A minimal
anisotropic stress test is therefore not a claim of full Mixmaster dynamics, but
a volume--shape robustness question: does the dilation grading survive when the
round representative is replaced by a homogeneous squashed \(\Sthree\)?  We answer
only this fixed-shape question and extract the first mode-coupling matrix needed
for a later compact-SVT implementation.

Let \(\sigma^i\) be left-invariant one-forms on \(SU(2)\simeq\Sthree\), normalized so that the isotropic metric has scalar harmonics satisfying
\begin{equation}
-\nabla^2_{\Sthree_A}Y_{n\lambda}=\frac{n(n+2)}{A^2}Y_{n\lambda}.
\end{equation}
Throughout this appendix \(L_i\) denote self-adjoint left-action angular-momentum
operators on scalar harmonics, not raw real vector-field generators.  With real
left-invariant vector fields \(X_i\), the convention is \(L_i=-iX_i\); all
matrix elements below are interpreted in this \(L_i\) normalization.
A diagonal Bianchi IX representative may be written as
\begin{equation}
\dd s_3^2=A^2\sum_{i=1}^3 e^{2\beta_i}(\sigma^i)^2,
\qquad
\sum_{i=1}^3\beta_i=0,
\label{eq:BIX-metric}
\end{equation}
where \(A\) is the volume scale and \(\beta_i\) are dimensionless shape variables.
The dilation operator is the volume-scale derivative at fixed shape,
\begin{equation}
\Dil=A\frac{\partial}{\partial A}\Big|_{\beta}.
\end{equation}

\begin{theorem}[Fixed-shape robustness of the dilation grading]
\label{thm:BIX-dilation-grading}
Let \(L_i=-iX_i\) be the self-adjoint generators associated with the
left-invariant vector fields \(X_i\), with the normalization
\begin{equation}
\sum_{i=1}^3L_i^2Y_{n\lambda}=n(n+2)Y_{n\lambda}
\end{equation}
in the isotropic round limit.  The scalar Bianchi IX Laplacian is then defined by
\begin{equation}
\Delta_{\rm IX}
=-\frac1{A^2}\sum_{i=1}^3 e^{-2\beta_i}L_i^2,
\label{eq:BIX-laplacian}
\end{equation}
so that \(-\Delta_{\rm IX}\) is nonnegative and reduces exactly to the convention
\(-\nabla^2_{\Sthree_A}Y_{n\lambda}=n(n+2)A^{-2}Y_{n\lambda}\) when
\(\beta_i=0\).  At fixed \(\beta_i\),
\begin{equation}
[\Dil,L_i]=0,
\qquad
[\Dil,A^2\Delta_{\rm IX}]=0.
\end{equation}
Moreover, any branch of the form
\begin{equation}
\varepsilon_s(A,\beta)=C_s(\beta)A^s
\end{equation}
remains a dilation eigenbranch:
\begin{equation}
\Dil\varepsilon_s=s\varepsilon_s.
\end{equation}
Thus fixed-shape Bianchi IX deformation preserves the dilation eigenvalue but
can promote the branch amplitude to a shape-dependent coefficient.
\end{theorem}

\begin{proof}
The generators \(L_i\) act on the group coordinate, while \(\Dil\) acts only on
the scale \(A\).  Hence \([\Dil,L_i]=0\).  In \(A^2\Delta_{\rm IX}\) the explicit
factor \(A^{-2}\) has been removed, and the remaining operator depends only on
the fixed shape parameters \(\beta_i\) and the generators \(L_i\).  Therefore it
commutes with \(\Dil\).  Finally,
\(\Dil[C_s(\beta)A^s]=sC_s(\beta)A^s\) because the shape variables are held fixed.
\end{proof}

\begin{remark}[Scope of the Bianchi IX test]
\Cref{thm:BIX-dilation-grading} is not a theorem about nonlinear Mixmaster
evolution, anisotropic backreaction, or stability of the physical universe near
the singularity.  It is the precise fixed-shape statement needed to test whether
the finite dilation grading is destroyed by homogeneous squashing of the compact
spatial representative.  The answer is no; the grading survives, while shape
perturbations can still mix compact harmonic modes within the transfer system.
\end{remark}

Choose an orthonormal scalar harmonic basis \(\{|n\lambda\rangle\}\) for the round
\(\Sthree\), where \(n\) labels the round eigenvalue \(n(n+2)\) and \(\lambda\)
collects the degeneracy labels.  Write a composite index \(I=(n,\lambda)\).  The
Bianchi IX Laplacian has matrix elements
\begin{equation}
(\Delta_{\rm IX})_{IJ}
=-\frac1{A^2}\sum_{i=1}^3 e^{-2\beta_i}\langle I|L_i^2|J\rangle .
\end{equation}
For small traceless anisotropy, \(e^{-2\beta_i}=1-2\beta_i+O(\beta^2)\), hence
\begin{equation}
(\Delta_{\rm IX})_{IJ}
=-\frac1{A^2}
\left[
 n(n+2)\delta_{IJ}
 -2\sum_{i=1}^3\beta_i\langle I|L_i^2|J\rangle
 +O(\beta^2)
\right],
\end{equation}
where the first term is diagonal in the round basis.  The first anisotropic
mode-coupling matrix is therefore
\begin{equation}
\boxed{
\mathcal M^{\rm IX}_{IJ}
:=-2\sum_{i=1}^3\beta_i\langle I|L_i^2|J\rangle .}
\label{eq:BIX-mode-coupling}
\end{equation}

\begin{remark}[Basis dependence of \(L_i^2\) matrix elements]
The displayed expression is basis invariant, but the numerical entries of \(\langle I|L_i^2|J\rangle\) are not universal arrays independent of representation.  They must be evaluated in the harmonic basis actually used by the transfer code.  In a Wigner-\(D\) basis, in a coupled \(SO(4)\) basis, and in a hyperspherical \((n,\ell,m)\) basis the sparsity pattern and Clebsch--Gordan bookkeeping may look different, although the resulting operator is the same.  The only invariant requirements used here are block structure within the round Laplacian eigenspaces, symmetry with respect to the chosen orthonormal inner product, and the isotropic reduction as \(\beta_i\to0\). With any different generator normalization or metric convention, the sign and numerical factor in \(\mathcal M^{\rm IX}_{IJ}\) must be recalibrated before implementation.
\end{remark}

\begin{corollary}[Bianchi IX deformation of the canonical VDE kernel]
\label{cor:BIX-VDE-kernel}
To first order in the homogeneous shape variables, the round canonical VDE
kernel \(\mathscr W_n\) is replaced by the matrix kernel
\begin{equation}
\boxed{
\mathscr W^{\rm IX}_{IJ}(a,\beta)=
(-\Omega_K^{\rm geom})a^{-2}
\left[n(n+2)\delta_{IJ}+\mathcal M^{\rm IX}_{IJ}\right]
+\frac54\left[F_{\rm VDE}-\Dil F_{\rm VDE}\right]\delta_{IJ}
+O(\beta^2).}
\label{eq:BIX-VDE-kernel}
\end{equation}
The mass-flow term remains diagonal because it is a homogeneous background
quantity.  The shape variables enter the scalar transfer problem through the
compact Laplacian matrix, not through a new sound speed, anisotropic-stress
parameter, or mass function for the canonical VDE source.
\end{corollary}

\begin{proof}
Replace the round Laplacian contribution \(n(n+2)/A^2\) in
\cref{eq:VDE-transfer-kernel} by the Bianchi IX Laplacian matrix expansion above.
The scalar mass-flow term is already fixed by the background potential and is
therefore proportional to \(\delta_{IJ}\) at this order.  No additional source
parameter is introduced.
\end{proof}

\begin{criterion}[Anisotropic robustness null]
A transfer-grade extension of canonical VDE to weak homogeneous anisotropy must
reduce, in the limit \(\beta_i\to0\), to the round closed-\(\Sthree\) kernel
\cref{eq:VDE-transfer-kernel}, and its first fixed-shape anisotropic correction
must agree with \cref{eq:BIX-VDE-kernel}.  Failure of this reduction would signal
an implementation error or an additional physical coupling beyond the minimal
canonical VDE completion.
\end{criterion}

\begin{remark}[Hopf labels versus dust dynamics]
The Hopf fibration may supply useful global labels for internal sectors, but it
does not replace the metric-dependent pressureless-current action.  A Hopf or
linking invariant can label a superselection sector; it does not by itself imply
\(T_{\mu\nu}=\rho u_\mu u_\nu\), \(c_s^2=0\), or
\(\pi_{\mu\nu}=0\).  Those properties remain supplied by the conserved-current
dust realization.  This preserves the separation
\(\text{topological label}\neq\text{local source}\neq\text{amplitude}\).
\end{remark}

\section{\texorpdfstring{$\Sfour$}{S4} Compensation by Gluing}
Consider the standard double
\begin{equation}
\Sfour=\Bfour_+\cup_{\mathrm{id}}\Bfour_-,
\end{equation}
where the common \(\Sthree\) boundary carries the opposite induced orientations. The following statement is only a boundary-orientation lemma; it does not prove cancellation of energies, actions, stresses, amplitudes, or equations of motion.

\begin{proposition}[Boundary-form cancellation, not dynamical compensation]
Let \(\omega_\pm\) be boundary forms satisfying
\begin{equation}
\iota_+^*\omega_+=\iota_-^*\omega_-.
\end{equation}
Then, with opposite induced orientations,
\begin{equation}
\int_{\partial\Bfour_+}\omega_+ + \int_{\partial\Bfour_-}\omega_-=0.
\end{equation}
\end{proposition}
\begin{proof}
The identity gluing identifies the values of the forms on the common boundary, while the induced orientations supply opposite signs. The two contributions are therefore equal in magnitude and opposite in sign.
\end{proof}

\begin{remark}[Chern--Simons boundary actions and amplitude limitation]
For a boundary connection \(\mathcal{A}\), one may form the Chern--Simons three-form \cite{ChernSimons1974}
\begin{equation}
\operatorname{CS}_3(\mathcal{A})=
\operatorname{tr}\left(\mathcal{A}\wedge\dd\mathcal{A}
+\frac{2}{3}\mathcal{A}\wedge\mathcal{A}\wedge\mathcal{A}\right),
\end{equation}
with
\begin{equation}
\dd\operatorname{CS}_3(\mathcal{A})=\operatorname{tr}(F_{\mathcal{A}}\wedge F_{\mathcal{A}}).
\end{equation}
For the standard double \(\Sfour=\Bfour_+\cup_{\mathrm{id}}\Bfour_-\), compatible boundary data and opposite induced orientations imply cancellation of the oriented boundary contributions up to the usual quantized large-gauge shift fixed by the normalization of the Chern--Simons coupling. This is a statement about oriented topological consistency of boundary data. It does not determine the dimensional amplitudes \(\sigma\) or \(\mu\), which require independent couplings or physical scales.
\end{remark}


\begin{nogotheorem}[No vacuum amplitude from Chern--Simons gluing alone]
A Chern--Simons boundary functional can supply an oriented topological phase or a quantized large-gauge shift, but it does not by itself determine the dimensionful vacuum energy density \(\sigma\).
\end{nogotheorem}
\begin{proof}
The Chern--Simons functional changes under a large gauge transformation by a quantized amount fixed by the normalization of the Chern--Simons coupling. This quantized shift is a statement about the boundary action or phase. The endpoint vacuum contribution is instead
\begin{equation}
E_q=\sigma V_3=2\pi^2\sigma A^3,
\end{equation}
where
\begin{equation}
[\sigma]=EL^{-3}.
\end{equation}
Equating \(E_q\) to a Chern--Simons shift requires an additional dimensionful coupling, time scale, cutoff, or Hamiltonian normalization. Therefore Chern--Simons gluing may constrain a topological sector or a phase, but it cannot determine \(\sigma\) without an independent dimensional input.
\end{proof}

Orientation supplies the sign but not equality of the boundary data. Functionals containing normal derivatives, extrinsic curvature, independent fields, or corner terms require additional matching conditions. The construction of a concrete \(\Sfour\) functional that selects the amplitudes \(\varepsilon_q\) and \(\varepsilon_w\) is therefore a separate completion condition.


\section{Suspension, \texorpdfstring{$SU(2)$}{SU(2)}, and Hopf Structure}
The exact suspension relations are
\begin{equation}
S^1\simeq\Sigma S^0,
\qquad
S^2\simeq\Sigma S^1,
\qquad
S^3\simeq\Sigma S^2,
\qquad
S^4\simeq\Sigma S^3.
\end{equation}
Separately, the Hopf fibration and the standard structural identifications are
\begin{equation}
S^1\hookrightarrow S^3\to S^2,
\qquad
\Sthree\simeq SU(2),
\qquad
T\Sthree\simeq \Sthree\times\mathbb{R}^3.
\end{equation}
For the chosen normalization of the left-invariant Maurer-Cartan forms,
\begin{equation}
\dd\sigma^i=-\frac{1}{2}\epsilon^i{}_{jk}\sigma^j\wedge\sigma^k.
\end{equation}
Parallelizability and the Hopf fibration provide global geometric structures capable of carrying directions, phases, and flows \cite{Hopf1931}, but they do not by themselves determine an energy or the amplitude of the residual sector.

\begin{hypothesis}
Global Hopf modes may enter the internal variables \(\xi^I\) of the filling energy. Their physical imprint would then have to appear through \(\QD\), anisotropic stress, a surface charge, or the spectrum of global modes, rather than through a topological identification alone.
\end{hypothesis}


\section{Light and a Causal Measure}
The radial null condition for the metric gives
\begin{equation}
\dd\chi=\frac{cN(t)}{A(t)}\dd t.
\end{equation}
Define the causal angular radius
\begin{equation}
\chi_c(t_i,t):=\min\left\{\pi,\int_{t_i}^t\frac{cN(t')}{A(t')}\,\dd t'\right\}.
\end{equation}
The causally accessible fraction of the three-volume is
\begin{equation}
I_c(t_i,t):=\frac{V_3(\chi_c,A)}{V_3(\pi,A)}
=\frac{\chi_c-\sin\chi_c\cos\chi_c}{\pi},
\qquad
0\le I_c\le1.
\end{equation}
The antipode is reached at accumulated conformal length \(\pi\), while a complete geodesic circuit requires \(2\pi\). Causal accessibility is not identified with entropic or quantum coherence.

\begin{hypothesis}
If \(\Efour=\Efour(A,I_c,\xi)\), then the causal variable contributes a definite exchange term,
\begin{equation}
Q_D^{(I_c)}=\frac{1}{V_3}\frac{\partial\Efour}{\partial I_c}\frac{\dd I_c}{\dd\tau}.
\end{equation}
\end{hypothesis}



\begin{thebibliography}{99}

\bibitem{Noether1918}
E. Noether, ``Invariante Variationsprobleme,'' \emph{Nachrichten von der Gesellschaft der Wissenschaften zu G\"ottingen, Mathematisch-Physikalische Klasse} (1918), 235--257.

\bibitem{LeeWald1990}
J. Lee and R. M. Wald, ``Local symmetries and constraints,'' \emph{Journal of Mathematical Physics} \textbf{31} (1990), 725--743, \href{https://doi.org/10.1063/1.528801}{doi:10.1063/1.528801}.

\bibitem{BrownYork1993}
J. D. Brown and J. W. York, ``Quasilocal energy and conserved charges derived from the gravitational action,'' \emph{Physical Review D} \textbf{47} (1993), 1407--1419, \href{https://doi.org/10.1103/PhysRevD.47.1407}{doi:10.1103/PhysRevD.47.1407}.

\bibitem{IyerWald1994}
V. Iyer and R. M. Wald, ``Some properties of Noether charge and a proposal for dynamical black hole entropy,'' \emph{Physical Review D} \textbf{50} (1994), 846--864, \href{https://doi.org/10.1103/PhysRevD.50.846}{doi:10.1103/PhysRevD.50.846}.

\bibitem{BarnichBrandt2002}
G. Barnich and F. Brandt, ``Covariant theory of asymptotic symmetries, conservation laws and central charges,'' \emph{Nuclear Physics B} \textbf{633} (2002), 3--82, \href{https://doi.org/10.1016/S0550-3213(02)00251-1}{doi:10.1016/S0550-3213(02)00251-1}.

\bibitem{GibbonsHawking1977}
G. W. Gibbons and S. W. Hawking, ``Action integrals and partition functions in quantum gravity,'' \emph{Physical Review D} \textbf{15} (1977), 2752--2756, \href{https://doi.org/10.1103/PhysRevD.15.2752}{doi:10.1103/PhysRevD.15.2752}.

\bibitem{DeWitt1967}
B. S. DeWitt, ``Quantum theory of gravity. I. The canonical theory,'' \emph{Physical Review} \textbf{160} (1967), 1113--1148, \href{https://doi.org/10.1103/PhysRev.160.1113}{doi:10.1103/PhysRev.160.1113}.

\bibitem{HawkingEllis1973}
S. W. Hawking and G. F. R. Ellis, \emph{The Large Scale Structure of Space-Time}, Cambridge University Press, Cambridge, 1973.

\bibitem{MTW1973}
C. W. Misner, K. S. Thorne, and J. A. Wheeler, \emph{Gravitation}, W. H. Freeman, San Francisco, 1973.

\bibitem{Weinberg2008}
S. Weinberg, \emph{Cosmology}, Oxford University Press, Oxford, 2008.

\bibitem{KodamaSasaki1984}
H. Kodama and M. Sasaki, ``Cosmological perturbation theory,'' \emph{Progress of Theoretical Physics Supplement} \textbf{78} (1984), 1--166, \href{https://doi.org/10.1143/PTPS.78.1}{doi:10.1143/PTPS.78.1}.

\bibitem{HalliwellHawking1985}
J. J. Halliwell and S. W. Hawking, ``Origin of structure in the Universe,'' \emph{Physical Review D} \textbf{31} (1985), 1777--1791, \href{https://doi.org/10.1103/PhysRevD.31.1777}{doi:10.1103/PhysRevD.31.1777}.

\bibitem{MukhanovFeldmanBrandenberger1992}
V. F. Mukhanov, H. A. Feldman, and R. H. Brandenberger, ``Theory of cosmological perturbations,'' \emph{Physics Reports} \textbf{215} (1992), 203--333, \href{https://doi.org/10.1016/0370-1573(92)90044-Z}{doi:10.1016/0370-1573(92)90044-Z}.

\bibitem{MalikWands2009}
K. A. Malik and D. Wands, ``Cosmological perturbations,'' \emph{Physics Reports} \textbf{475} (2009), 1--51, \href{https://doi.org/10.1016/j.physrep.2009.03.001}{doi:10.1016/j.physrep.2009.03.001}.

\bibitem{Hopf1931}
H. Hopf, ``\"Uber die Abbildungen der dreidimensionalen Sph\"are auf die Kugelfl\"ache,'' \emph{Mathematische Annalen} \textbf{104} (1931), 637--665, \href{https://doi.org/10.1007/BF01457962}{doi:10.1007/BF01457962}.

\bibitem{CornishSpergelStarkman1998}
N. J. Cornish, D. N. Spergel, and G. D. Starkman, ``Circles in the sky: finding topology with the microwave background radiation,'' \emph{Classical and Quantum Gravity} \textbf{15} (1998), 2657--2670, \href{https://doi.org/10.1088/0264-9381/15/9/005}{doi:10.1088/0264-9381/15/9/005}.

\bibitem{LuminetEtAl2003}
J.-P. Luminet, J. R. Weeks, A. Riazuelo, R. Lehoucq, and J.-P. Uzan, ``Dodecahedral space topology as an explanation for weak wide-angle temperature correlations in the cosmic microwave background,'' \emph{Nature} \textbf{425} (2003), 593--595, \href{https://doi.org/10.1038/nature01944}{doi:10.1038/nature01944}.


\bibitem{KonigRothMecke2004}
P.-M. K\"onig, R. Roth, and K. R. Mecke, ``Morphological thermodynamics of fluids: Shape dependence of free energies,'' \emph{Physical Review Letters} \textbf{93} (2004), 160601, \href{https://doi.org/10.1103/PhysRevLett.93.160601}{doi:10.1103/PhysRevLett.93.160601}.

\bibitem{Hadwiger1957}
H. Hadwiger, \emph{Vorlesungen \"uber Inhalt, Oberfl\"ache und Isoperimetrie}, Springer, Berlin, 1957.

\bibitem{Schneider2014}
R. Schneider, \emph{Convex Bodies: The Brunn-Minkowski Theory}, 2nd expanded ed., Cambridge University Press, Cambridge, 2014.

\bibitem{BhatiaDavis2000}
R. Bhatia and C. Davis, ``A better bound on the variance,''
\emph{American Mathematical Monthly} \textbf{107} (2000), 353--357,
\href{https://doi.org/10.1080/00029890.2000.12005203}{doi:10.1080/00029890.2000.12005203}.

\bibitem{Brown1993Fluid}
J. D. Brown, ``Action functionals for relativistic perfect fluids,'' \emph{Classical and Quantum Gravity} \textbf{10} (1993), 1579--1606, \href{https://doi.org/10.1088/0264-9381/10/8/017}{doi:10.1088/0264-9381/10/8/017}.

\bibitem{ChernSimons1974}
S.-S. Chern and J. Simons, ``Characteristic forms and geometric invariants,'' \emph{Annals of Mathematics} \textbf{99} (1974), 48--69, \href{https://doi.org/10.2307/1971013}{doi:10.2307/1971013}.

\bibitem{ClarksonBassettLu2008}
C. Clarkson, B. Bassett, and T. H.-C. Lu, ``A general test of the Copernican Principle,'' \emph{Physical Review Letters} \textbf{101} (2008), 011301, \href{https://doi.org/10.1103/PhysRevLett.101.011301}{doi:10.1103/PhysRevLett.101.011301}.

\bibitem{BrownKuchar1995}
J. D. Brown and K. V. Kucha\v{r}, ``Dust as a standard of space and time in canonical quantum gravity,'' \emph{Physical Review D} \textbf{51} (1995), 5600--5629, \href{https://doi.org/10.1103/PhysRevD.51.5600}{doi:10.1103/PhysRevD.51.5600}.



\bibitem{SahlenLiddleParkinson2005}
M. Sahl\'en, A. R. Liddle, and D. Parkinson, ``Direct reconstruction of the quintessence potential,'' \emph{Physical Review D} \textbf{72} (2005), 083511, \href{https://doi.org/10.1103/PhysRevD.72.083511}{doi:10.1103/PhysRevD.72.083511}, arXiv:astro-ph/0506696.

\bibitem{BlasLesgourguesTram2011}
D. Blas, J. Lesgourgues, and T. Tram, ``The Cosmic Linear Anisotropy Solving System (CLASS). Part II: Approximation schemes,'' \emph{Journal of Cosmology and Astroparticle Physics} \textbf{07} (2011), 034, \href{https://doi.org/10.1088/1475-7516/2011/07/034}{doi:10.1088/1475-7516/2011/07/034}, arXiv:1104.2933 [astro-ph.CO].

\bibitem{LesgourguesTram2014}
J. Lesgourgues and T. Tram, ``Fast and accurate CMB computations in non-flat FLRW universes,'' \emph{Journal of Cosmology and Astroparticle Physics} \textbf{09} (2014), 032, \href{https://doi.org/10.1088/1475-7516/2014/09/032}{doi:10.1088/1475-7516/2014/09/032}, arXiv:1312.2697 [astro-ph.CO].

\bibitem{AnHanZhang2025}
H. An, C. Han, and B. Zhang, ``Topological defects as effective dynamical dark energy,'' \emph{Physical Review D} \textbf{113} (2026), 043543, \href{https://doi.org/10.1103/2m4c-1zl5}{doi:10.1103/2m4c-1zl5}, arXiv:2506.10075 [hep-ph].

\bibitem{DESY6ThreeByTwo2026}
Dark Energy Survey Collaboration, T. M. C. Abbott et al., ``Dark Energy Survey Year 6 results: Cosmological constraints from galaxy clustering and weak lensing,'' arXiv:2601.14559 [astro-ph.CO] (2026).

\bibitem{DESIExtendedDE2025}
DESI Collaboration, K. Lodha et al., ``Extended Dark Energy analysis using DESI DR2 BAO measurements,'' arXiv:2503.14743 [astro-ph.CO] (2025).

\bibitem{CMBComp2026}
A. Srivastav, P. Bansal, and D. Huterer, ``CMBComp: A simple and accurate compressed CMB likelihood for dark energy, curvature, and massive neutrinos,'' arXiv:2606.18455 [astro-ph.CO] (2026).

\bibitem{PulidoCervantes2026}
H. Pulido-Hernandez and J. L. Cervantes-Cota, ``Negative masses and spatial curvature: Alleviating neutrino mass tensions in \(\Lambda\)CDM and extended cosmologies,'' arXiv:2603.13208 [astro-ph.CO] (2026).

\bibitem{CosmicTopologyReview2026}
C. J. Copi et al., ``The topology of the Universe,'' arXiv:2606.24886 [astro-ph.CO] (2026).

\bibitem{COMPACTLensSpaces2025}
S. Saha et al. (COMPACT Collaboration), ``Cosmic topology. Part Ic. Limits on lens spaces from circle searches,'' \emph{Journal of Cosmology and Astroparticle Physics} \textbf{01} (2025), 004, \href{https://doi.org/10.1088/1475-7516/2025/01/004}{doi:10.1088/1475-7516/2025/01/004}, arXiv:2409.02226 [astro-ph.CO].

\bibitem{DESY6DynamicalDE2026}
Dark Energy Survey Collaboration, T. M. C. Abbott et al., ``Constraints on Dynamical Dark Energy from Multiple Probes in the Full Dark Energy Survey,'' arXiv:2605.27221 [astro-ph.CO] (2026).

\bibitem{ACTDR6Extended2025}
E. Calabrese et al. (Atacama Cosmology Telescope Collaboration), ``The Atacama Cosmology Telescope: DR6 constraints on extended cosmological models,'' \emph{Journal of Cosmology and Astroparticle Physics} \textbf{11} (2025), 063, \href{https://doi.org/10.1088/1475-7516/2025/11/063}{doi:10.1088/1475-7516/2025/11/063}, arXiv:2503.14454 [astro-ph.CO].

\bibitem{EuclidTimeline2026}
Euclid Consortium, ``Timeline of the Euclid Survey,'' updated June 19, 2026, \href{https://www.euclid-ec.org/public/timeline/}{Euclid Consortium public timeline}.

\bibitem{EuclidCosmology2026}
Euclid Consortium, ``Euclid starts seeing darkness,'' June 5, 2026, \href{https://www.euclid-ec.org/euclid-starts-seeing-darkness/}{Euclid Consortium science communication}, noting that the first main cosmology results are expected in 2027.

\bibitem{Planck2018Parameters}
Planck Collaboration, ``Planck 2018 results. VI. Cosmological parameters,'' \emph{Astronomy \& Astrophysics} \textbf{641} (2020), A6, \href{https://doi.org/10.1051/0004-6361/201833910}{doi:10.1051/0004-6361/201833910}.

\bibitem{DESI2025DR2BAO}
DESI Collaboration, M. Abdul Karim et al., ``DESI DR2 results. II. Measurements of baryon acoustic oscillations and cosmological constraints,'' \emph{Physical Review D} \textbf{112} (2025), 083515, \href{https://doi.org/10.1103/tr6y-kpc6}{doi:10.1103/tr6y-kpc6}, arXiv:2503.14738 [astro-ph.CO].

\bibitem{Chernoff1954}
H. Chernoff, ``On the distribution of the likelihood ratio,'' \emph{Annals of Mathematical Statistics} \textbf{25} (1954), 573--578, \href{https://doi.org/10.1214/aoms/1177728725}{doi:10.1214/aoms/1177728725}.

\bibitem{SelfLiang1987}
S. G. Self and K.-Y. Liang, ``Asymptotic properties of maximum likelihood estimators and likelihood ratio tests under nonstandard conditions,'' \emph{Journal of the American Statistical Association} \textbf{82} (1987), 605--610, \href{https://doi.org/10.1080/01621459.1987.10478472}{doi:10.1080/01621459.1987.10478472}.

\bibitem{Planck2016Topology}
Planck Collaboration, ``Planck 2015 results. XVIII. Background geometry and topology of the Universe,'' \emph{Astronomy \& Astrophysics} \textbf{594} (2016), A18, \href{https://doi.org/10.1051/0004-6361/201525829}{doi:10.1051/0004-6361/201525829}.

\end{thebibliography}

\end{document}
